\begin{enumerate}[label=\Alph*.]
\item Expecting global search to answer precise record lookups wastes cost and may lose detail.
\item Confusing open-world semantics with database constraints produces invalid assumptions.
\item Reporting only downstream answer scores can hide a brittle graph with unverifiable facts.
\item Using untyped generic edges turns the graph into an expensive document store.
\end{enumerate}
\noindent\textbf{Answer.} Expecting global search to answer precise record lookups wastes cost and may lose detail.
\par\textit{Global search summarizes graph communities and combines relevant community reports for corpus-level questions. Tradeoff: It addresses whole-corpus themes but has expensive indexing and lossy summaries. Watch for: Expecting global search to answer precise record lookups wastes cost and may lose detail.}
\paragraph{FC-0478 [medium]} Define GraphRAG global search in interview-ready language.
\noindent\textbf{Answer.} Global search summarizes graph communities and combines relevant community reports for corpus-level questions.
\par\textit{Global search summarizes graph communities and combines relevant community reports for corpus-level questions.}
\paragraph{FC-0479 [hard]} State the main tradeoff and production pitfall for GraphRAG global search.
\noindent\textbf{Answer.} Tradeoff: It addresses whole-corpus themes but has expensive indexing and lossy summaries. Pitfall: Expecting global search to answer precise record lookups wastes cost and may lose detail.
\par\textit{Tradeoff: It addresses whole-corpus themes but has expensive indexing and lossy summaries. Pitfall: Expecting global search to answer precise record lookups wastes cost and may lose detail.}
\paragraph{LA-0480 [hard]} Explain GraphRAG global search from first principles, then show how you would evaluate and operate it in production.
\noindent\textbf{Answer.} Start with the mechanism: Global search summarizes graph communities and combines relevant community reports for corpus-level questions. Make the tradeoff explicit: It addresses whole-corpus themes but has expensive indexing and lossy summaries. Define workload-specific offline metrics and a latency/cost budget, test important slices, instrument the critical path, and create a rollback criterion. The production review must explicitly guard against this failure: Expecting global search to answer precise record lookups wastes cost and may lose detail.
\par\textit{A strong answer connects mechanism, measurable tradeoffs, failure isolation, observability, and rollback rather than listing product features.}
\paragraph{MCQ-0481 [medium]} Which statement best describes KAG?
\begin{enumerate}[label=\Alph*.]
\item Knowledge Augmented Generation combines schema-aware graphs, chunk mutual indexing, alignment, and logical-form-guided reasoning.
\item Construction extracts entities and relations, resolves identities, applies schema, validates facts, and records provenance.
\item Nodes, typed relationships, and properties represent entities and their connections for traversal and pattern matching.
\item Local search expands from query-relevant entities through nearby relationships and associated text.
\end{enumerate}
\noindent\textbf{Answer.} Knowledge Augmented Generation combines schema-aware graphs, chunk mutual indexing, alignment, and logical-form-guided reasoning.
\par\textit{Knowledge Augmented Generation combines schema-aware graphs, chunk mutual indexing, alignment, and logical-form-guided reasoning. Tradeoff: It improves domain logic and multi-hop control at substantial modeling and maintenance cost. Watch for: Adopting KAG without stable domain schema simply relocates ambiguity.}
\paragraph{MCQ-0482 [hard]} What is the central engineering tradeoff of KAG?
\begin{enumerate}[label=\Alph*.]
\item They are intuitive for operational queries but semantics depend on disciplined labels and schema.
\item It supports entity-centric multi-hop questions but depends on extraction and linking quality.
\item Automation scales coverage but introduces extraction and entity-resolution errors.
\item It improves domain logic and multi-hop control at substantial modeling and maintenance cost.
\end{enumerate}
\noindent\textbf{Answer.} It improves domain logic and multi-hop control at substantial modeling and maintenance cost.
\par\textit{Knowledge Augmented Generation combines schema-aware graphs, chunk mutual indexing, alignment, and logical-form-guided reasoning. Tradeoff: It improves domain logic and multi-hop control at substantial modeling and maintenance cost. Watch for: Adopting KAG without stable domain schema simply relocates ambiguity.}
\paragraph{MCQ-0483 [hard]} Which failure mode should an interviewer expect you to identify for KAG?
\begin{enumerate}[label=\Alph*.]
\item Unbounded neighborhood expansion floods context with high-degree hubs.
\item Using untyped generic edges turns the graph into an expensive document store.
\item Adopting KAG without stable domain schema simply relocates ambiguity.
\item Allowing generated edges without source spans prevents later verification.
\end{enumerate}
\noindent\textbf{Answer.} Adopting KAG without stable domain schema simply relocates ambiguity.
\par\textit{Knowledge Augmented Generation combines schema-aware graphs, chunk mutual indexing, alignment, and logical-form-guided reasoning. Tradeoff: It improves domain logic and multi-hop control at substantial modeling and maintenance cost. Watch for: Adopting KAG without stable domain schema simply relocates ambiguity.}
\paragraph{MCQ-0484 [medium]} A design review is specifically concerned with this risk: Adopting KAG without stable domain schema simply relocates ambiguity. Which topic should be reviewed first?
\begin{enumerate}[label=\Alph*.]
\item Knowledge graph construction
\item Property graphs
\item GraphRAG local search
\item KAG
\end{enumerate}
\noindent\textbf{Answer.} KAG
\par\textit{Knowledge Augmented Generation combines schema-aware graphs, chunk mutual indexing, alignment, and logical-form-guided reasoning. Tradeoff: It improves domain logic and multi-hop control at substantial modeling and maintenance cost. Watch for: Adopting KAG without stable domain schema simply relocates ambiguity.}
\paragraph{MCQ-0485 [hard]} Which design note most accurately balances the benefits and costs of KAG?
\begin{enumerate}[label=\Alph*.]
\item It supports entity-centric multi-hop questions but depends on extraction and linking quality.
\item They are intuitive for operational queries but semantics depend on disciplined labels and schema.
\item Automation scales coverage but introduces extraction and entity-resolution errors.
\item It improves domain logic and multi-hop control at substantial modeling and maintenance cost.
\end{enumerate}
\noindent\textbf{Answer.} It improves domain logic and multi-hop control at substantial modeling and maintenance cost.
\par\textit{Knowledge Augmented Generation combines schema-aware graphs, chunk mutual indexing, alignment, and logical-form-guided reasoning. Tradeoff: It improves domain logic and multi-hop control at substantial modeling and maintenance cost. Watch for: Adopting KAG without stable domain schema simply relocates ambiguity.}
\paragraph{MCQ-0486 [medium]} Which explanation would be strongest in a system-design interview about KAG?
\begin{enumerate}[label=\Alph*.]
\item Construction extracts entities and relations, resolves identities, applies schema, validates facts, and records provenance.
\item Local search expands from query-relevant entities through nearby relationships and associated text.
\item Knowledge Augmented Generation combines schema-aware graphs, chunk mutual indexing, alignment, and logical-form-guided reasoning.
\item Nodes, typed relationships, and properties represent entities and their connections for traversal and pattern matching.
\end{enumerate}
\noindent\textbf{Answer.} Knowledge Augmented Generation combines schema-aware graphs, chunk mutual indexing, alignment, and logical-form-guided reasoning.
\par\textit{Knowledge Augmented Generation combines schema-aware graphs, chunk mutual indexing, alignment, and logical-form-guided reasoning. Tradeoff: It improves domain logic and multi-hop control at substantial modeling and maintenance cost. Watch for: Adopting KAG without stable domain schema simply relocates ambiguity.}
\paragraph{MCQ-0487 [hard]} Which production warning is most directly associated with KAG?
\begin{enumerate}[label=\Alph*.]
\item Using untyped generic edges turns the graph into an expensive document store.
\item Allowing generated edges without source spans prevents later verification.
\item Unbounded neighborhood expansion floods context with high-degree hubs.
\item Adopting KAG without stable domain schema simply relocates ambiguity.
\end{enumerate}
\noindent\textbf{Answer.} Adopting KAG without stable domain schema simply relocates ambiguity.
\par\textit{Knowledge Augmented Generation combines schema-aware graphs, chunk mutual indexing, alignment, and logical-form-guided reasoning. Tradeoff: It improves domain logic and multi-hop control at substantial modeling and maintenance cost. Watch for: Adopting KAG without stable domain schema simply relocates ambiguity.}
\paragraph{FC-0488 [medium]} Define KAG in interview-ready language.
\noindent\textbf{Answer.} Knowledge Augmented Generation combines schema-aware graphs, chunk mutual indexing, alignment, and logical-form-guided reasoning.
\par\textit{Knowledge Augmented Generation combines schema-aware graphs, chunk mutual indexing, alignment, and logical-form-guided reasoning.}
\paragraph{FC-0489 [hard]} State the main tradeoff and production pitfall for KAG.
\noindent\textbf{Answer.} Tradeoff: It improves domain logic and multi-hop control at substantial modeling and maintenance cost. Pitfall: Adopting KAG without stable domain schema simply relocates ambiguity.
\par\textit{Tradeoff: It improves domain logic and multi-hop control at substantial modeling and maintenance cost. Pitfall: Adopting KAG without stable domain schema simply relocates ambiguity.}
\paragraph{LA-0490 [hard]} Explain KAG from first principles, then show how you would evaluate and operate it in production.
\noindent\textbf{Answer.} Start with the mechanism: Knowledge Augmented Generation combines schema-aware graphs, chunk mutual indexing, alignment, and logical-form-guided reasoning. Make the tradeoff explicit: It improves domain logic and multi-hop control at substantial modeling and maintenance cost. Define workload-specific offline metrics and a latency/cost budget, test important slices, instrument the critical path, and create a rollback criterion. The production review must explicitly guard against this failure: Adopting KAG without stable domain schema simply relocates ambiguity.
\par\textit{A strong answer connects mechanism, measurable tradeoffs, failure isolation, observability, and rollback rather than listing product features.}
\paragraph{MCQ-0491 [medium]} Which statement best describes Graph quality evaluation?
\begin{enumerate}[label=\Alph*.]
\item Cypher expresses graph patterns, predicates, paths, updates, and aggregation for property graphs.
\item Evaluation checks entity precision, relation correctness, completeness, consistency, provenance, and downstream QA value.
\item Facts carry validity intervals, event time, and version history so queries can reason about change.
\item Construction extracts entities and relations, resolves identities, applies schema, validates facts, and records provenance.
\end{enumerate}
\noindent\textbf{Answer.} Evaluation checks entity precision, relation correctness, completeness, consistency, provenance, and downstream QA value.
\par\textit{Evaluation checks entity precision, relation correctness, completeness, consistency, provenance, and downstream QA value. Tradeoff: Intrinsic graph metrics catch construction defects but may not predict application utility. Watch for: Reporting only downstream answer scores can hide a brittle graph with unverifiable facts.}
\paragraph{MCQ-0492 [hard]} What is the central engineering tradeoff of Graph quality evaluation?
\begin{enumerate}[label=\Alph*.]
\item Automation scales coverage but introduces extraction and entity-resolution errors.
\item Declarative queries are readable, while unbounded path patterns can be expensive.
\item Intrinsic graph metrics catch construction defects but may not predict application utility.
\item Temporal accuracy improves, while updates, conflicts, and query logic become harder.
\end{enumerate}
\noindent\textbf{Answer.} Intrinsic graph metrics catch construction defects but may not predict application utility.
\par\textit{Evaluation checks entity precision, relation correctness, completeness, consistency, provenance, and downstream QA value. Tradeoff: Intrinsic graph metrics catch construction defects but may not predict application utility. Watch for: Reporting only downstream answer scores can hide a brittle graph with unverifiable facts.}
\paragraph{MCQ-0493 [hard]} Which failure mode should an interviewer expect you to identify for Graph quality evaluation?
\begin{enumerate}[label=\Alph*.]
\item Allowing generated edges without source spans prevents later verification.
\item Omitting labels or relationship types can trigger broad scans.
\item Reporting only downstream answer scores can hide a brittle graph with unverifiable facts.
\item Overwriting old facts destroys the ability to answer as-of questions.
\end{enumerate}
\noindent\textbf{Answer.} Reporting only downstream answer scores can hide a brittle graph with unverifiable facts.
\par\textit{Evaluation checks entity precision, relation correctness, completeness, consistency, provenance, and downstream QA value. Tradeoff: Intrinsic graph metrics catch construction defects but may not predict application utility. Watch for: Reporting only downstream answer scores can hide a brittle graph with unverifiable facts.}
\paragraph{MCQ-0494 [medium]} A design review is specifically concerned with this risk: Reporting only downstream answer scores can hide a brittle graph with unverifiable facts. Which topic should be reviewed first?
\begin{enumerate}[label=\Alph*.]
\item Knowledge graph construction
\item Temporal knowledge graphs
\item Cypher
\item Graph quality evaluation
\end{enumerate}
\noindent\textbf{Answer.} Graph quality evaluation
\par\textit{Evaluation checks entity precision, relation correctness, completeness, consistency, provenance, and downstream QA value. Tradeoff: Intrinsic graph metrics catch construction defects but may not predict application utility. Watch for: Reporting only downstream answer scores can hide a brittle graph with unverifiable facts.}
\paragraph{MCQ-0495 [hard]} Which design note most accurately balances the benefits and costs of Graph quality evaluation?
\begin{enumerate}[label=\Alph*.]
\item Automation scales coverage but introduces extraction and entity-resolution errors.
\item Declarative queries are readable, while unbounded path patterns can be expensive.
\item Intrinsic graph metrics catch construction defects but may not predict application utility.
\item Temporal accuracy improves, while updates, conflicts, and query logic become harder.
\end{enumerate}
\noindent\textbf{Answer.} Intrinsic graph metrics catch construction defects but may not predict application utility.
\par\textit{Evaluation checks entity precision, relation correctness, completeness, consistency, provenance, and downstream QA value. Tradeoff: Intrinsic graph metrics catch construction defects but may not predict application utility. Watch for: Reporting only downstream answer scores can hide a brittle graph with unverifiable facts.}
\paragraph{MCQ-0496 [medium]} Which explanation would be strongest in a system-design interview about Graph quality evaluation?
\begin{enumerate}[label=\Alph*.]
\item Facts carry validity intervals, event time, and version history so queries can reason about change.
\item Cypher expresses graph patterns, predicates, paths, updates, and aggregation for property graphs.
\item Construction extracts entities and relations, resolves identities, applies schema, validates facts, and records provenance.
\item Evaluation checks entity precision, relation correctness, completeness, consistency, provenance, and downstream QA value.
\end{enumerate}
\noindent\textbf{Answer.} Evaluation checks entity precision, relation correctness, completeness, consistency, provenance, and downstream QA value.
\par\textit{Evaluation checks entity precision, relation correctness, completeness, consistency, provenance, and downstream QA value. Tradeoff: Intrinsic graph metrics catch construction defects but may not predict application utility. Watch for: Reporting only downstream answer scores can hide a brittle graph with unverifiable facts.}
\paragraph{MCQ-0497 [hard]} Which production warning is most directly associated with Graph quality evaluation?
\begin{enumerate}[label=\Alph*.]
\item Overwriting old facts destroys the ability to answer as-of questions.
\item Omitting labels or relationship types can trigger broad scans.
\item Allowing generated edges without source spans prevents later verification.
\item Reporting only downstream answer scores can hide a brittle graph with unverifiable facts.
\end{enumerate}
\noindent\textbf{Answer.} Reporting only downstream answer scores can hide a brittle graph with unverifiable facts.
\par\textit{Evaluation checks entity precision, relation correctness, completeness, consistency, provenance, and downstream QA value. Tradeoff: Intrinsic graph metrics catch construction defects but may not predict application utility. Watch for: Reporting only downstream answer scores can hide a brittle graph with unverifiable facts.}
\paragraph{FC-0498 [medium]} Define Graph quality evaluation in interview-ready language.
\noindent\textbf{Answer.} Evaluation checks entity precision, relation correctness, completeness, consistency, provenance, and downstream QA value.
\par\textit{Evaluation checks entity precision, relation correctness, completeness, consistency, provenance, and downstream QA value.}
\paragraph{FC-0499 [hard]} State the main tradeoff and production pitfall for Graph quality evaluation.
\noindent\textbf{Answer.} Tradeoff: Intrinsic graph metrics catch construction defects but may not predict application utility. Pitfall: Reporting only downstream answer scores can hide a brittle graph with unverifiable facts.
\par\textit{Tradeoff: Intrinsic graph metrics catch construction defects but may not predict application utility. Pitfall: Reporting only downstream answer scores can hide a brittle graph with unverifiable facts.}
\paragraph{LA-0500 [hard]} Explain Graph quality evaluation from first principles, then show how you would evaluate and operate it in production.
\noindent\textbf{Answer.} Start with the mechanism: Evaluation checks entity precision, relation correctness, completeness, consistency, provenance, and downstream QA value. Make the tradeoff explicit: Intrinsic graph metrics catch construction defects but may not predict application utility. Define workload-specific offline metrics and a latency/cost budget, test important slices, instrument the critical path, and create a rollback criterion. The production review must explicitly guard against this failure: Reporting only downstream answer scores can hide a brittle graph with unverifiable facts.
\par\textit{A strong answer connects mechanism, measurable tradeoffs, failure isolation, observability, and rollback rather than listing product features.}
\paragraph{MCQ-0501 [medium]} Which statement best describes Temporal knowledge graphs?
\begin{enumerate}[label=\Alph*.]
\item Local search expands from query-relevant entities through nearby relationships and associated text.
\item RDF triples plus vocabularies such as RDFS or OWL provide globally identified semantics and inference rules.
\item Evaluation checks entity precision, relation correctness, completeness, consistency, provenance, and downstream QA value.
\item Facts carry validity intervals, event time, and version history so queries can reason about change.
\end{enumerate}
\noindent\textbf{Answer.} Facts carry validity intervals, event time, and version history so queries can reason about change.
\par\textit{Facts carry validity intervals, event time, and version history so queries can reason about change. Tradeoff: Temporal accuracy improves, while updates, conflicts, and query logic become harder. Watch for: Overwriting old facts destroys the ability to answer as-of questions.}
\paragraph{MCQ-0502 [hard]} What is the central engineering tradeoff of Temporal knowledge graphs?
\begin{enumerate}[label=\Alph*.]
\item Intrinsic graph metrics catch construction defects but may not predict application utility.
\item Temporal accuracy improves, while updates, conflicts, and query logic become harder.
\item It supports entity-centric multi-hop questions but depends on extraction and linking quality.
\item Formal interoperability grows, with more modeling and reasoning complexity.
\end{enumerate}
\noindent\textbf{Answer.} Temporal accuracy improves, while updates, conflicts, and query logic become harder.
\par\textit{Facts carry validity intervals, event time, and version history so queries can reason about change. Tradeoff: Temporal accuracy improves, while updates, conflicts, and query logic become harder. Watch for: Overwriting old facts destroys the ability to answer as-of questions.}
\paragraph{MCQ-0503 [hard]} Which failure mode should an interviewer expect you to identify for Temporal knowledge graphs?
\begin{enumerate}[label=\Alph*.]
\item Reporting only downstream answer scores can hide a brittle graph with unverifiable facts.
\item Overwriting old facts destroys the ability to answer as-of questions.
\item Confusing open-world semantics with database constraints produces invalid assumptions.
\item Unbounded neighborhood expansion floods context with high-degree hubs.
\end{enumerate}
\noindent\textbf{Answer.} Overwriting old facts destroys the ability to answer as-of questions.
\par\textit{Facts carry validity intervals, event time, and version history so queries can reason about change. Tradeoff: Temporal accuracy improves, while updates, conflicts, and query logic become harder. Watch for: Overwriting old facts destroys the ability to answer as-of questions.}
\paragraph{MCQ-0504 [medium]} A design review is specifically concerned with this risk: Overwriting old facts destroys the ability to answer as-of questions. Which topic should be reviewed first?
\begin{enumerate}[label=\Alph*.]
\item GraphRAG local search
\item Graph quality evaluation
\item RDF and ontologies
\item Temporal knowledge graphs
\end{enumerate}
\noindent\textbf{Answer.} Temporal knowledge graphs
\par\textit{Facts carry validity intervals, event time, and version history so queries can reason about change. Tradeoff: Temporal accuracy improves, while updates, conflicts, and query logic become harder. Watch for: Overwriting old facts destroys the ability to answer as-of questions.}
\paragraph{MCQ-0505 [hard]} Which design note most accurately balances the benefits and costs of Temporal knowledge graphs?
\begin{enumerate}[label=\Alph*.]
\item It supports entity-centric multi-hop questions but depends on extraction and linking quality.
\item Temporal accuracy improves, while updates, conflicts, and query logic become harder.
\item Formal interoperability grows, with more modeling and reasoning complexity.
\item Intrinsic graph metrics catch construction defects but may not predict application utility.
\end{enumerate}
\noindent\textbf{Answer.} Temporal accuracy improves, while updates, conflicts, and query logic become harder.
\par\textit{Facts carry validity intervals, event time, and version history so queries can reason about change. Tradeoff: Temporal accuracy improves, while updates, conflicts, and query logic become harder. Watch for: Overwriting old facts destroys the ability to answer as-of questions.}
\paragraph{MCQ-0506 [medium]} Which explanation would be strongest in a system-design interview about Temporal knowledge graphs?
\begin{enumerate}[label=\Alph*.]
\item Evaluation checks entity precision, relation correctness, completeness, consistency, provenance, and downstream QA value.
\item RDF triples plus vocabularies such as RDFS or OWL provide globally identified semantics and inference rules.
\item Local search expands from query-relevant entities through nearby relationships and associated text.
\item Facts carry validity intervals, event time, and version history so queries can reason about change.
\end{enumerate}
\noindent\textbf{Answer.} Facts carry validity intervals, event time, and version history so queries can reason about change.
\par\textit{Facts carry validity intervals, event time, and version history so queries can reason about change. Tradeoff: Temporal accuracy improves, while updates, conflicts, and query logic become harder. Watch for: Overwriting old facts destroys the ability to answer as-of questions.}
\paragraph{MCQ-0507 [hard]} Which production warning is most directly associated with Temporal knowledge graphs?
\begin{enumerate}[label=\Alph*.]
\item Unbounded neighborhood expansion floods context with high-degree hubs.
\item Overwriting old facts destroys the ability to answer as-of questions.
\item Confusing open-world semantics with database constraints produces invalid assumptions.
\item Reporting only downstream answer scores can hide a brittle graph with unverifiable facts.
\end{enumerate}
\noindent\textbf{Answer.} Overwriting old facts destroys the ability to answer as-of questions.
\par\textit{Facts carry validity intervals, event time, and version history so queries can reason about change. Tradeoff: Temporal accuracy improves, while updates, conflicts, and query logic become harder. Watch for: Overwriting old facts destroys the ability to answer as-of questions.}
\paragraph{FC-0508 [medium]} Define Temporal knowledge graphs in interview-ready language.
\noindent\textbf{Answer.} Facts carry validity intervals, event time, and version history so queries can reason about change.
\par\textit{Facts carry validity intervals, event time, and version history so queries can reason about change.}
\paragraph{FC-0509 [hard]} State the main tradeoff and production pitfall for Temporal knowledge graphs.
\noindent\textbf{Answer.} Tradeoff: Temporal accuracy improves, while updates, conflicts, and query logic become harder. Pitfall: Overwriting old facts destroys the ability to answer as-of questions.
\par\textit{Tradeoff: Temporal accuracy improves, while updates, conflicts, and query logic become harder. Pitfall: Overwriting old facts destroys the ability to answer as-of questions.}
\paragraph{LA-0510 [hard]} Explain Temporal knowledge graphs from first principles, then show how you would evaluate and operate it in production.
\noindent\textbf{Answer.} Start with the mechanism: Facts carry validity intervals, event time, and version history so queries can reason about change. Make the tradeoff explicit: Temporal accuracy improves, while updates, conflicts, and query logic become harder. Define workload-specific offline metrics and a latency/cost budget, test important slices, instrument the critical path, and create a rollback criterion. The production review must explicitly guard against this failure: Overwriting old facts destroys the ability to answer as-of questions.
\par\textit{A strong answer connects mechanism, measurable tradeoffs, failure isolation, observability, and rollback rather than listing product features.}
\subsection{Agents, MCP, and Control}
\paragraph{MCQ-0511 [medium]} Which statement best describes Agent loop?
\begin{enumerate}[label=\Alph*.]
\item Context engineering selects instructions, memory, evidence, tools, schemas, and state for each model turn.
\item Explicit states and transitions constrain allowed agent behavior and make recovery inspectable.
\item A sandbox isolates untrusted code with resource limits, scoped files, network policy, timeouts, and audit logs.
\item An agent repeatedly observes state, chooses an action, invokes a tool, and incorporates the result until stopping.
\end{enumerate}
\noindent\textbf{Answer.} An agent repeatedly observes state, chooses an action, invokes a tool, and incorporates the result until stopping.
\par\textit{An agent repeatedly observes state, chooses an action, invokes a tool, and incorporates the result until stopping. Tradeoff: Flexibility handles open-ended tasks, but loop length and behavior are less predictable than workflows. Watch for: Omitting hard budgets and stop conditions creates runaway calls.}
\paragraph{MCQ-0512 [hard]} What is the central engineering tradeoff of Agent loop?
\begin{enumerate}[label=\Alph*.]
\item Isolation enables capability while imposing startup and compatibility overhead.
\item Targeted context improves reliability and cost, but selection logic becomes core application code.
\item Flexibility handles open-ended tasks, but loop length and behavior are less predictable than workflows.
\item Predictability and testability improve at the cost of less open-ended autonomy.
\end{enumerate}
\noindent\textbf{Answer.} Flexibility handles open-ended tasks, but loop length and behavior are less predictable than workflows.
\par\textit{An agent repeatedly observes state, chooses an action, invokes a tool, and incorporates the result until stopping. Tradeoff: Flexibility handles open-ended tasks, but loop length and behavior are less predictable than workflows. Watch for: Omitting hard budgets and stop conditions creates runaway calls.}
\paragraph{MCQ-0513 [hard]} Which failure mode should an interviewer expect you to identify for Agent loop?
\begin{enumerate}[label=\Alph*.]
\item Omitting hard budgets and stop conditions creates runaway calls.
\item Encoding side effects inside transition guards breaks replay safety.
\item A container without kernel, credential, and network boundaries is not a sufficient sandbox.
\item Appending everything causes attention dilution and instruction conflicts.
\end{enumerate}
\noindent\textbf{Answer.} Omitting hard budgets and stop conditions creates runaway calls.
\par\textit{An agent repeatedly observes state, chooses an action, invokes a tool, and incorporates the result until stopping. Tradeoff: Flexibility handles open-ended tasks, but loop length and behavior are less predictable than workflows. Watch for: Omitting hard budgets and stop conditions creates runaway calls.}
\paragraph{MCQ-0514 [medium]} A design review is specifically concerned with this risk: Omitting hard budgets and stop conditions creates runaway calls. Which topic should be reviewed first?
\begin{enumerate}[label=\Alph*.]
\item Agent loop
\item State machines
\item Context engineering
\item Sandbox runs
\end{enumerate}
\noindent\textbf{Answer.} Agent loop
\par\textit{An agent repeatedly observes state, chooses an action, invokes a tool, and incorporates the result until stopping. Tradeoff: Flexibility handles open-ended tasks, but loop length and behavior are less predictable than workflows. Watch for: Omitting hard budgets and stop conditions creates runaway calls.}
\paragraph{MCQ-0515 [hard]} Which design note most accurately balances the benefits and costs of Agent loop?
\begin{enumerate}[label=\Alph*.]
\item Flexibility handles open-ended tasks, but loop length and behavior are less predictable than workflows.
\item Isolation enables capability while imposing startup and compatibility overhead.
\item Targeted context improves reliability and cost, but selection logic becomes core application code.
\item Predictability and testability improve at the cost of less open-ended autonomy.
\end{enumerate}
\noindent\textbf{Answer.} Flexibility handles open-ended tasks, but loop length and behavior are less predictable than workflows.
\par\textit{An agent repeatedly observes state, chooses an action, invokes a tool, and incorporates the result until stopping. Tradeoff: Flexibility handles open-ended tasks, but loop length and behavior are less predictable than workflows. Watch for: Omitting hard budgets and stop conditions creates runaway calls.}
\paragraph{MCQ-0516 [medium]} Which explanation would be strongest in a system-design interview about Agent loop?
\begin{enumerate}[label=\Alph*.]
\item A sandbox isolates untrusted code with resource limits, scoped files, network policy, timeouts, and audit logs.
\item Explicit states and transitions constrain allowed agent behavior and make recovery inspectable.
\item Context engineering selects instructions, memory, evidence, tools, schemas, and state for each model turn.
\item An agent repeatedly observes state, chooses an action, invokes a tool, and incorporates the result until stopping.
\end{enumerate}
\noindent\textbf{Answer.} An agent repeatedly observes state, chooses an action, invokes a tool, and incorporates the result until stopping.
\par\textit{An agent repeatedly observes state, chooses an action, invokes a tool, and incorporates the result until stopping. Tradeoff: Flexibility handles open-ended tasks, but loop length and behavior are less predictable than workflows. Watch for: Omitting hard budgets and stop conditions creates runaway calls.}
\paragraph{MCQ-0517 [hard]} Which production warning is most directly associated with Agent loop?
\begin{enumerate}[label=\Alph*.]
\item Omitting hard budgets and stop conditions creates runaway calls.
\item Appending everything causes attention dilution and instruction conflicts.
\item Encoding side effects inside transition guards breaks replay safety.
\item A container without kernel, credential, and network boundaries is not a sufficient sandbox.
\end{enumerate}
\noindent\textbf{Answer.} Omitting hard budgets and stop conditions creates runaway calls.
\par\textit{An agent repeatedly observes state, chooses an action, invokes a tool, and incorporates the result until stopping. Tradeoff: Flexibility handles open-ended tasks, but loop length and behavior are less predictable than workflows. Watch for: Omitting hard budgets and stop conditions creates runaway calls.}
\paragraph{FC-0518 [medium]} Define Agent loop in interview-ready language.
\noindent\textbf{Answer.} An agent repeatedly observes state, chooses an action, invokes a tool, and incorporates the result until stopping.
\par\textit{An agent repeatedly observes state, chooses an action, invokes a tool, and incorporates the result until stopping.}
\paragraph{FC-0519 [hard]} State the main tradeoff and production pitfall for Agent loop.
\noindent\textbf{Answer.} Tradeoff: Flexibility handles open-ended tasks, but loop length and behavior are less predictable than workflows. Pitfall: Omitting hard budgets and stop conditions creates runaway calls.
\par\textit{Tradeoff: Flexibility handles open-ended tasks, but loop length and behavior are less predictable than workflows. Pitfall: Omitting hard budgets and stop conditions creates runaway calls.}
\paragraph{LA-0520 [hard]} Explain Agent loop from first principles, then show how you would evaluate and operate it in production.
\noindent\textbf{Answer.} Start with the mechanism: An agent repeatedly observes state, chooses an action, invokes a tool, and incorporates the result until stopping. Make the tradeoff explicit: Flexibility handles open-ended tasks, but loop length and behavior are less predictable than workflows. Define workload-specific offline metrics and a latency/cost budget, test important slices, instrument the critical path, and create a rollback criterion. The production review must explicitly guard against this failure: Omitting hard budgets and stop conditions creates runaway calls.
\par\textit{A strong answer connects mechanism, measurable tradeoffs, failure isolation, observability, and rollback rather than listing product features.}
\paragraph{MCQ-0521 [medium]} Which statement best describes ReAct?
\begin{enumerate}[label=\Alph*.]
\item ReAct interleaves reasoning with actions and observations to update plans using external evidence.
\item High-risk actions pause for review, edit, approval, rejection, or direct human response.
\item Models emit structured arguments for declared functions that a host validates and executes.
\item Working, episodic, semantic, and procedural memories preserve relevant state beyond a single model call.
\end{enumerate}
\noindent\textbf{Answer.} ReAct interleaves reasoning with actions and observations to update plans using external evidence.
\par\textit{ReAct interleaves reasoning with actions and observations to update plans using external evidence. Tradeoff: Tool feedback can correct model knowledge, while longer trajectories add latency and error opportunities. Watch for: Treating every thought as reliable state allows early mistakes to compound.}
\paragraph{MCQ-0522 [hard]} What is the central engineering tradeoff of ReAct?
\begin{enumerate}[label=\Alph*.]
\item Memory supports continuity but retrieval, deletion, privacy, and stale facts become product concerns.
\item Oversight reduces irreversible mistakes but adds latency and reviewer load.
\item It narrows the action interface but does not make proposed actions safe or correct.
\item Tool feedback can correct model knowledge, while longer trajectories add latency and error opportunities.
\end{enumerate}
\noindent\textbf{Answer.} Tool feedback can correct model knowledge, while longer trajectories add latency and error opportunities.
\par\textit{ReAct interleaves reasoning with actions and observations to update plans using external evidence. Tradeoff: Tool feedback can correct model knowledge, while longer trajectories add latency and error opportunities. Watch for: Treating every thought as reliable state allows early mistakes to compound.}
\paragraph{MCQ-0523 [hard]} Which failure mode should an interviewer expect you to identify for ReAct?
\begin{enumerate}[label=\Alph*.]
\item Treating every thought as reliable state allows early mistakes to compound.
\item Executing model-produced arguments without schema and authorization checks invites damage.
\item Writing every interaction to long-term memory creates noise and privacy exposure.
\item Asking for approval after an action has already caused side effects is theater.
\end{enumerate}
\noindent\textbf{Answer.} Treating every thought as reliable state allows early mistakes to compound.
\par\textit{ReAct interleaves reasoning with actions and observations to update plans using external evidence. Tradeoff: Tool feedback can correct model knowledge, while longer trajectories add latency and error opportunities. Watch for: Treating every thought as reliable state allows early mistakes to compound.}
\paragraph{MCQ-0524 [medium]} A design review is specifically concerned with this risk: Treating every thought as reliable state allows early mistakes to compound. Which topic should be reviewed first?
\begin{enumerate}[label=\Alph*.]
\item Agent memory
\item Tool calling
\item ReAct
\item Human in the loop
\end{enumerate}
\noindent\textbf{Answer.} ReAct
\par\textit{ReAct interleaves reasoning with actions and observations to update plans using external evidence. Tradeoff: Tool feedback can correct model knowledge, while longer trajectories add latency and error opportunities. Watch for: Treating every thought as reliable state allows early mistakes to compound.}
\paragraph{MCQ-0525 [hard]} Which design note most accurately balances the benefits and costs of ReAct?
\begin{enumerate}[label=\Alph*.]
\item Memory supports continuity but retrieval, deletion, privacy, and stale facts become product concerns.
\item Oversight reduces irreversible mistakes but adds latency and reviewer load.
\item It narrows the action interface but does not make proposed actions safe or correct.
\item Tool feedback can correct model knowledge, while longer trajectories add latency and error opportunities.
\end{enumerate}
\noindent\textbf{Answer.} Tool feedback can correct model knowledge, while longer trajectories add latency and error opportunities.
\par\textit{ReAct interleaves reasoning with actions and observations to update plans using external evidence. Tradeoff: Tool feedback can correct model knowledge, while longer trajectories add latency and error opportunities. Watch for: Treating every thought as reliable state allows early mistakes to compound.}
\paragraph{MCQ-0526 [medium]} Which explanation would be strongest in a system-design interview about ReAct?
\begin{enumerate}[label=\Alph*.]
\item High-risk actions pause for review, edit, approval, rejection, or direct human response.
\item Models emit structured arguments for declared functions that a host validates and executes.
\item Working, episodic, semantic, and procedural memories preserve relevant state beyond a single model call.
\item ReAct interleaves reasoning with actions and observations to update plans using external evidence.
\end{enumerate}
\noindent\textbf{Answer.} ReAct interleaves reasoning with actions and observations to update plans using external evidence.
\par\textit{ReAct interleaves reasoning with actions and observations to update plans using external evidence. Tradeoff: Tool feedback can correct model knowledge, while longer trajectories add latency and error opportunities. Watch for: Treating every thought as reliable state allows early mistakes to compound.}
\paragraph{MCQ-0527 [hard]} Which production warning is most directly associated with ReAct?
\begin{enumerate}[label=\Alph*.]
\item Writing every interaction to long-term memory creates noise and privacy exposure.
\item Executing model-produced arguments without schema and authorization checks invites damage.
\item Treating every thought as reliable state allows early mistakes to compound.
\item Asking for approval after an action has already caused side effects is theater.
\end{enumerate}
\noindent\textbf{Answer.} Treating every thought as reliable state allows early mistakes to compound.
\par\textit{ReAct interleaves reasoning with actions and observations to update plans using external evidence. Tradeoff: Tool feedback can correct model knowledge, while longer trajectories add latency and error opportunities. Watch for: Treating every thought as reliable state allows early mistakes to compound.}
\paragraph{FC-0528 [medium]} Define ReAct in interview-ready language.
\noindent\textbf{Answer.} ReAct interleaves reasoning with actions and observations to update plans using external evidence.
\par\textit{ReAct interleaves reasoning with actions and observations to update plans using external evidence.}
\paragraph{FC-0529 [hard]} State the main tradeoff and production pitfall for ReAct.
\noindent\textbf{Answer.} Tradeoff: Tool feedback can correct model knowledge, while longer trajectories add latency and error opportunities. Pitfall: Treating every thought as reliable state allows early mistakes to compound.
\par\textit{Tradeoff: Tool feedback can correct model knowledge, while longer trajectories add latency and error opportunities. Pitfall: Treating every thought as reliable state allows early mistakes to compound.}
\paragraph{LA-0530 [hard]} Explain ReAct from first principles, then show how you would evaluate and operate it in production.
\noindent\textbf{Answer.} Start with the mechanism: ReAct interleaves reasoning with actions and observations to update plans using external evidence. Make the tradeoff explicit: Tool feedback can correct model knowledge, while longer trajectories add latency and error opportunities. Define workload-specific offline metrics and a latency/cost budget, test important slices, instrument the critical path, and create a rollback criterion. The production review must explicitly guard against this failure: Treating every thought as reliable state allows early mistakes to compound.
\par\textit{A strong answer connects mechanism, measurable tradeoffs, failure isolation, observability, and rollback rather than listing product features.}
\paragraph{MCQ-0531 [medium]} Which statement best describes Planning?
\begin{enumerate}[label=\Alph*.]
\item A handoff transfers a task and structured context to a specialized agent or deterministic service.
\item Planning decomposes goals into ordered, conditional, or revisable steps before or during execution.
\item Context engineering selects instructions, memory, evidence, tools, schemas, and state for each model turn.
\item An idempotent tool can be retried with the same operation key without duplicating the intended effect.
\end{enumerate}
\noindent\textbf{Answer.} Planning decomposes goals into ordered, conditional, or revisable steps before or during execution.
\par\textit{Planning decomposes goals into ordered, conditional, or revisable steps before or during execution. Tradeoff: Upfront plans improve structure but become stale in uncertain environments. Watch for: Forcing a long immutable plan prevents recovery after new observations.}
\paragraph{MCQ-0532 [hard]} What is the central engineering tradeoff of Planning?
\begin{enumerate}[label=\Alph*.]
\item Upfront plans improve structure but become stale in uncertain environments.
\item Reliable retries require deduplication state and carefully defined semantics.
\item Specialization reduces prompt overload but adds routing and ownership ambiguity.
\item Targeted context improves reliability and cost, but selection logic becomes core application code.
\end{enumerate}
\noindent\textbf{Answer.} Upfront plans improve structure but become stale in uncertain environments.
\par\textit{Planning decomposes goals into ordered, conditional, or revisable steps before or during execution. Tradeoff: Upfront plans improve structure but become stale in uncertain environments. Watch for: Forcing a long immutable plan prevents recovery after new observations.}
\paragraph{MCQ-0533 [hard]} Which failure mode should an interviewer expect you to identify for Planning?
\begin{enumerate}[label=\Alph*.]
\item Appending everything causes attention dilution and instruction conflicts.
\item Assuming HTTP method alone guarantees business idempotency causes duplicate orders or messages.
\item Passing the entire transcript leaks irrelevant data and confuses the recipient.
\item Forcing a long immutable plan prevents recovery after new observations.
\end{enumerate}
\noindent\textbf{Answer.} Forcing a long immutable plan prevents recovery after new observations.
\par\textit{Planning decomposes goals into ordered, conditional, or revisable steps before or during execution. Tradeoff: Upfront plans improve structure but become stale in uncertain environments. Watch for: Forcing a long immutable plan prevents recovery after new observations.}
\paragraph{MCQ-0534 [medium]} A design review is specifically concerned with this risk: Forcing a long immutable plan prevents recovery after new observations. Which topic should be reviewed first?
\begin{enumerate}[label=\Alph*.]
\item Agent handoffs
\item Context engineering
\item Planning
\item Idempotent tools
\end{enumerate}
\noindent\textbf{Answer.} Planning
\par\textit{Planning decomposes goals into ordered, conditional, or revisable steps before or during execution. Tradeoff: Upfront plans improve structure but become stale in uncertain environments. Watch for: Forcing a long immutable plan prevents recovery after new observations.}
\paragraph{MCQ-0535 [hard]} Which design note most accurately balances the benefits and costs of Planning?
\begin{enumerate}[label=\Alph*.]
\item Upfront plans improve structure but become stale in uncertain environments.
\item Specialization reduces prompt overload but adds routing and ownership ambiguity.
\item Reliable retries require deduplication state and carefully defined semantics.
\item Targeted context improves reliability and cost, but selection logic becomes core application code.
\end{enumerate}
\noindent\textbf{Answer.} Upfront plans improve structure but become stale in uncertain environments.
\par\textit{Planning decomposes goals into ordered, conditional, or revisable steps before or during execution. Tradeoff: Upfront plans improve structure but become stale in uncertain environments. Watch for: Forcing a long immutable plan prevents recovery after new observations.}
\paragraph{MCQ-0536 [medium]} Which explanation would be strongest in a system-design interview about Planning?
\begin{enumerate}[label=\Alph*.]
\item An idempotent tool can be retried with the same operation key without duplicating the intended effect.
\item A handoff transfers a task and structured context to a specialized agent or deterministic service.
\item Planning decomposes goals into ordered, conditional, or revisable steps before or during execution.
\item Context engineering selects instructions, memory, evidence, tools, schemas, and state for each model turn.
\end{enumerate}
\noindent\textbf{Answer.} Planning decomposes goals into ordered, conditional, or revisable steps before or during execution.
\par\textit{Planning decomposes goals into ordered, conditional, or revisable steps before or during execution. Tradeoff: Upfront plans improve structure but become stale in uncertain environments. Watch for: Forcing a long immutable plan prevents recovery after new observations.}
\paragraph{MCQ-0537 [hard]} Which production warning is most directly associated with Planning?
\begin{enumerate}[label=\Alph*.]
\item Passing the entire transcript leaks irrelevant data and confuses the recipient.
\item Appending everything causes attention dilution and instruction conflicts.
\item Forcing a long immutable plan prevents recovery after new observations.
\item Assuming HTTP method alone guarantees business idempotency causes duplicate orders or messages.
\end{enumerate}
\noindent\textbf{Answer.} Forcing a long immutable plan prevents recovery after new observations.
\par\textit{Planning decomposes goals into ordered, conditional, or revisable steps before or during execution. Tradeoff: Upfront plans improve structure but become stale in uncertain environments. Watch for: Forcing a long immutable plan prevents recovery after new observations.}
\paragraph{FC-0538 [medium]} Define Planning in interview-ready language.
\noindent\textbf{Answer.} Planning decomposes goals into ordered, conditional, or revisable steps before or during execution.
\par\textit{Planning decomposes goals into ordered, conditional, or revisable steps before or during execution.}
\paragraph{FC-0539 [hard]} State the main tradeoff and production pitfall for Planning.
\noindent\textbf{Answer.} Tradeoff: Upfront plans improve structure but become stale in uncertain environments. Pitfall: Forcing a long immutable plan prevents recovery after new observations.
\par\textit{Tradeoff: Upfront plans improve structure but become stale in uncertain environments. Pitfall: Forcing a long immutable plan prevents recovery after new observations.}
\paragraph{LA-0540 [hard]} Explain Planning from first principles, then show how you would evaluate and operate it in production.
\noindent\textbf{Answer.} Start with the mechanism: Planning decomposes goals into ordered, conditional, or revisable steps before or during execution. Make the tradeoff explicit: Upfront plans improve structure but become stale in uncertain environments. Define workload-specific offline metrics and a latency/cost budget, test important slices, instrument the critical path, and create a rollback criterion. The production review must explicitly guard against this failure: Forcing a long immutable plan prevents recovery after new observations.
\par\textit{A strong answer connects mechanism, measurable tradeoffs, failure isolation, observability, and rollback rather than listing product features.}
\paragraph{MCQ-0541 [medium]} Which statement best describes Reflection?
\begin{enumerate}[label=\Alph*.]
\item High-risk actions pause for review, edit, approval, rejection, or direct human response.
\item A handoff transfers a task and structured context to a specialized agent or deterministic service.
\item Models emit structured arguments for declared functions that a host validates and executes.
\item Reflection critiques outcomes or trajectories and stores actionable lessons for another attempt.
\end{enumerate}
\noindent\textbf{Answer.} Reflection critiques outcomes or trajectories and stores actionable lessons for another attempt.
\par\textit{Reflection critiques outcomes or trajectories and stores actionable lessons for another attempt. Tradeoff: It can improve recovery without weight updates, while self-critique may repeat the same bias. Watch for: Saving vague reflections adds context without changing behavior.}
\paragraph{MCQ-0542 [hard]} What is the central engineering tradeoff of Reflection?
\begin{enumerate}[label=\Alph*.]
\item It narrows the action interface but does not make proposed actions safe or correct.
\item Specialization reduces prompt overload but adds routing and ownership ambiguity.
\item Oversight reduces irreversible mistakes but adds latency and reviewer load.
\item It can improve recovery without weight updates, while self-critique may repeat the same bias.
\end{enumerate}
\noindent\textbf{Answer.} It can improve recovery without weight updates, while self-critique may repeat the same bias.
\par\textit{Reflection critiques outcomes or trajectories and stores actionable lessons for another attempt. Tradeoff: It can improve recovery without weight updates, while self-critique may repeat the same bias. Watch for: Saving vague reflections adds context without changing behavior.}
\paragraph{MCQ-0543 [hard]} Which failure mode should an interviewer expect you to identify for Reflection?
\begin{enumerate}[label=\Alph*.]
\item Saving vague reflections adds context without changing behavior.
\item Asking for approval after an action has already caused side effects is theater.
\item Passing the entire transcript leaks irrelevant data and confuses the recipient.
\item Executing model-produced arguments without schema and authorization checks invites damage.
\end{enumerate}
\noindent\textbf{Answer.} Saving vague reflections adds context without changing behavior.
\par\textit{Reflection critiques outcomes or trajectories and stores actionable lessons for another attempt. Tradeoff: It can improve recovery without weight updates, while self-critique may repeat the same bias. Watch for: Saving vague reflections adds context without changing behavior.}
\paragraph{MCQ-0544 [medium]} A design review is specifically concerned with this risk: Saving vague reflections adds context without changing behavior. Which topic should be reviewed first?
\begin{enumerate}[label=\Alph*.]
\item Agent handoffs
\item Tool calling
\item Human in the loop
\item Reflection
\end{enumerate}
\noindent\textbf{Answer.} Reflection
\par\textit{Reflection critiques outcomes or trajectories and stores actionable lessons for another attempt. Tradeoff: It can improve recovery without weight updates, while self-critique may repeat the same bias. Watch for: Saving vague reflections adds context without changing behavior.}
\paragraph{MCQ-0545 [hard]} Which design note most accurately balances the benefits and costs of Reflection?
\begin{enumerate}[label=\Alph*.]
\item Oversight reduces irreversible mistakes but adds latency and reviewer load.
\item It narrows the action interface but does not make proposed actions safe or correct.
\item It can improve recovery without weight updates, while self-critique may repeat the same bias.
\item Specialization reduces prompt overload but adds routing and ownership ambiguity.
\end{enumerate}
\noindent\textbf{Answer.} It can improve recovery without weight updates, while self-critique may repeat the same bias.
\par\textit{Reflection critiques outcomes or trajectories and stores actionable lessons for another attempt. Tradeoff: It can improve recovery without weight updates, while self-critique may repeat the same bias. Watch for: Saving vague reflections adds context without changing behavior.}
\paragraph{MCQ-0546 [medium]} Which explanation would be strongest in a system-design interview about Reflection?
\begin{enumerate}[label=\Alph*.]
\item High-risk actions pause for review, edit, approval, rejection, or direct human response.
\item Models emit structured arguments for declared functions that a host validates and executes.
\item A handoff transfers a task and structured context to a specialized agent or deterministic service.
\item Reflection critiques outcomes or trajectories and stores actionable lessons for another attempt.
\end{enumerate}
\noindent\textbf{Answer.} Reflection critiques outcomes or trajectories and stores actionable lessons for another attempt.
\par\textit{Reflection critiques outcomes or trajectories and stores actionable lessons for another attempt. Tradeoff: It can improve recovery without weight updates, while self-critique may repeat the same bias. Watch for: Saving vague reflections adds context without changing behavior.}
\paragraph{MCQ-0547 [hard]} Which production warning is most directly associated with Reflection?
\begin{enumerate}[label=\Alph*.]
\item Saving vague reflections adds context without changing behavior.
\item Executing model-produced arguments without schema and authorization checks invites damage.
\item Asking for approval after an action has already caused side effects is theater.
\item Passing the entire transcript leaks irrelevant data and confuses the recipient.
\end{enumerate}
\noindent\textbf{Answer.} Saving vague reflections adds context without changing behavior.
\par\textit{Reflection critiques outcomes or trajectories and stores actionable lessons for another attempt. Tradeoff: It can improve recovery without weight updates, while self-critique may repeat the same bias. Watch for: Saving vague reflections adds context without changing behavior.}
\paragraph{FC-0548 [medium]} Define Reflection in interview-ready language.
\noindent\textbf{Answer.} Reflection critiques outcomes or trajectories and stores actionable lessons for another attempt.
\par\textit{Reflection critiques outcomes or trajectories and stores actionable lessons for another attempt.}
\paragraph{FC-0549 [hard]} State the main tradeoff and production pitfall for Reflection.
\noindent\textbf{Answer.} Tradeoff: It can improve recovery without weight updates, while self-critique may repeat the same bias. Pitfall: Saving vague reflections adds context without changing behavior.
\par\textit{Tradeoff: It can improve recovery without weight updates, while self-critique may repeat the same bias. Pitfall: Saving vague reflections adds context without changing behavior.}
\paragraph{LA-0550 [hard]} Explain Reflection from first principles, then show how you would evaluate and operate it in production.
\noindent\textbf{Answer.} Start with the mechanism: Reflection critiques outcomes or trajectories and stores actionable lessons for another attempt. Make the tradeoff explicit: It can improve recovery without weight updates, while self-critique may repeat the same bias. Define workload-specific offline metrics and a latency/cost budget, test important slices, instrument the critical path, and create a rollback criterion. The production review must explicitly guard against this failure: Saving vague reflections adds context without changing behavior.
\par\textit{A strong answer connects mechanism, measurable tradeoffs, failure isolation, observability, and rollback rather than listing product features.}
\paragraph{MCQ-0551 [medium]} Which statement best describes Tool calling?
\begin{enumerate}[label=\Alph*.]
\item Models emit structured arguments for declared functions that a host validates and executes.
\item Working, episodic, semantic, and procedural memories preserve relevant state beyond a single model call.
\item A sandbox isolates untrusted code with resource limits, scoped files, network policy, timeouts, and audit logs.
\item Planning decomposes goals into ordered, conditional, or revisable steps before or during execution.
\end{enumerate}
\noindent\textbf{Answer.} Models emit structured arguments for declared functions that a host validates and executes.
\par\textit{Models emit structured arguments for declared functions that a host validates and executes. Tradeoff: It narrows the action interface but does not make proposed actions safe or correct. Watch for: Executing model-produced arguments without schema and authorization checks invites damage.}
\paragraph{MCQ-0552 [hard]} What is the central engineering tradeoff of Tool calling?
\begin{enumerate}[label=\Alph*.]
\item Memory supports continuity but retrieval, deletion, privacy, and stale facts become product concerns.
\item Upfront plans improve structure but become stale in uncertain environments.
\item Isolation enables capability while imposing startup and compatibility overhead.
\item It narrows the action interface but does not make proposed actions safe or correct.
\end{enumerate}
\noindent\textbf{Answer.} It narrows the action interface but does not make proposed actions safe or correct.
\par\textit{Models emit structured arguments for declared functions that a host validates and executes. Tradeoff: It narrows the action interface but does not make proposed actions safe or correct. Watch for: Executing model-produced arguments without schema and authorization checks invites damage.}
\paragraph{MCQ-0553 [hard]} Which failure mode should an interviewer expect you to identify for Tool calling?
\begin{enumerate}[label=\Alph*.]
\item A container without kernel, credential, and network boundaries is not a sufficient sandbox.
\item Executing model-produced arguments without schema and authorization checks invites damage.
\item Forcing a long immutable plan prevents recovery after new observations.
\item Writing every interaction to long-term memory creates noise and privacy exposure.
\end{enumerate}
\noindent\textbf{Answer.} Executing model-produced arguments without schema and authorization checks invites damage.
\par\textit{Models emit structured arguments for declared functions that a host validates and executes. Tradeoff: It narrows the action interface but does not make proposed actions safe or correct. Watch for: Executing model-produced arguments without schema and authorization checks invites damage.}
\paragraph{MCQ-0554 [medium]} A design review is specifically concerned with this risk: Executing model-produced arguments without schema and authorization checks invites damage. Which topic should be reviewed first?
\begin{enumerate}[label=\Alph*.]
\item Sandbox runs
\item Planning
\item Agent memory
\item Tool calling
\end{enumerate}
\noindent\textbf{Answer.} Tool calling
\par\textit{Models emit structured arguments for declared functions that a host validates and executes. Tradeoff: It narrows the action interface but does not make proposed actions safe or correct. Watch for: Executing model-produced arguments without schema and authorization checks invites damage.}
\paragraph{MCQ-0555 [hard]} Which design note most accurately balances the benefits and costs of Tool calling?
\begin{enumerate}[label=\Alph*.]
\item It narrows the action interface but does not make proposed actions safe or correct.
\item Upfront plans improve structure but become stale in uncertain environments.
\item Isolation enables capability while imposing startup and compatibility overhead.
\item Memory supports continuity but retrieval, deletion, privacy, and stale facts become product concerns.
\end{enumerate}
\noindent\textbf{Answer.} It narrows the action interface but does not make proposed actions safe or correct.
\par\textit{Models emit structured arguments for declared functions that a host validates and executes. Tradeoff: It narrows the action interface but does not make proposed actions safe or correct. Watch for: Executing model-produced arguments without schema and authorization checks invites damage.}
\paragraph{MCQ-0556 [medium]} Which explanation would be strongest in a system-design interview about Tool calling?
\begin{enumerate}[label=\Alph*.]
\item Planning decomposes goals into ordered, conditional, or revisable steps before or during execution.
\item Models emit structured arguments for declared functions that a host validates and executes.
\item A sandbox isolates untrusted code with resource limits, scoped files, network policy, timeouts, and audit logs.
\item Working, episodic, semantic, and procedural memories preserve relevant state beyond a single model call.
\end{enumerate}
\noindent\textbf{Answer.} Models emit structured arguments for declared functions that a host validates and executes.
\par\textit{Models emit structured arguments for declared functions that a host validates and executes. Tradeoff: It narrows the action interface but does not make proposed actions safe or correct. Watch for: Executing model-produced arguments without schema and authorization checks invites damage.}
\paragraph{MCQ-0557 [hard]} Which production warning is most directly associated with Tool calling?
\begin{enumerate}[label=\Alph*.]
\item Writing every interaction to long-term memory creates noise and privacy exposure.
\item Forcing a long immutable plan prevents recovery after new observations.
\item A container without kernel, credential, and network boundaries is not a sufficient sandbox.
\item Executing model-produced arguments without schema and authorization checks invites damage.
\end{enumerate}
\noindent\textbf{Answer.} Executing model-produced arguments without schema and authorization checks invites damage.
\par\textit{Models emit structured arguments for declared functions that a host validates and executes. Tradeoff: It narrows the action interface but does not make proposed actions safe or correct. Watch for: Executing model-produced arguments without schema and authorization checks invites damage.}
\paragraph{FC-0558 [medium]} Define Tool calling in interview-ready language.
\noindent\textbf{Answer.} Models emit structured arguments for declared functions that a host validates and executes.
\par\textit{Models emit structured arguments for declared functions that a host validates and executes.}
\paragraph{FC-0559 [hard]} State the main tradeoff and production pitfall for Tool calling.
\noindent\textbf{Answer.} Tradeoff: It narrows the action interface but does not make proposed actions safe or correct. Pitfall: Executing model-produced arguments without schema and authorization checks invites damage.
\par\textit{Tradeoff: It narrows the action interface but does not make proposed actions safe or correct. Pitfall: Executing model-produced arguments without schema and authorization checks invites damage.}
\paragraph{LA-0560 [hard]} Explain Tool calling from first principles, then show how you would evaluate and operate it in production.
\noindent\textbf{Answer.} Start with the mechanism: Models emit structured arguments for declared functions that a host validates and executes. Make the tradeoff explicit: It narrows the action interface but does not make proposed actions safe or correct. Define workload-specific offline metrics and a latency/cost budget, test important slices, instrument the critical path, and create a rollback criterion. The production review must explicitly guard against this failure: Executing model-produced arguments without schema and authorization checks invites damage.
\par\textit{A strong answer connects mechanism, measurable tradeoffs, failure isolation, observability, and rollback rather than listing product features.}
\paragraph{MCQ-0561 [medium]} Which statement best describes Model Context Protocol?
\begin{enumerate}[label=\Alph*.]
\item MCP standardizes client-server discovery and invocation of tools, resources, prompts, and related capabilities.
\item ReAct interleaves reasoning with actions and observations to update plans using external evidence.
\item Models emit structured arguments for declared functions that a host validates and executes.
\item Autonomy is graduated from suggest-only through approval-gated execution to bounded autonomous operation.
\end{enumerate}
\noindent\textbf{Answer.} MCP standardizes client-server discovery and invocation of tools, resources, prompts, and related capabilities.
\par\textit{MCP standardizes client-server discovery and invocation of tools, resources, prompts, and related capabilities. Tradeoff: Interoperability improves, while trust boundaries, consent, and capability scoping remain host responsibilities. Watch for: Treating a discovered MCP server as trusted creates prompt-injection and privilege risks.}
\paragraph{MCQ-0562 [hard]} What is the central engineering tradeoff of Model Context Protocol?
\begin{enumerate}[label=\Alph*.]
\item It narrows the action interface but does not make proposed actions safe or correct.
\item Tool feedback can correct model knowledge, while longer trajectories add latency and error opportunities.
\item More autonomy reduces human effort while increasing required assurance and containment.
\item Interoperability improves, while trust boundaries, consent, and capability scoping remain host responsibilities.
\end{enumerate}
\noindent\textbf{Answer.} Interoperability improves, while trust boundaries, consent, and capability scoping remain host responsibilities.
\par\textit{MCP standardizes client-server discovery and invocation of tools, resources, prompts, and related capabilities. Tradeoff: Interoperability improves, while trust boundaries, consent, and capability scoping remain host responsibilities. Watch for: Treating a discovered MCP server as trusted creates prompt-injection and privilege risks.}
\paragraph{MCQ-0563 [hard]} Which failure mode should an interviewer expect you to identify for Model Context Protocol?
\begin{enumerate}[label=\Alph*.]
\item Treating every thought as reliable state allows early mistakes to compound.
\item Executing model-produced arguments without schema and authorization checks invites damage.
\item Assigning one global autonomy level ignores action-specific risk.
\item Treating a discovered MCP server as trusted creates prompt-injection and privilege risks.
\end{enumerate}
\noindent\textbf{Answer.} Treating a discovered MCP server as trusted creates prompt-injection and privilege risks.
\par\textit{MCP standardizes client-server discovery and invocation of tools, resources, prompts, and related capabilities. Tradeoff: Interoperability improves, while trust boundaries, consent, and capability scoping remain host responsibilities. Watch for: Treating a discovered MCP server as trusted creates prompt-injection and privilege risks.}
\paragraph{MCQ-0564 [medium]} A design review is specifically concerned with this risk: Treating a discovered MCP server as trusted creates prompt-injection and privilege risks. Which topic should be reviewed first?
\begin{enumerate}[label=\Alph*.]
\item ReAct
\item Autonomy levels
\item Model Context Protocol
\item Tool calling
\end{enumerate}
\noindent\textbf{Answer.} Model Context Protocol
\par\textit{MCP standardizes client-server discovery and invocation of tools, resources, prompts, and related capabilities. Tradeoff: Interoperability improves, while trust boundaries, consent, and capability scoping remain host responsibilities. Watch for: Treating a discovered MCP server as trusted creates prompt-injection and privilege risks.}
\paragraph{MCQ-0565 [hard]} Which design note most accurately balances the benefits and costs of Model Context Protocol?
\begin{enumerate}[label=\Alph*.]
\item Interoperability improves, while trust boundaries, consent, and capability scoping remain host responsibilities.
\item Tool feedback can correct model knowledge, while longer trajectories add latency and error opportunities.
\item More autonomy reduces human effort while increasing required assurance and containment.
\item It narrows the action interface but does not make proposed actions safe or correct.
\end{enumerate}
\noindent\textbf{Answer.} Interoperability improves, while trust boundaries, consent, and capability scoping remain host responsibilities.
\par\textit{MCP standardizes client-server discovery and invocation of tools, resources, prompts, and related capabilities. Tradeoff: Interoperability improves, while trust boundaries, consent, and capability scoping remain host responsibilities. Watch for: Treating a discovered MCP server as trusted creates prompt-injection and privilege risks.}
\paragraph{MCQ-0566 [medium]} Which explanation would be strongest in a system-design interview about Model Context Protocol?
\begin{enumerate}[label=\Alph*.]
\item Models emit structured arguments for declared functions that a host validates and executes.
\item Autonomy is graduated from suggest-only through approval-gated execution to bounded autonomous operation.
\item ReAct interleaves reasoning with actions and observations to update plans using external evidence.
\item MCP standardizes client-server discovery and invocation of tools, resources, prompts, and related capabilities.
\end{enumerate}
\noindent\textbf{Answer.} MCP standardizes client-server discovery and invocation of tools, resources, prompts, and related capabilities.
\par\textit{MCP standardizes client-server discovery and invocation of tools, resources, prompts, and related capabilities. Tradeoff: Interoperability improves, while trust boundaries, consent, and capability scoping remain host responsibilities. Watch for: Treating a discovered MCP server as trusted creates prompt-injection and privilege risks.}
\paragraph{MCQ-0567 [hard]} Which production warning is most directly associated with Model Context Protocol?
\begin{enumerate}[label=\Alph*.]
\item Assigning one global autonomy level ignores action-specific risk.
\item Executing model-produced arguments without schema and authorization checks invites damage.
\item Treating a discovered MCP server as trusted creates prompt-injection and privilege risks.
\item Treating every thought as reliable state allows early mistakes to compound.
\end{enumerate}
\noindent\textbf{Answer.} Treating a discovered MCP server as trusted creates prompt-injection and privilege risks.
\par\textit{MCP standardizes client-server discovery and invocation of tools, resources, prompts, and related capabilities. Tradeoff: Interoperability improves, while trust boundaries, consent, and capability scoping remain host responsibilities. Watch for: Treating a discovered MCP server as trusted creates prompt-injection and privilege risks.}
\paragraph{FC-0568 [medium]} Define Model Context Protocol in interview-ready language.
\noindent\textbf{Answer.} MCP standardizes client-server discovery and invocation of tools, resources, prompts, and related capabilities.
\par\textit{MCP standardizes client-server discovery and invocation of tools, resources, prompts, and related capabilities.}
\paragraph{FC-0569 [hard]} State the main tradeoff and production pitfall for Model Context Protocol.
\noindent\textbf{Answer.} Tradeoff: Interoperability improves, while trust boundaries, consent, and capability scoping remain host responsibilities. Pitfall: Treating a discovered MCP server as trusted creates prompt-injection and privilege risks.
\par\textit{Tradeoff: Interoperability improves, while trust boundaries, consent, and capability scoping remain host responsibilities. Pitfall: Treating a discovered MCP server as trusted creates prompt-injection and privilege risks.}
\paragraph{LA-0570 [hard]} Explain Model Context Protocol from first principles, then show how you would evaluate and operate it in production.
\noindent\textbf{Answer.} Start with the mechanism: MCP standardizes client-server discovery and invocation of tools, resources, prompts, and related capabilities. Make the tradeoff explicit: Interoperability improves, while trust boundaries, consent, and capability scoping remain host responsibilities. Define workload-specific offline metrics and a latency/cost budget, test important slices, instrument the critical path, and create a rollback criterion. The production review must explicitly guard against this failure: Treating a discovered MCP server as trusted creates prompt-injection and privilege risks.
\par\textit{A strong answer connects mechanism, measurable tradeoffs, failure isolation, observability, and rollback rather than listing product features.}
\paragraph{MCQ-0571 [medium]} Which statement best describes Agent memory?
\begin{enumerate}[label=\Alph*.]
\item Reflection critiques outcomes or trajectories and stores actionable lessons for another attempt.
\item High-risk actions pause for review, edit, approval, rejection, or direct human response.
\item Working, episodic, semantic, and procedural memories preserve relevant state beyond a single model call.
\item An idempotent tool can be retried with the same operation key without duplicating the intended effect.
\end{enumerate}
\noindent\textbf{Answer.} Working, episodic, semantic, and procedural memories preserve relevant state beyond a single model call.
\par\textit{Working, episodic, semantic, and procedural memories preserve relevant state beyond a single model call. Tradeoff: Memory supports continuity but retrieval, deletion, privacy, and stale facts become product concerns. Watch for: Writing every interaction to long-term memory creates noise and privacy exposure.}
\paragraph{MCQ-0572 [hard]} What is the central engineering tradeoff of Agent memory?
\begin{enumerate}[label=\Alph*.]
\item Reliable retries require deduplication state and carefully defined semantics.
\item It can improve recovery without weight updates, while self-critique may repeat the same bias.
\item Oversight reduces irreversible mistakes but adds latency and reviewer load.
\item Memory supports continuity but retrieval, deletion, privacy, and stale facts become product concerns.
\end{enumerate}
\noindent\textbf{Answer.} Memory supports continuity but retrieval, deletion, privacy, and stale facts become product concerns.
\par\textit{Working, episodic, semantic, and procedural memories preserve relevant state beyond a single model call. Tradeoff: Memory supports continuity but retrieval, deletion, privacy, and stale facts become product concerns. Watch for: Writing every interaction to long-term memory creates noise and privacy exposure.}
\paragraph{MCQ-0573 [hard]} Which failure mode should an interviewer expect you to identify for Agent memory?
\begin{enumerate}[label=\Alph*.]
\item Assuming HTTP method alone guarantees business idempotency causes duplicate orders or messages.
\item Asking for approval after an action has already caused side effects is theater.
\item Writing every interaction to long-term memory creates noise and privacy exposure.
\item Saving vague reflections adds context without changing behavior.
\end{enumerate}
\noindent\textbf{Answer.} Writing every interaction to long-term memory creates noise and privacy exposure.
\par\textit{Working, episodic, semantic, and procedural memories preserve relevant state beyond a single model call. Tradeoff: Memory supports continuity but retrieval, deletion, privacy, and stale facts become product concerns. Watch for: Writing every interaction to long-term memory creates noise and privacy exposure.}
\paragraph{MCQ-0574 [medium]} A design review is specifically concerned with this risk: Writing every interaction to long-term memory creates noise and privacy exposure. Which topic should be reviewed first?
\begin{enumerate}[label=\Alph*.]
\item Agent memory
\item Human in the loop
\item Idempotent tools
\item Reflection
\end{enumerate}
\noindent\textbf{Answer.} Agent memory
\par\textit{Working, episodic, semantic, and procedural memories preserve relevant state beyond a single model call. Tradeoff: Memory supports continuity but retrieval, deletion, privacy, and stale facts become product concerns. Watch for: Writing every interaction to long-term memory creates noise and privacy exposure.}
\paragraph{MCQ-0575 [hard]} Which design note most accurately balances the benefits and costs of Agent memory?
\begin{enumerate}[label=\Alph*.]
\item Reliable retries require deduplication state and carefully defined semantics.
\item Memory supports continuity but retrieval, deletion, privacy, and stale facts become product concerns.
\item It can improve recovery without weight updates, while self-critique may repeat the same bias.
\item Oversight reduces irreversible mistakes but adds latency and reviewer load.
\end{enumerate}
\noindent\textbf{Answer.} Memory supports continuity but retrieval, deletion, privacy, and stale facts become product concerns.
\par\textit{Working, episodic, semantic, and procedural memories preserve relevant state beyond a single model call. Tradeoff: Memory supports continuity but retrieval, deletion, privacy, and stale facts become product concerns. Watch for: Writing every interaction to long-term memory creates noise and privacy exposure.}
\paragraph{MCQ-0576 [medium]} Which explanation would be strongest in a system-design interview about Agent memory?
\begin{enumerate}[label=\Alph*.]
\item An idempotent tool can be retried with the same operation key without duplicating the intended effect.
\item Reflection critiques outcomes or trajectories and stores actionable lessons for another attempt.
\item Working, episodic, semantic, and procedural memories preserve relevant state beyond a single model call.
\item High-risk actions pause for review, edit, approval, rejection, or direct human response.
\end{enumerate}
\noindent\textbf{Answer.} Working, episodic, semantic, and procedural memories preserve relevant state beyond a single model call.
\par\textit{Working, episodic, semantic, and procedural memories preserve relevant state beyond a single model call. Tradeoff: Memory supports continuity but retrieval, deletion, privacy, and stale facts become product concerns. Watch for: Writing every interaction to long-term memory creates noise and privacy exposure.}
\paragraph{MCQ-0577 [hard]} Which production warning is most directly associated with Agent memory?
\begin{enumerate}[label=\Alph*.]
\item Writing every interaction to long-term memory creates noise and privacy exposure.
\item Asking for approval after an action has already caused side effects is theater.
\item Assuming HTTP method alone guarantees business idempotency causes duplicate orders or messages.
\item Saving vague reflections adds context without changing behavior.
\end{enumerate}
\noindent\textbf{Answer.} Writing every interaction to long-term memory creates noise and privacy exposure.
\par\textit{Working, episodic, semantic, and procedural memories preserve relevant state beyond a single model call. Tradeoff: Memory supports continuity but retrieval, deletion, privacy, and stale facts become product concerns. Watch for: Writing every interaction to long-term memory creates noise and privacy exposure.}
\paragraph{FC-0578 [medium]} Define Agent memory in interview-ready language.
\noindent\textbf{Answer.} Working, episodic, semantic, and procedural memories preserve relevant state beyond a single model call.
\par\textit{Working, episodic, semantic, and procedural memories preserve relevant state beyond a single model call.}
\paragraph{FC-0579 [hard]} State the main tradeoff and production pitfall for Agent memory.
\noindent\textbf{Answer.} Tradeoff: Memory supports continuity but retrieval, deletion, privacy, and stale facts become product concerns. Pitfall: Writing every interaction to long-term memory creates noise and privacy exposure.
\par\textit{Tradeoff: Memory supports continuity but retrieval, deletion, privacy, and stale facts become product concerns. Pitfall: Writing every interaction to long-term memory creates noise and privacy exposure.}
\paragraph{LA-0580 [hard]} Explain Agent memory from first principles, then show how you would evaluate and operate it in production.
\noindent\textbf{Answer.} Start with the mechanism: Working, episodic, semantic, and procedural memories preserve relevant state beyond a single model call. Make the tradeoff explicit: Memory supports continuity but retrieval, deletion, privacy, and stale facts become product concerns. Define workload-specific offline metrics and a latency/cost budget, test important slices, instrument the critical path, and create a rollback criterion. The production review must explicitly guard against this failure: Writing every interaction to long-term memory creates noise and privacy exposure.
\par\textit{A strong answer connects mechanism, measurable tradeoffs, failure isolation, observability, and rollback rather than listing product features.}
\paragraph{MCQ-0581 [medium]} Which statement best describes Human in the loop?
\begin{enumerate}[label=\Alph*.]
\item Context engineering selects instructions, memory, evidence, tools, schemas, and state for each model turn.
\item MCP standardizes client-server discovery and invocation of tools, resources, prompts, and related capabilities.
\item High-risk actions pause for review, edit, approval, rejection, or direct human response.
\item ReAct interleaves reasoning with actions and observations to update plans using external evidence.
\end{enumerate}
\noindent\textbf{Answer.} High-risk actions pause for review, edit, approval, rejection, or direct human response.
\par\textit{High-risk actions pause for review, edit, approval, rejection, or direct human response. Tradeoff: Oversight reduces irreversible mistakes but adds latency and reviewer load. Watch for: Asking for approval after an action has already caused side effects is theater.}
\paragraph{MCQ-0582 [hard]} What is the central engineering tradeoff of Human in the loop?
\begin{enumerate}[label=\Alph*.]
\item Oversight reduces irreversible mistakes but adds latency and reviewer load.
\item Targeted context improves reliability and cost, but selection logic becomes core application code.
\item Interoperability improves, while trust boundaries, consent, and capability scoping remain host responsibilities.
\item Tool feedback can correct model knowledge, while longer trajectories add latency and error opportunities.
\end{enumerate}
\noindent\textbf{Answer.} Oversight reduces irreversible mistakes but adds latency and reviewer load.
\par\textit{High-risk actions pause for review, edit, approval, rejection, or direct human response. Tradeoff: Oversight reduces irreversible mistakes but adds latency and reviewer load. Watch for: Asking for approval after an action has already caused side effects is theater.}
\paragraph{MCQ-0583 [hard]} Which failure mode should an interviewer expect you to identify for Human in the loop?
\begin{enumerate}[label=\Alph*.]
\item Treating every thought as reliable state allows early mistakes to compound.
\item Treating a discovered MCP server as trusted creates prompt-injection and privilege risks.
\item Appending everything causes attention dilution and instruction conflicts.
\item Asking for approval after an action has already caused side effects is theater.
\end{enumerate}
\noindent\textbf{Answer.} Asking for approval after an action has already caused side effects is theater.
\par\textit{High-risk actions pause for review, edit, approval, rejection, or direct human response. Tradeoff: Oversight reduces irreversible mistakes but adds latency and reviewer load. Watch for: Asking for approval after an action has already caused side effects is theater.}
\paragraph{MCQ-0584 [medium]} A design review is specifically concerned with this risk: Asking for approval after an action has already caused side effects is theater. Which topic should be reviewed first?
\begin{enumerate}[label=\Alph*.]
\item ReAct
\item Context engineering
\item Human in the loop
\item Model Context Protocol
\end{enumerate}
\noindent\textbf{Answer.} Human in the loop
\par\textit{High-risk actions pause for review, edit, approval, rejection, or direct human response. Tradeoff: Oversight reduces irreversible mistakes but adds latency and reviewer load. Watch for: Asking for approval after an action has already caused side effects is theater.}
\paragraph{MCQ-0585 [hard]} Which design note most accurately balances the benefits and costs of Human in the loop?
\begin{enumerate}[label=\Alph*.]
\item Targeted context improves reliability and cost, but selection logic becomes core application code.
\item Interoperability improves, while trust boundaries, consent, and capability scoping remain host responsibilities.
\item Oversight reduces irreversible mistakes but adds latency and reviewer load.
\item Tool feedback can correct model knowledge, while longer trajectories add latency and error opportunities.
\end{enumerate}
\noindent\textbf{Answer.} Oversight reduces irreversible mistakes but adds latency and reviewer load.
\par\textit{High-risk actions pause for review, edit, approval, rejection, or direct human response. Tradeoff: Oversight reduces irreversible mistakes but adds latency and reviewer load. Watch for: Asking for approval after an action has already caused side effects is theater.}
\paragraph{MCQ-0586 [medium]} Which explanation would be strongest in a system-design interview about Human in the loop?
\begin{enumerate}[label=\Alph*.]
\item High-risk actions pause for review, edit, approval, rejection, or direct human response.
\item MCP standardizes client-server discovery and invocation of tools, resources, prompts, and related capabilities.
\item Context engineering selects instructions, memory, evidence, tools, schemas, and state for each model turn.
\item ReAct interleaves reasoning with actions and observations to update plans using external evidence.
\end{enumerate}
\noindent\textbf{Answer.} High-risk actions pause for review, edit, approval, rejection, or direct human response.
\par\textit{High-risk actions pause for review, edit, approval, rejection, or direct human response. Tradeoff: Oversight reduces irreversible mistakes but adds latency and reviewer load. Watch for: Asking for approval after an action has already caused side effects is theater.}
\paragraph{MCQ-0587 [hard]} Which production warning is most directly associated with Human in the loop?
\begin{enumerate}[label=\Alph*.]
\item Appending everything causes attention dilution and instruction conflicts.
\item Treating a discovered MCP server as trusted creates prompt-injection and privilege risks.
\item Asking for approval after an action has already caused side effects is theater.
\item Treating every thought as reliable state allows early mistakes to compound.
\end{enumerate}
\noindent\textbf{Answer.} Asking for approval after an action has already caused side effects is theater.
\par\textit{High-risk actions pause for review, edit, approval, rejection, or direct human response. Tradeoff: Oversight reduces irreversible mistakes but adds latency and reviewer load. Watch for: Asking for approval after an action has already caused side effects is theater.}
\paragraph{FC-0588 [medium]} Define Human in the loop in interview-ready language.
\noindent\textbf{Answer.} High-risk actions pause for review, edit, approval, rejection, or direct human response.
\par\textit{High-risk actions pause for review, edit, approval, rejection, or direct human response.}
\paragraph{FC-0589 [hard]} State the main tradeoff and production pitfall for Human in the loop.
\noindent\textbf{Answer.} Tradeoff: Oversight reduces irreversible mistakes but adds latency and reviewer load. Pitfall: Asking for approval after an action has already caused side effects is theater.
\par\textit{Tradeoff: Oversight reduces irreversible mistakes but adds latency and reviewer load. Pitfall: Asking for approval after an action has already caused side effects is theater.}
\paragraph{LA-0590 [hard]} Explain Human in the loop from first principles, then show how you would evaluate and operate it in production.
\noindent\textbf{Answer.} Start with the mechanism: High-risk actions pause for review, edit, approval, rejection, or direct human response. Make the tradeoff explicit: Oversight reduces irreversible mistakes but adds latency and reviewer load. Define workload-specific offline metrics and a latency/cost budget, test important slices, instrument the critical path, and create a rollback criterion. The production review must explicitly guard against this failure: Asking for approval after an action has already caused side effects is theater.
\par\textit{A strong answer connects mechanism, measurable tradeoffs, failure isolation, observability, and rollback rather than listing product features.}
\paragraph{MCQ-0591 [medium]} Which statement best describes State machines?
\begin{enumerate}[label=\Alph*.]
\item Reflection critiques outcomes or trajectories and stores actionable lessons for another attempt.
\item Explicit states and transitions constrain allowed agent behavior and make recovery inspectable.
\item Models emit structured arguments for declared functions that a host validates and executes.
\item Planning decomposes goals into ordered, conditional, or revisable steps before or during execution.
\end{enumerate}
\noindent\textbf{Answer.} Explicit states and transitions constrain allowed agent behavior and make recovery inspectable.
\par\textit{Explicit states and transitions constrain allowed agent behavior and make recovery inspectable. Tradeoff: Predictability and testability improve at the cost of less open-ended autonomy. Watch for: Encoding side effects inside transition guards breaks replay safety.}
\paragraph{MCQ-0592 [hard]} What is the central engineering tradeoff of State machines?
\begin{enumerate}[label=\Alph*.]
\item It narrows the action interface but does not make proposed actions safe or correct.
\item Upfront plans improve structure but become stale in uncertain environments.
\item It can improve recovery without weight updates, while self-critique may repeat the same bias.
\item Predictability and testability improve at the cost of less open-ended autonomy.
\end{enumerate}
\noindent\textbf{Answer.} Predictability and testability improve at the cost of less open-ended autonomy.
\par\textit{Explicit states and transitions constrain allowed agent behavior and make recovery inspectable. Tradeoff: Predictability and testability improve at the cost of less open-ended autonomy. Watch for: Encoding side effects inside transition guards breaks replay safety.}
\paragraph{MCQ-0593 [hard]} Which failure mode should an interviewer expect you to identify for State machines?
\begin{enumerate}[label=\Alph*.]
\item Encoding side effects inside transition guards breaks replay safety.
\item Forcing a long immutable plan prevents recovery after new observations.
\item Executing model-produced arguments without schema and authorization checks invites damage.
\item Saving vague reflections adds context without changing behavior.
\end{enumerate}
\noindent\textbf{Answer.} Encoding side effects inside transition guards breaks replay safety.
\par\textit{Explicit states and transitions constrain allowed agent behavior and make recovery inspectable. Tradeoff: Predictability and testability improve at the cost of less open-ended autonomy. Watch for: Encoding side effects inside transition guards breaks replay safety.}
\paragraph{MCQ-0594 [medium]} A design review is specifically concerned with this risk: Encoding side effects inside transition guards breaks replay safety. Which topic should be reviewed first?
\begin{enumerate}[label=\Alph*.]
\item State machines
\item Planning
\item Reflection
\item Tool calling
\end{enumerate}
\noindent\textbf{Answer.} State machines
\par\textit{Explicit states and transitions constrain allowed agent behavior and make recovery inspectable. Tradeoff: Predictability and testability improve at the cost of less open-ended autonomy. Watch for: Encoding side effects inside transition guards breaks replay safety.}
\paragraph{MCQ-0595 [hard]} Which design note most accurately balances the benefits and costs of State machines?
\begin{enumerate}[label=\Alph*.]
\item It narrows the action interface but does not make proposed actions safe or correct.
\item Predictability and testability improve at the cost of less open-ended autonomy.
\item Upfront plans improve structure but become stale in uncertain environments.
\item It can improve recovery without weight updates, while self-critique may repeat the same bias.
\end{enumerate}
\noindent\textbf{Answer.} Predictability and testability improve at the cost of less open-ended autonomy.
\par\textit{Explicit states and transitions constrain allowed agent behavior and make recovery inspectable. Tradeoff: Predictability and testability improve at the cost of less open-ended autonomy. Watch for: Encoding side effects inside transition guards breaks replay safety.}
\paragraph{MCQ-0596 [medium]} Which explanation would be strongest in a system-design interview about State machines?
\begin{enumerate}[label=\Alph*.]
\item Explicit states and transitions constrain allowed agent behavior and make recovery inspectable.
\item Models emit structured arguments for declared functions that a host validates and executes.
\item Reflection critiques outcomes or trajectories and stores actionable lessons for another attempt.
\item Planning decomposes goals into ordered, conditional, or revisable steps before or during execution.
\end{enumerate}
\noindent\textbf{Answer.} Explicit states and transitions constrain allowed agent behavior and make recovery inspectable.
\par\textit{Explicit states and transitions constrain allowed agent behavior and make recovery inspectable. Tradeoff: Predictability and testability improve at the cost of less open-ended autonomy. Watch for: Encoding side effects inside transition guards breaks replay safety.}
\paragraph{MCQ-0597 [hard]} Which production warning is most directly associated with State machines?
\begin{enumerate}[label=\Alph*.]
\item Saving vague reflections adds context without changing behavior.
\item Encoding side effects inside transition guards breaks replay safety.
\item Forcing a long immutable plan prevents recovery after new observations.
\item Executing model-produced arguments without schema and authorization checks invites damage.
\end{enumerate}
\noindent\textbf{Answer.} Encoding side effects inside transition guards breaks replay safety.
\par\textit{Explicit states and transitions constrain allowed agent behavior and make recovery inspectable. Tradeoff: Predictability and testability improve at the cost of less open-ended autonomy. Watch for: Encoding side effects inside transition guards breaks replay safety.}
\paragraph{FC-0598 [medium]} Define State machines in interview-ready language.
\noindent\textbf{Answer.} Explicit states and transitions constrain allowed agent behavior and make recovery inspectable.
\par\textit{Explicit states and transitions constrain allowed agent behavior and make recovery inspectable.}
\paragraph{FC-0599 [hard]} State the main tradeoff and production pitfall for State machines.
\noindent\textbf{Answer.} Tradeoff: Predictability and testability improve at the cost of less open-ended autonomy. Pitfall: Encoding side effects inside transition guards breaks replay safety.
\par\textit{Tradeoff: Predictability and testability improve at the cost of less open-ended autonomy. Pitfall: Encoding side effects inside transition guards breaks replay safety.}
\paragraph{LA-0600 [hard]} Explain State machines from first principles, then show how you would evaluate and operate it in production.
\noindent\textbf{Answer.} Start with the mechanism: Explicit states and transitions constrain allowed agent behavior and make recovery inspectable. Make the tradeoff explicit: Predictability and testability improve at the cost of less open-ended autonomy. Define workload-specific offline metrics and a latency/cost budget, test important slices, instrument the critical path, and create a rollback criterion. The production review must explicitly guard against this failure: Encoding side effects inside transition guards breaks replay safety.
\par\textit{A strong answer connects mechanism, measurable tradeoffs, failure isolation, observability, and rollback rather than listing product features.}
\paragraph{MCQ-0601 [medium]} Which statement best describes Idempotent tools?
\begin{enumerate}[label=\Alph*.]
\item An idempotent tool can be retried with the same operation key without duplicating the intended effect.
\item ReAct interleaves reasoning with actions and observations to update plans using external evidence.
\item MCP standardizes client-server discovery and invocation of tools, resources, prompts, and related capabilities.
\item Explicit states and transitions constrain allowed agent behavior and make recovery inspectable.
\end{enumerate}
\noindent\textbf{Answer.} An idempotent tool can be retried with the same operation key without duplicating the intended effect.
\par\textit{An idempotent tool can be retried with the same operation key without duplicating the intended effect. Tradeoff: Reliable retries require deduplication state and carefully defined semantics. Watch for: Assuming HTTP method alone guarantees business idempotency causes duplicate orders or messages.}
\paragraph{MCQ-0602 [hard]} What is the central engineering tradeoff of Idempotent tools?
\begin{enumerate}[label=\Alph*.]
\item Predictability and testability improve at the cost of less open-ended autonomy.
\item Interoperability improves, while trust boundaries, consent, and capability scoping remain host responsibilities.
\item Tool feedback can correct model knowledge, while longer trajectories add latency and error opportunities.
\item Reliable retries require deduplication state and carefully defined semantics.
\end{enumerate}
\noindent\textbf{Answer.} Reliable retries require deduplication state and carefully defined semantics.
\par\textit{An idempotent tool can be retried with the same operation key without duplicating the intended effect. Tradeoff: Reliable retries require deduplication state and carefully defined semantics. Watch for: Assuming HTTP method alone guarantees business idempotency causes duplicate orders or messages.}
\paragraph{MCQ-0603 [hard]} Which failure mode should an interviewer expect you to identify for Idempotent tools?
\begin{enumerate}[label=\Alph*.]
\item Treating a discovered MCP server as trusted creates prompt-injection and privilege risks.
\item Treating every thought as reliable state allows early mistakes to compound.
\item Assuming HTTP method alone guarantees business idempotency causes duplicate orders or messages.
\item Encoding side effects inside transition guards breaks replay safety.
\end{enumerate}
\noindent\textbf{Answer.} Assuming HTTP method alone guarantees business idempotency causes duplicate orders or messages.
\par\textit{An idempotent tool can be retried with the same operation key without duplicating the intended effect. Tradeoff: Reliable retries require deduplication state and carefully defined semantics. Watch for: Assuming HTTP method alone guarantees business idempotency causes duplicate orders or messages.}
\paragraph{MCQ-0604 [medium]} A design review is specifically concerned with this risk: Assuming HTTP method alone guarantees business idempotency causes duplicate orders or messages. Which topic should be reviewed first?
\begin{enumerate}[label=\Alph*.]
\item State machines
\item Idempotent tools
\item Model Context Protocol
\item ReAct
\end{enumerate}
\noindent\textbf{Answer.} Idempotent tools
\par\textit{An idempotent tool can be retried with the same operation key without duplicating the intended effect. Tradeoff: Reliable retries require deduplication state and carefully defined semantics. Watch for: Assuming HTTP method alone guarantees business idempotency causes duplicate orders or messages.}
\paragraph{MCQ-0605 [hard]} Which design note most accurately balances the benefits and costs of Idempotent tools?
\begin{enumerate}[label=\Alph*.]
\item Tool feedback can correct model knowledge, while longer trajectories add latency and error opportunities.
\item Reliable retries require deduplication state and carefully defined semantics.
\item Predictability and testability improve at the cost of less open-ended autonomy.
\item Interoperability improves, while trust boundaries, consent, and capability scoping remain host responsibilities.
\end{enumerate}
\noindent\textbf{Answer.} Reliable retries require deduplication state and carefully defined semantics.
\par\textit{An idempotent tool can be retried with the same operation key without duplicating the intended effect. Tradeoff: Reliable retries require deduplication state and carefully defined semantics. Watch for: Assuming HTTP method alone guarantees business idempotency causes duplicate orders or messages.}
\paragraph{MCQ-0606 [medium]} Which explanation would be strongest in a system-design interview about Idempotent tools?
\begin{enumerate}[label=\Alph*.]
\item An idempotent tool can be retried with the same operation key without duplicating the intended effect.
\item ReAct interleaves reasoning with actions and observations to update plans using external evidence.
\item MCP standardizes client-server discovery and invocation of tools, resources, prompts, and related capabilities.
\item Explicit states and transitions constrain allowed agent behavior and make recovery inspectable.
\end{enumerate}
\noindent\textbf{Answer.} An idempotent tool can be retried with the same operation key without duplicating the intended effect.
\par\textit{An idempotent tool can be retried with the same operation key without duplicating the intended effect. Tradeoff: Reliable retries require deduplication state and carefully defined semantics. Watch for: Assuming HTTP method alone guarantees business idempotency causes duplicate orders or messages.}
\paragraph{MCQ-0607 [hard]} Which production warning is most directly associated with Idempotent tools?
\begin{enumerate}[label=\Alph*.]
\item Encoding side effects inside transition guards breaks replay safety.
\item Treating a discovered MCP server as trusted creates prompt-injection and privilege risks.
\item Assuming HTTP method alone guarantees business idempotency causes duplicate orders or messages.
\item Treating every thought as reliable state allows early mistakes to compound.
\end{enumerate}
\noindent\textbf{Answer.} Assuming HTTP method alone guarantees business idempotency causes duplicate orders or messages.
\par\textit{An idempotent tool can be retried with the same operation key without duplicating the intended effect. Tradeoff: Reliable retries require deduplication state and carefully defined semantics. Watch for: Assuming HTTP method alone guarantees business idempotency causes duplicate orders or messages.}
\paragraph{FC-0608 [medium]} Define Idempotent tools in interview-ready language.
\noindent\textbf{Answer.} An idempotent tool can be retried with the same operation key without duplicating the intended effect.
\par\textit{An idempotent tool can be retried with the same operation key without duplicating the intended effect.}
\paragraph{FC-0609 [hard]} State the main tradeoff and production pitfall for Idempotent tools.
\noindent\textbf{Answer.} Tradeoff: Reliable retries require deduplication state and carefully defined semantics. Pitfall: Assuming HTTP method alone guarantees business idempotency causes duplicate orders or messages.
\par\textit{Tradeoff: Reliable retries require deduplication state and carefully defined semantics. Pitfall: Assuming HTTP method alone guarantees business idempotency causes duplicate orders or messages.}
\paragraph{LA-0610 [hard]} Explain Idempotent tools from first principles, then show how you would evaluate and operate it in production.
\noindent\textbf{Answer.} Start with the mechanism: An idempotent tool can be retried with the same operation key without duplicating the intended effect. Make the tradeoff explicit: Reliable retries require deduplication state and carefully defined semantics. Define workload-specific offline metrics and a latency/cost budget, test important slices, instrument the critical path, and create a rollback criterion. The production review must explicitly guard against this failure: Assuming HTTP method alone guarantees business idempotency causes duplicate orders or messages.
\par\textit{A strong answer connects mechanism, measurable tradeoffs, failure isolation, observability, and rollback rather than listing product features.}
\paragraph{MCQ-0611 [medium]} Which statement best describes Agent handoffs?
\begin{enumerate}[label=\Alph*.]
\item A sandbox isolates untrusted code with resource limits, scoped files, network policy, timeouts, and audit logs.
\item MCP standardizes client-server discovery and invocation of tools, resources, prompts, and related capabilities.
\item A handoff transfers a task and structured context to a specialized agent or deterministic service.
\item ReAct interleaves reasoning with actions and observations to update plans using external evidence.
\end{enumerate}
\noindent\textbf{Answer.} A handoff transfers a task and structured context to a specialized agent or deterministic service.
\par\textit{A handoff transfers a task and structured context to a specialized agent or deterministic service. Tradeoff: Specialization reduces prompt overload but adds routing and ownership ambiguity. Watch for: Passing the entire transcript leaks irrelevant data and confuses the recipient.}
\paragraph{MCQ-0612 [hard]} What is the central engineering tradeoff of Agent handoffs?
\begin{enumerate}[label=\Alph*.]
\item Interoperability improves, while trust boundaries, consent, and capability scoping remain host responsibilities.
\item Specialization reduces prompt overload but adds routing and ownership ambiguity.
\item Tool feedback can correct model knowledge, while longer trajectories add latency and error opportunities.
\item Isolation enables capability while imposing startup and compatibility overhead.
\end{enumerate}
\noindent\textbf{Answer.} Specialization reduces prompt overload but adds routing and ownership ambiguity.
\par\textit{A handoff transfers a task and structured context to a specialized agent or deterministic service. Tradeoff: Specialization reduces prompt overload but adds routing and ownership ambiguity. Watch for: Passing the entire transcript leaks irrelevant data and confuses the recipient.}
\paragraph{MCQ-0613 [hard]} Which failure mode should an interviewer expect you to identify for Agent handoffs?
\begin{enumerate}[label=\Alph*.]
\item Passing the entire transcript leaks irrelevant data and confuses the recipient.
\item A container without kernel, credential, and network boundaries is not a sufficient sandbox.
\item Treating every thought as reliable state allows early mistakes to compound.
\item Treating a discovered MCP server as trusted creates prompt-injection and privilege risks.
\end{enumerate}
\noindent\textbf{Answer.} Passing the entire transcript leaks irrelevant data and confuses the recipient.
\par\textit{A handoff transfers a task and structured context to a specialized agent or deterministic service. Tradeoff: Specialization reduces prompt overload but adds routing and ownership ambiguity. Watch for: Passing the entire transcript leaks irrelevant data and confuses the recipient.}
\paragraph{MCQ-0614 [medium]} A design review is specifically concerned with this risk: Passing the entire transcript leaks irrelevant data and confuses the recipient. Which topic should be reviewed first?
\begin{enumerate}[label=\Alph*.]
\item Agent handoffs
\item Model Context Protocol
\item ReAct
\item Sandbox runs
\end{enumerate}
\noindent\textbf{Answer.} Agent handoffs
\par\textit{A handoff transfers a task and structured context to a specialized agent or deterministic service. Tradeoff: Specialization reduces prompt overload but adds routing and ownership ambiguity. Watch for: Passing the entire transcript leaks irrelevant data and confuses the recipient.}
\paragraph{MCQ-0615 [hard]} Which design note most accurately balances the benefits and costs of Agent handoffs?
\begin{enumerate}[label=\Alph*.]
\item Tool feedback can correct model knowledge, while longer trajectories add latency and error opportunities.
\item Isolation enables capability while imposing startup and compatibility overhead.
\item Interoperability improves, while trust boundaries, consent, and capability scoping remain host responsibilities.
\item Specialization reduces prompt overload but adds routing and ownership ambiguity.
\end{enumerate}
\noindent\textbf{Answer.} Specialization reduces prompt overload but adds routing and ownership ambiguity.
\par\textit{A handoff transfers a task and structured context to a specialized agent or deterministic service. Tradeoff: Specialization reduces prompt overload but adds routing and ownership ambiguity. Watch for: Passing the entire transcript leaks irrelevant data and confuses the recipient.}
\paragraph{MCQ-0616 [medium]} Which explanation would be strongest in a system-design interview about Agent handoffs?
\begin{enumerate}[label=\Alph*.]
\item MCP standardizes client-server discovery and invocation of tools, resources, prompts, and related capabilities.
\item ReAct interleaves reasoning with actions and observations to update plans using external evidence.
\item A sandbox isolates untrusted code with resource limits, scoped files, network policy, timeouts, and audit logs.
\item A handoff transfers a task and structured context to a specialized agent or deterministic service.
\end{enumerate}
\noindent\textbf{Answer.} A handoff transfers a task and structured context to a specialized agent or deterministic service.
\par\textit{A handoff transfers a task and structured context to a specialized agent or deterministic service. Tradeoff: Specialization reduces prompt overload but adds routing and ownership ambiguity. Watch for: Passing the entire transcript leaks irrelevant data and confuses the recipient.}
\paragraph{MCQ-0617 [hard]} Which production warning is most directly associated with Agent handoffs?
\begin{enumerate}[label=\Alph*.]
\item Passing the entire transcript leaks irrelevant data and confuses the recipient.
\item Treating a discovered MCP server as trusted creates prompt-injection and privilege risks.
\item A container without kernel, credential, and network boundaries is not a sufficient sandbox.
\item Treating every thought as reliable state allows early mistakes to compound.
\end{enumerate}
\noindent\textbf{Answer.} Passing the entire transcript leaks irrelevant data and confuses the recipient.
\par\textit{A handoff transfers a task and structured context to a specialized agent or deterministic service. Tradeoff: Specialization reduces prompt overload but adds routing and ownership ambiguity. Watch for: Passing the entire transcript leaks irrelevant data and confuses the recipient.}
\paragraph{FC-0618 [medium]} Define Agent handoffs in interview-ready language.
\noindent\textbf{Answer.} A handoff transfers a task and structured context to a specialized agent or deterministic service.
\par\textit{A handoff transfers a task and structured context to a specialized agent or deterministic service.}
\paragraph{FC-0619 [hard]} State the main tradeoff and production pitfall for Agent handoffs.
\noindent\textbf{Answer.} Tradeoff: Specialization reduces prompt overload but adds routing and ownership ambiguity. Pitfall: Passing the entire transcript leaks irrelevant data and confuses the recipient.
\par\textit{Tradeoff: Specialization reduces prompt overload but adds routing and ownership ambiguity. Pitfall: Passing the entire transcript leaks irrelevant data and confuses the recipient.}
\paragraph{LA-0620 [hard]} Explain Agent handoffs from first principles, then show how you would evaluate and operate it in production.
\noindent\textbf{Answer.} Start with the mechanism: A handoff transfers a task and structured context to a specialized agent or deterministic service. Make the tradeoff explicit: Specialization reduces prompt overload but adds routing and ownership ambiguity. Define workload-specific offline metrics and a latency/cost budget, test important slices, instrument the critical path, and create a rollback criterion. The production review must explicitly guard against this failure: Passing the entire transcript leaks irrelevant data and confuses the recipient.
\par\textit{A strong answer connects mechanism, measurable tradeoffs, failure isolation, observability, and rollback rather than listing product features.}
\paragraph{MCQ-0621 [medium]} Which statement best describes Autonomy levels?
\begin{enumerate}[label=\Alph*.]
\item A handoff transfers a task and structured context to a specialized agent or deterministic service.
\item Autonomy is graduated from suggest-only through approval-gated execution to bounded autonomous operation.
\item An idempotent tool can be retried with the same operation key without duplicating the intended effect.
\item Context engineering selects instructions, memory, evidence, tools, schemas, and state for each model turn.
\end{enumerate}
\noindent\textbf{Answer.} Autonomy is graduated from suggest-only through approval-gated execution to bounded autonomous operation.
\par\textit{Autonomy is graduated from suggest-only through approval-gated execution to bounded autonomous operation. Tradeoff: More autonomy reduces human effort while increasing required assurance and containment. Watch for: Assigning one global autonomy level ignores action-specific risk.}
\paragraph{MCQ-0622 [hard]} What is the central engineering tradeoff of Autonomy levels?
\begin{enumerate}[label=\Alph*.]
\item More autonomy reduces human effort while increasing required assurance and containment.
\item Specialization reduces prompt overload but adds routing and ownership ambiguity.
\item Reliable retries require deduplication state and carefully defined semantics.
\item Targeted context improves reliability and cost, but selection logic becomes core application code.
\end{enumerate}
\noindent\textbf{Answer.} More autonomy reduces human effort while increasing required assurance and containment.
\par\textit{Autonomy is graduated from suggest-only through approval-gated execution to bounded autonomous operation. Tradeoff: More autonomy reduces human effort while increasing required assurance and containment. Watch for: Assigning one global autonomy level ignores action-specific risk.}
\paragraph{MCQ-0623 [hard]} Which failure mode should an interviewer expect you to identify for Autonomy levels?
\begin{enumerate}[label=\Alph*.]
\item Appending everything causes attention dilution and instruction conflicts.
\item Passing the entire transcript leaks irrelevant data and confuses the recipient.
\item Assigning one global autonomy level ignores action-specific risk.
\item Assuming HTTP method alone guarantees business idempotency causes duplicate orders or messages.
\end{enumerate}
\noindent\textbf{Answer.} Assigning one global autonomy level ignores action-specific risk.
\par\textit{Autonomy is graduated from suggest-only through approval-gated execution to bounded autonomous operation. Tradeoff: More autonomy reduces human effort while increasing required assurance and containment. Watch for: Assigning one global autonomy level ignores action-specific risk.}
\paragraph{MCQ-0624 [medium]} A design review is specifically concerned with this risk: Assigning one global autonomy level ignores action-specific risk. Which topic should be reviewed first?
\begin{enumerate}[label=\Alph*.]
\item Autonomy levels
\item Context engineering
\item Idempotent tools
\item Agent handoffs
\end{enumerate}
\noindent\textbf{Answer.} Autonomy levels
\par\textit{Autonomy is graduated from suggest-only through approval-gated execution to bounded autonomous operation. Tradeoff: More autonomy reduces human effort while increasing required assurance and containment. Watch for: Assigning one global autonomy level ignores action-specific risk.}
\paragraph{MCQ-0625 [hard]} Which design note most accurately balances the benefits and costs of Autonomy levels?
\begin{enumerate}[label=\Alph*.]
\item Reliable retries require deduplication state and carefully defined semantics.
\item Targeted context improves reliability and cost, but selection logic becomes core application code.
\item More autonomy reduces human effort while increasing required assurance and containment.
\item Specialization reduces prompt overload but adds routing and ownership ambiguity.
\end{enumerate}
\noindent\textbf{Answer.} More autonomy reduces human effort while increasing required assurance and containment.
\par\textit{Autonomy is graduated from suggest-only through approval-gated execution to bounded autonomous operation. Tradeoff: More autonomy reduces human effort while increasing required assurance and containment. Watch for: Assigning one global autonomy level ignores action-specific risk.}
\paragraph{MCQ-0626 [medium]} Which explanation would be strongest in a system-design interview about Autonomy levels?
\begin{enumerate}[label=\Alph*.]
\item Autonomy is graduated from suggest-only through approval-gated execution to bounded autonomous operation.
\item An idempotent tool can be retried with the same operation key without duplicating the intended effect.
\item A handoff transfers a task and structured context to a specialized agent or deterministic service.
\item Context engineering selects instructions, memory, evidence, tools, schemas, and state for each model turn.
\end{enumerate}
\noindent\textbf{Answer.} Autonomy is graduated from suggest-only through approval-gated execution to bounded autonomous operation.
\par\textit{Autonomy is graduated from suggest-only through approval-gated execution to bounded autonomous operation. Tradeoff: More autonomy reduces human effort while increasing required assurance and containment. Watch for: Assigning one global autonomy level ignores action-specific risk.}
\paragraph{MCQ-0627 [hard]} Which production warning is most directly associated with Autonomy levels?
\begin{enumerate}[label=\Alph*.]
\item Appending everything causes attention dilution and instruction conflicts.
\item Passing the entire transcript leaks irrelevant data and confuses the recipient.
\item Assuming HTTP method alone guarantees business idempotency causes duplicate orders or messages.
\item Assigning one global autonomy level ignores action-specific risk.
\end{enumerate}
\noindent\textbf{Answer.} Assigning one global autonomy level ignores action-specific risk.
\par\textit{Autonomy is graduated from suggest-only through approval-gated execution to bounded autonomous operation. Tradeoff: More autonomy reduces human effort while increasing required assurance and containment. Watch for: Assigning one global autonomy level ignores action-specific risk.}
\paragraph{FC-0628 [medium]} Define Autonomy levels in interview-ready language.
\noindent\textbf{Answer.} Autonomy is graduated from suggest-only through approval-gated execution to bounded autonomous operation.
\par\textit{Autonomy is graduated from suggest-only through approval-gated execution to bounded autonomous operation.}
\paragraph{FC-0629 [hard]} State the main tradeoff and production pitfall for Autonomy levels.
\noindent\textbf{Answer.} Tradeoff: More autonomy reduces human effort while increasing required assurance and containment. Pitfall: Assigning one global autonomy level ignores action-specific risk.
\par\textit{Tradeoff: More autonomy reduces human effort while increasing required assurance and containment. Pitfall: Assigning one global autonomy level ignores action-specific risk.}
\paragraph{LA-0630 [hard]} Explain Autonomy levels from first principles, then show how you would evaluate and operate it in production.
\noindent\textbf{Answer.} Start with the mechanism: Autonomy is graduated from suggest-only through approval-gated execution to bounded autonomous operation. Make the tradeoff explicit: More autonomy reduces human effort while increasing required assurance and containment. Define workload-specific offline metrics and a latency/cost budget, test important slices, instrument the critical path, and create a rollback criterion. The production review must explicitly guard against this failure: Assigning one global autonomy level ignores action-specific risk.
\par\textit{A strong answer connects mechanism, measurable tradeoffs, failure isolation, observability, and rollback rather than listing product features.}
\paragraph{MCQ-0631 [medium]} Which statement best describes Context engineering?
\begin{enumerate}[label=\Alph*.]
\item Planning decomposes goals into ordered, conditional, or revisable steps before or during execution.
\item Reflection critiques outcomes or trajectories and stores actionable lessons for another attempt.
\item Context engineering selects instructions, memory, evidence, tools, schemas, and state for each model turn.
\item MCP standardizes client-server discovery and invocation of tools, resources, prompts, and related capabilities.
\end{enumerate}
\noindent\textbf{Answer.} Context engineering selects instructions, memory, evidence, tools, schemas, and state for each model turn.
\par\textit{Context engineering selects instructions, memory, evidence, tools, schemas, and state for each model turn. Tradeoff: Targeted context improves reliability and cost, but selection logic becomes core application code. Watch for: Appending everything causes attention dilution and instruction conflicts.}
\paragraph{MCQ-0632 [hard]} What is the central engineering tradeoff of Context engineering?
\begin{enumerate}[label=\Alph*.]
\item Interoperability improves, while trust boundaries, consent, and capability scoping remain host responsibilities.
\item It can improve recovery without weight updates, while self-critique may repeat the same bias.
\item Targeted context improves reliability and cost, but selection logic becomes core application code.
\item Upfront plans improve structure but become stale in uncertain environments.
\end{enumerate}
\noindent\textbf{Answer.} Targeted context improves reliability and cost, but selection logic becomes core application code.
\par\textit{Context engineering selects instructions, memory, evidence, tools, schemas, and state for each model turn. Tradeoff: Targeted context improves reliability and cost, but selection logic becomes core application code. Watch for: Appending everything causes attention dilution and instruction conflicts.}
\paragraph{MCQ-0633 [hard]} Which failure mode should an interviewer expect you to identify for Context engineering?
\begin{enumerate}[label=\Alph*.]
\item Forcing a long immutable plan prevents recovery after new observations.
\item Treating a discovered MCP server as trusted creates prompt-injection and privilege risks.
\item Appending everything causes attention dilution and instruction conflicts.
\item Saving vague reflections adds context without changing behavior.
\end{enumerate}
\noindent\textbf{Answer.} Appending everything causes attention dilution and instruction conflicts.
\par\textit{Context engineering selects instructions, memory, evidence, tools, schemas, and state for each model turn. Tradeoff: Targeted context improves reliability and cost, but selection logic becomes core application code. Watch for: Appending everything causes attention dilution and instruction conflicts.}
\paragraph{MCQ-0634 [medium]} A design review is specifically concerned with this risk: Appending everything causes attention dilution and instruction conflicts. Which topic should be reviewed first?
\begin{enumerate}[label=\Alph*.]
\item Planning
\item Context engineering
\item Model Context Protocol
\item Reflection
\end{enumerate}
\noindent\textbf{Answer.} Context engineering
\par\textit{Context engineering selects instructions, memory, evidence, tools, schemas, and state for each model turn. Tradeoff: Targeted context improves reliability and cost, but selection logic becomes core application code. Watch for: Appending everything causes attention dilution and instruction conflicts.}
\paragraph{MCQ-0635 [hard]} Which design note most accurately balances the benefits and costs of Context engineering?
\begin{enumerate}[label=\Alph*.]
\item Interoperability improves, while trust boundaries, consent, and capability scoping remain host responsibilities.
\item It can improve recovery without weight updates, while self-critique may repeat the same bias.
\item Targeted context improves reliability and cost, but selection logic becomes core application code.
\item Upfront plans improve structure but become stale in uncertain environments.
\end{enumerate}
\noindent\textbf{Answer.} Targeted context improves reliability and cost, but selection logic becomes core application code.
\par\textit{Context engineering selects instructions, memory, evidence, tools, schemas, and state for each model turn. Tradeoff: Targeted context improves reliability and cost, but selection logic becomes core application code. Watch for: Appending everything causes attention dilution and instruction conflicts.}
\paragraph{MCQ-0636 [medium]} Which explanation would be strongest in a system-design interview about Context engineering?
\begin{enumerate}[label=\Alph*.]
\item Planning decomposes goals into ordered, conditional, or revisable steps before or during execution.
\item Context engineering selects instructions, memory, evidence, tools, schemas, and state for each model turn.
\item Reflection critiques outcomes or trajectories and stores actionable lessons for another attempt.
\item MCP standardizes client-server discovery and invocation of tools, resources, prompts, and related capabilities.
\end{enumerate}
\noindent\textbf{Answer.} Context engineering selects instructions, memory, evidence, tools, schemas, and state for each model turn.
\par\textit{Context engineering selects instructions, memory, evidence, tools, schemas, and state for each model turn. Tradeoff: Targeted context improves reliability and cost, but selection logic becomes core application code. Watch for: Appending everything causes attention dilution and instruction conflicts.}
\paragraph{MCQ-0637 [hard]} Which production warning is most directly associated with Context engineering?
\begin{enumerate}[label=\Alph*.]
\item Saving vague reflections adds context without changing behavior.
\item Treating a discovered MCP server as trusted creates prompt-injection and privilege risks.
\item Appending everything causes attention dilution and instruction conflicts.
\item Forcing a long immutable plan prevents recovery after new observations.
\end{enumerate}
\noindent\textbf{Answer.} Appending everything causes attention dilution and instruction conflicts.
\par\textit{Context engineering selects instructions, memory, evidence, tools, schemas, and state for each model turn. Tradeoff: Targeted context improves reliability and cost, but selection logic becomes core application code. Watch for: Appending everything causes attention dilution and instruction conflicts.}
\paragraph{FC-0638 [medium]} Define Context engineering in interview-ready language.
\noindent\textbf{Answer.} Context engineering selects instructions, memory, evidence, tools, schemas, and state for each model turn.
\par\textit{Context engineering selects instructions, memory, evidence, tools, schemas, and state for each model turn.}
\paragraph{FC-0639 [hard]} State the main tradeoff and production pitfall for Context engineering.
\noindent\textbf{Answer.} Tradeoff: Targeted context improves reliability and cost, but selection logic becomes core application code. Pitfall: Appending everything causes attention dilution and instruction conflicts.
\par\textit{Tradeoff: Targeted context improves reliability and cost, but selection logic becomes core application code. Pitfall: Appending everything causes attention dilution and instruction conflicts.}
\paragraph{LA-0640 [hard]} Explain Context engineering from first principles, then show how you would evaluate and operate it in production.
\noindent\textbf{Answer.} Start with the mechanism: Context engineering selects instructions, memory, evidence, tools, schemas, and state for each model turn. Make the tradeoff explicit: Targeted context improves reliability and cost, but selection logic becomes core application code. Define workload-specific offline metrics and a latency/cost budget, test important slices, instrument the critical path, and create a rollback criterion. The production review must explicitly guard against this failure: Appending everything causes attention dilution and instruction conflicts.
\par\textit{A strong answer connects mechanism, measurable tradeoffs, failure isolation, observability, and rollback rather than listing product features.}
\paragraph{MCQ-0641 [medium]} Which statement best describes Sandbox runs?
\begin{enumerate}[label=\Alph*.]
\item A sandbox isolates untrusted code with resource limits, scoped files, network policy, timeouts, and audit logs.
\item MCP standardizes client-server discovery and invocation of tools, resources, prompts, and related capabilities.
\item A handoff transfers a task and structured context to a specialized agent or deterministic service.
\item Reflection critiques outcomes or trajectories and stores actionable lessons for another attempt.
\end{enumerate}
\noindent\textbf{Answer.} A sandbox isolates untrusted code with resource limits, scoped files, network policy, timeouts, and audit logs.
\par\textit{A sandbox isolates untrusted code with resource limits, scoped files, network policy, timeouts, and audit logs. Tradeoff: Isolation enables capability while imposing startup and compatibility overhead. Watch for: A container without kernel, credential, and network boundaries is not a sufficient sandbox.}
\paragraph{MCQ-0642 [hard]} What is the central engineering tradeoff of Sandbox runs?
\begin{enumerate}[label=\Alph*.]
\item Interoperability improves, while trust boundaries, consent, and capability scoping remain host responsibilities.
\item Specialization reduces prompt overload but adds routing and ownership ambiguity.
\item Isolation enables capability while imposing startup and compatibility overhead.
\item It can improve recovery without weight updates, while self-critique may repeat the same bias.
\end{enumerate}
\noindent\textbf{Answer.} Isolation enables capability while imposing startup and compatibility overhead.
\par\textit{A sandbox isolates untrusted code with resource limits, scoped files, network policy, timeouts, and audit logs. Tradeoff: Isolation enables capability while imposing startup and compatibility overhead. Watch for: A container without kernel, credential, and network boundaries is not a sufficient sandbox.}
\paragraph{MCQ-0643 [hard]} Which failure mode should an interviewer expect you to identify for Sandbox runs?
\begin{enumerate}[label=\Alph*.]
\item Saving vague reflections adds context without changing behavior.
\item A container without kernel, credential, and network boundaries is not a sufficient sandbox.
\item Treating a discovered MCP server as trusted creates prompt-injection and privilege risks.
\item Passing the entire transcript leaks irrelevant data and confuses the recipient.
\end{enumerate}
\noindent\textbf{Answer.} A container without kernel, credential, and network boundaries is not a sufficient sandbox.
\par\textit{A sandbox isolates untrusted code with resource limits, scoped files, network policy, timeouts, and audit logs. Tradeoff: Isolation enables capability while imposing startup and compatibility overhead. Watch for: A container without kernel, credential, and network boundaries is not a sufficient sandbox.}
\paragraph{MCQ-0644 [medium]} A design review is specifically concerned with this risk: A container without kernel, credential, and network boundaries is not a sufficient sandbox. Which topic should be reviewed first?
\begin{enumerate}[label=\Alph*.]
\item Sandbox runs
\item Model Context Protocol
\item Reflection
\item Agent handoffs
\end{enumerate}
\noindent\textbf{Answer.} Sandbox runs
\par\textit{A sandbox isolates untrusted code with resource limits, scoped files, network policy, timeouts, and audit logs. Tradeoff: Isolation enables capability while imposing startup and compatibility overhead. Watch for: A container without kernel, credential, and network boundaries is not a sufficient sandbox.}
\paragraph{MCQ-0645 [hard]} Which design note most accurately balances the benefits and costs of Sandbox runs?
\begin{enumerate}[label=\Alph*.]
\item Isolation enables capability while imposing startup and compatibility overhead.
\item Specialization reduces prompt overload but adds routing and ownership ambiguity.
\item It can improve recovery without weight updates, while self-critique may repeat the same bias.
\item Interoperability improves, while trust boundaries, consent, and capability scoping remain host responsibilities.
\end{enumerate}
\noindent\textbf{Answer.} Isolation enables capability while imposing startup and compatibility overhead.
\par\textit{A sandbox isolates untrusted code with resource limits, scoped files, network policy, timeouts, and audit logs. Tradeoff: Isolation enables capability while imposing startup and compatibility overhead. Watch for: A container without kernel, credential, and network boundaries is not a sufficient sandbox.}
\paragraph{MCQ-0646 [medium]} Which explanation would be strongest in a system-design interview about Sandbox runs?
\begin{enumerate}[label=\Alph*.]
\item Reflection critiques outcomes or trajectories and stores actionable lessons for another attempt.
\item A handoff transfers a task and structured context to a specialized agent or deterministic service.
\item A sandbox isolates untrusted code with resource limits, scoped files, network policy, timeouts, and audit logs.
\item MCP standardizes client-server discovery and invocation of tools, resources, prompts, and related capabilities.
\end{enumerate}
\noindent\textbf{Answer.} A sandbox isolates untrusted code with resource limits, scoped files, network policy, timeouts, and audit logs.
\par\textit{A sandbox isolates untrusted code with resource limits, scoped files, network policy, timeouts, and audit logs. Tradeoff: Isolation enables capability while imposing startup and compatibility overhead. Watch for: A container without kernel, credential, and network boundaries is not a sufficient sandbox.}
\paragraph{MCQ-0647 [hard]} Which production warning is most directly associated with Sandbox runs?
\begin{enumerate}[label=\Alph*.]
\item Passing the entire transcript leaks irrelevant data and confuses the recipient.
\item A container without kernel, credential, and network boundaries is not a sufficient sandbox.
\item Treating a discovered MCP server as trusted creates prompt-injection and privilege risks.
\item Saving vague reflections adds context without changing behavior.
\end{enumerate}
\noindent\textbf{Answer.} A container without kernel, credential, and network boundaries is not a sufficient sandbox.
\par\textit{A sandbox isolates untrusted code with resource limits, scoped files, network policy, timeouts, and audit logs. Tradeoff: Isolation enables capability while imposing startup and compatibility overhead. Watch for: A container without kernel, credential, and network boundaries is not a sufficient sandbox.}
\paragraph{FC-0648 [medium]} Define Sandbox runs in interview-ready language.
\noindent\textbf{Answer.} A sandbox isolates untrusted code with resource limits, scoped files, network policy, timeouts, and audit logs.
\par\textit{A sandbox isolates untrusted code with resource limits, scoped files, network policy, timeouts, and audit logs.}
\paragraph{FC-0649 [hard]} State the main tradeoff and production pitfall for Sandbox runs.
\noindent\textbf{Answer.} Tradeoff: Isolation enables capability while imposing startup and compatibility overhead. Pitfall: A container without kernel, credential, and network boundaries is not a sufficient sandbox.
\par\textit{Tradeoff: Isolation enables capability while imposing startup and compatibility overhead. Pitfall: A container without kernel, credential, and network boundaries is not a sufficient sandbox.}
\paragraph{LA-0650 [hard]} Explain Sandbox runs from first principles, then show how you would evaluate and operate it in production.
\noindent\textbf{Answer.} Start with the mechanism: A sandbox isolates untrusted code with resource limits, scoped files, network policy, timeouts, and audit logs. Make the tradeoff explicit: Isolation enables capability while imposing startup and compatibility overhead. Define workload-specific offline metrics and a latency/cost budget, test important slices, instrument the critical path, and create a rollback criterion. The production review must explicitly guard against this failure: A container without kernel, credential, and network boundaries is not a sufficient sandbox.
\par\textit{A strong answer connects mechanism, measurable tradeoffs, failure isolation, observability, and rollback rather than listing product features.}
\subsection{Orchestration Frameworks}
\paragraph{MCQ-0651 [medium]} Which statement best describes LangGraph?
\begin{enumerate}[label=\Alph*.]
\item Temporal persists workflow history and replays deterministic workflow code around durable activities.
\item AutoGen structures multi-agent conversations and tool-using participants for collaborative task execution.
\item CrewAI organizes role-based agents, tasks, crews, and flows for sequential or coordinated work.
\item LangGraph models long-running agent workflows as stateful graphs with checkpoints, interrupts, streaming, and durable execution.
\end{enumerate}
\noindent\textbf{Answer.} LangGraph models long-running agent workflows as stateful graphs with checkpoints, interrupts, streaming, and durable execution.
\par\textit{LangGraph models long-running agent workflows as stateful graphs with checkpoints, interrupts, streaming, and durable execution. Tradeoff: Fine-grained control aids production reliability but exposes more state and graph design decisions. Watch for: Non-idempotent node side effects make checkpoint replay unsafe.}
\paragraph{MCQ-0652 [hard]} What is the central engineering tradeoff of LangGraph?
\begin{enumerate}[label=\Alph*.]
\item It provides robust retries and long-running execution but requires deterministic workflow discipline.
\item Role decomposition enables complex workflows but multiplies calls, coordination failures, and evaluation cost.
\item Fine-grained control aids production reliability but exposes more state and graph design decisions.
\item It gives accessible orchestration but still needs explicit state, permissions, and observability.
\end{enumerate}
\noindent\textbf{Answer.} Fine-grained control aids production reliability but exposes more state and graph design decisions.
\par\textit{LangGraph models long-running agent workflows as stateful graphs with checkpoints, interrupts, streaming, and durable execution. Tradeoff: Fine-grained control aids production reliability but exposes more state and graph design decisions. Watch for: Non-idempotent node side effects make checkpoint replay unsafe.}
\paragraph{MCQ-0653 [hard]} Which failure mode should an interviewer expect you to identify for LangGraph?
\begin{enumerate}[label=\Alph*.]
\item Adding agents without independent capabilities creates expensive echo chambers.
\item Non-idempotent node side effects make checkpoint replay unsafe.
\item Reading wall-clock time or random state directly in replayed code causes nondeterminism.
\item Persona descriptions cannot substitute for tool-level access control.
\end{enumerate}
\noindent\textbf{Answer.} Non-idempotent node side effects make checkpoint replay unsafe.
\par\textit{LangGraph models long-running agent workflows as stateful graphs with checkpoints, interrupts, streaming, and durable execution. Tradeoff: Fine-grained control aids production reliability but exposes more state and graph design decisions. Watch for: Non-idempotent node side effects make checkpoint replay unsafe.}
\paragraph{MCQ-0654 [medium]} A design review is specifically concerned with this risk: Non-idempotent node side effects make checkpoint replay unsafe. Which topic should be reviewed first?
\begin{enumerate}[label=\Alph*.]
\item Temporal
\item CrewAI
\item LangGraph
\item AutoGen
\end{enumerate}
\noindent\textbf{Answer.} LangGraph
\par\textit{LangGraph models long-running agent workflows as stateful graphs with checkpoints, interrupts, streaming, and durable execution. Tradeoff: Fine-grained control aids production reliability but exposes more state and graph design decisions. Watch for: Non-idempotent node side effects make checkpoint replay unsafe.}
\paragraph{MCQ-0655 [hard]} Which design note most accurately balances the benefits and costs of LangGraph?
\begin{enumerate}[label=\Alph*.]
\item It gives accessible orchestration but still needs explicit state, permissions, and observability.
\item Fine-grained control aids production reliability but exposes more state and graph design decisions.
\item Role decomposition enables complex workflows but multiplies calls, coordination failures, and evaluation cost.
\item It provides robust retries and long-running execution but requires deterministic workflow discipline.
\end{enumerate}
\noindent\textbf{Answer.} Fine-grained control aids production reliability but exposes more state and graph design decisions.
\par\textit{LangGraph models long-running agent workflows as stateful graphs with checkpoints, interrupts, streaming, and durable execution. Tradeoff: Fine-grained control aids production reliability but exposes more state and graph design decisions. Watch for: Non-idempotent node side effects make checkpoint replay unsafe.}
\paragraph{MCQ-0656 [medium]} Which explanation would be strongest in a system-design interview about LangGraph?
\begin{enumerate}[label=\Alph*.]
\item CrewAI organizes role-based agents, tasks, crews, and flows for sequential or coordinated work.
\item AutoGen structures multi-agent conversations and tool-using participants for collaborative task execution.
\item LangGraph models long-running agent workflows as stateful graphs with checkpoints, interrupts, streaming, and durable execution.
\item Temporal persists workflow history and replays deterministic workflow code around durable activities.
\end{enumerate}
\noindent\textbf{Answer.} LangGraph models long-running agent workflows as stateful graphs with checkpoints, interrupts, streaming, and durable execution.
\par\textit{LangGraph models long-running agent workflows as stateful graphs with checkpoints, interrupts, streaming, and durable execution. Tradeoff: Fine-grained control aids production reliability but exposes more state and graph design decisions. Watch for: Non-idempotent node side effects make checkpoint replay unsafe.}
\paragraph{MCQ-0657 [hard]} Which production warning is most directly associated with LangGraph?
\begin{enumerate}[label=\Alph*.]
\item Reading wall-clock time or random state directly in replayed code causes nondeterminism.
\item Adding agents without independent capabilities creates expensive echo chambers.
\item Non-idempotent node side effects make checkpoint replay unsafe.
\item Persona descriptions cannot substitute for tool-level access control.
\end{enumerate}
\noindent\textbf{Answer.} Non-idempotent node side effects make checkpoint replay unsafe.
\par\textit{LangGraph models long-running agent workflows as stateful graphs with checkpoints, interrupts, streaming, and durable execution. Tradeoff: Fine-grained control aids production reliability but exposes more state and graph design decisions. Watch for: Non-idempotent node side effects make checkpoint replay unsafe.}
\paragraph{FC-0658 [medium]} Define LangGraph in interview-ready language.
\noindent\textbf{Answer.} LangGraph models long-running agent workflows as stateful graphs with checkpoints, interrupts, streaming, and durable execution.
\par\textit{LangGraph models long-running agent workflows as stateful graphs with checkpoints, interrupts, streaming, and durable execution.}
\paragraph{FC-0659 [hard]} State the main tradeoff and production pitfall for LangGraph.
\noindent\textbf{Answer.} Tradeoff: Fine-grained control aids production reliability but exposes more state and graph design decisions. Pitfall: Non-idempotent node side effects make checkpoint replay unsafe.
\par\textit{Tradeoff: Fine-grained control aids production reliability but exposes more state and graph design decisions. Pitfall: Non-idempotent node side effects make checkpoint replay unsafe.}
\paragraph{LA-0660 [hard]} Explain LangGraph from first principles, then show how you would evaluate and operate it in production.
\noindent\textbf{Answer.} Start with the mechanism: LangGraph models long-running agent workflows as stateful graphs with checkpoints, interrupts, streaming, and durable execution. Make the tradeoff explicit: Fine-grained control aids production reliability but exposes more state and graph design decisions. Define workload-specific offline metrics and a latency/cost budget, test important slices, instrument the critical path, and create a rollback criterion. The production review must explicitly guard against this failure: Non-idempotent node side effects make checkpoint replay unsafe.
\par\textit{A strong answer connects mechanism, measurable tradeoffs, failure isolation, observability, and rollback rather than listing product features.}
\paragraph{MCQ-0661 [medium]} Which statement best describes LlamaIndex workflows?
\begin{enumerate}[label=\Alph*.]
\item CrewAI organizes role-based agents, tasks, crews, and flows for sequential or coordinated work.
\item General-purpose data orchestrators schedule DAGs, assets, or flows with retries, lineage, and operational UIs.
\item Haystack connects typed components into retrieval, generation, routing, and agent pipelines.
\item Event-driven workflows coordinate retrieval, tools, agents, and structured steps in the LlamaIndex ecosystem.
\end{enumerate}
\noindent\textbf{Answer.} Event-driven workflows coordinate retrieval, tools, agents, and structured steps in the LlamaIndex ecosystem.
\par\textit{Event-driven workflows coordinate retrieval, tools, agents, and structured steps in the LlamaIndex ecosystem. Tradeoff: Strong data and indexing integrations speed RAG work but can couple application flow to framework abstractions. Watch for: Hiding all retrieval behind defaults makes evaluation and migration difficult.}
\paragraph{MCQ-0662 [hard]} What is the central engineering tradeoff of LlamaIndex workflows?
\begin{enumerate}[label=\Alph*.]
\item Explicit components support testing, while custom graphs still require lifecycle discipline.
\item They excel at batch pipelines but are not automatically low-latency interactive agent runtimes.
\item Strong data and indexing integrations speed RAG work but can couple application flow to framework abstractions.
\item It gives accessible orchestration but still needs explicit state, permissions, and observability.
\end{enumerate}
\noindent\textbf{Answer.} Strong data and indexing integrations speed RAG work but can couple application flow to framework abstractions.
\par\textit{Event-driven workflows coordinate retrieval, tools, agents, and structured steps in the LlamaIndex ecosystem. Tradeoff: Strong data and indexing integrations speed RAG work but can couple application flow to framework abstractions. Watch for: Hiding all retrieval behind defaults makes evaluation and migration difficult.}
\paragraph{MCQ-0663 [hard]} Which failure mode should an interviewer expect you to identify for LlamaIndex workflows?
\begin{enumerate}[label=\Alph*.]
\item Running per-token or per-chat-turn work as heavy batch tasks adds unacceptable overhead.
\item Persona descriptions cannot substitute for tool-level access control.
\item Allowing incompatible document schemas across components creates runtime surprises.
\item Hiding all retrieval behind defaults makes evaluation and migration difficult.
\end{enumerate}
\noindent\textbf{Answer.} Hiding all retrieval behind defaults makes evaluation and migration difficult.
\par\textit{Event-driven workflows coordinate retrieval, tools, agents, and structured steps in the LlamaIndex ecosystem. Tradeoff: Strong data and indexing integrations speed RAG work but can couple application flow to framework abstractions. Watch for: Hiding all retrieval behind defaults makes evaluation and migration difficult.}
\paragraph{MCQ-0664 [medium]} A design review is specifically concerned with this risk: Hiding all retrieval behind defaults makes evaluation and migration difficult. Which topic should be reviewed first?
\begin{enumerate}[label=\Alph*.]
\item LlamaIndex workflows
\item Airflow, Dagster, and Prefect
\item CrewAI
\item Haystack pipelines
\end{enumerate}
\noindent\textbf{Answer.} LlamaIndex workflows
\par\textit{Event-driven workflows coordinate retrieval, tools, agents, and structured steps in the LlamaIndex ecosystem. Tradeoff: Strong data and indexing integrations speed RAG work but can couple application flow to framework abstractions. Watch for: Hiding all retrieval behind defaults makes evaluation and migration difficult.}
\paragraph{MCQ-0665 [hard]} Which design note most accurately balances the benefits and costs of LlamaIndex workflows?
\begin{enumerate}[label=\Alph*.]
\item They excel at batch pipelines but are not automatically low-latency interactive agent runtimes.
\item Strong data and indexing integrations speed RAG work but can couple application flow to framework abstractions.
\item It gives accessible orchestration but still needs explicit state, permissions, and observability.
\item Explicit components support testing, while custom graphs still require lifecycle discipline.
\end{enumerate}
\noindent\textbf{Answer.} Strong data and indexing integrations speed RAG work but can couple application flow to framework abstractions.
\par\textit{Event-driven workflows coordinate retrieval, tools, agents, and structured steps in the LlamaIndex ecosystem. Tradeoff: Strong data and indexing integrations speed RAG work but can couple application flow to framework abstractions. Watch for: Hiding all retrieval behind defaults makes evaluation and migration difficult.}
\paragraph{MCQ-0666 [medium]} Which explanation would be strongest in a system-design interview about LlamaIndex workflows?
\begin{enumerate}[label=\Alph*.]
\item CrewAI organizes role-based agents, tasks, crews, and flows for sequential or coordinated work.
\item General-purpose data orchestrators schedule DAGs, assets, or flows with retries, lineage, and operational UIs.
\item Haystack connects typed components into retrieval, generation, routing, and agent pipelines.
\item Event-driven workflows coordinate retrieval, tools, agents, and structured steps in the LlamaIndex ecosystem.
\end{enumerate}
\noindent\textbf{Answer.} Event-driven workflows coordinate retrieval, tools, agents, and structured steps in the LlamaIndex ecosystem.
\par\textit{Event-driven workflows coordinate retrieval, tools, agents, and structured steps in the LlamaIndex ecosystem. Tradeoff: Strong data and indexing integrations speed RAG work but can couple application flow to framework abstractions. Watch for: Hiding all retrieval behind defaults makes evaluation and migration difficult.}
\paragraph{MCQ-0667 [hard]} Which production warning is most directly associated with LlamaIndex workflows?
\begin{enumerate}[label=\Alph*.]
\item Persona descriptions cannot substitute for tool-level access control.
\item Allowing incompatible document schemas across components creates runtime surprises.
\item Hiding all retrieval behind defaults makes evaluation and migration difficult.
\item Running per-token or per-chat-turn work as heavy batch tasks adds unacceptable overhead.
\end{enumerate}
\noindent\textbf{Answer.} Hiding all retrieval behind defaults makes evaluation and migration difficult.
\par\textit{Event-driven workflows coordinate retrieval, tools, agents, and structured steps in the LlamaIndex ecosystem. Tradeoff: Strong data and indexing integrations speed RAG work but can couple application flow to framework abstractions. Watch for: Hiding all retrieval behind defaults makes evaluation and migration difficult.}
\paragraph{FC-0668 [medium]} Define LlamaIndex workflows in interview-ready language.
\noindent\textbf{Answer.} Event-driven workflows coordinate retrieval, tools, agents, and structured steps in the LlamaIndex ecosystem.
\par\textit{Event-driven workflows coordinate retrieval, tools, agents, and structured steps in the LlamaIndex ecosystem.}
\paragraph{FC-0669 [hard]} State the main tradeoff and production pitfall for LlamaIndex workflows.
\noindent\textbf{Answer.} Tradeoff: Strong data and indexing integrations speed RAG work but can couple application flow to framework abstractions. Pitfall: Hiding all retrieval behind defaults makes evaluation and migration difficult.
\par\textit{Tradeoff: Strong data and indexing integrations speed RAG work but can couple application flow to framework abstractions. Pitfall: Hiding all retrieval behind defaults makes evaluation and migration difficult.}
\paragraph{LA-0670 [hard]} Explain LlamaIndex workflows from first principles, then show how you would evaluate and operate it in production.
\noindent\textbf{Answer.} Start with the mechanism: Event-driven workflows coordinate retrieval, tools, agents, and structured steps in the LlamaIndex ecosystem. Make the tradeoff explicit: Strong data and indexing integrations speed RAG work but can couple application flow to framework abstractions. Define workload-specific offline metrics and a latency/cost budget, test important slices, instrument the critical path, and create a rollback criterion. The production review must explicitly guard against this failure: Hiding all retrieval behind defaults makes evaluation and migration difficult.
\par\textit{A strong answer connects mechanism, measurable tradeoffs, failure isolation, observability, and rollback rather than listing product features.}
\paragraph{MCQ-0671 [medium]} Which statement best describes Haystack pipelines?
\begin{enumerate}[label=\Alph*.]
\item Event-driven workflows coordinate retrieval, tools, agents, and structured steps in the LlamaIndex ecosystem.
\item General-purpose data orchestrators schedule DAGs, assets, or flows with retries, lineage, and operational UIs.
\item Haystack connects typed components into retrieval, generation, routing, and agent pipelines.
\item AutoGen structures multi-agent conversations and tool-using participants for collaborative task execution.
\end{enumerate}
\noindent\textbf{Answer.} Haystack connects typed components into retrieval, generation, routing, and agent pipelines.
\par\textit{Haystack connects typed components into retrieval, generation, routing, and agent pipelines. Tradeoff: Explicit components support testing, while custom graphs still require lifecycle discipline. Watch for: Allowing incompatible document schemas across components creates runtime surprises.}
\paragraph{MCQ-0672 [hard]} What is the central engineering tradeoff of Haystack pipelines?
\begin{enumerate}[label=\Alph*.]
\item Explicit components support testing, while custom graphs still require lifecycle discipline.
\item Role decomposition enables complex workflows but multiplies calls, coordination failures, and evaluation cost.
\item They excel at batch pipelines but are not automatically low-latency interactive agent runtimes.
\item Strong data and indexing integrations speed RAG work but can couple application flow to framework abstractions.
\end{enumerate}
\noindent\textbf{Answer.} Explicit components support testing, while custom graphs still require lifecycle discipline.
\par\textit{Haystack connects typed components into retrieval, generation, routing, and agent pipelines. Tradeoff: Explicit components support testing, while custom graphs still require lifecycle discipline. Watch for: Allowing incompatible document schemas across components creates runtime surprises.}
\paragraph{MCQ-0673 [hard]} Which failure mode should an interviewer expect you to identify for Haystack pipelines?
\begin{enumerate}[label=\Alph*.]
\item Adding agents without independent capabilities creates expensive echo chambers.
\item Hiding all retrieval behind defaults makes evaluation and migration difficult.
\item Running per-token or per-chat-turn work as heavy batch tasks adds unacceptable overhead.
\item Allowing incompatible document schemas across components creates runtime surprises.
\end{enumerate}
\noindent\textbf{Answer.} Allowing incompatible document schemas across components creates runtime surprises.
\par\textit{Haystack connects typed components into retrieval, generation, routing, and agent pipelines. Tradeoff: Explicit components support testing, while custom graphs still require lifecycle discipline. Watch for: Allowing incompatible document schemas across components creates runtime surprises.}
\paragraph{MCQ-0674 [medium]} A design review is specifically concerned with this risk: Allowing incompatible document schemas across components creates runtime surprises. Which topic should be reviewed first?
\begin{enumerate}[label=\Alph*.]
\item LlamaIndex workflows
\item Airflow, Dagster, and Prefect
\item AutoGen
\item Haystack pipelines
\end{enumerate}
\noindent\textbf{Answer.} Haystack pipelines
\par\textit{Haystack connects typed components into retrieval, generation, routing, and agent pipelines. Tradeoff: Explicit components support testing, while custom graphs still require lifecycle discipline. Watch for: Allowing incompatible document schemas across components creates runtime surprises.}
\paragraph{MCQ-0675 [hard]} Which design note most accurately balances the benefits and costs of Haystack pipelines?
\begin{enumerate}[label=\Alph*.]
\item Explicit components support testing, while custom graphs still require lifecycle discipline.
\item Role decomposition enables complex workflows but multiplies calls, coordination failures, and evaluation cost.
\item Strong data and indexing integrations speed RAG work but can couple application flow to framework abstractions.
\item They excel at batch pipelines but are not automatically low-latency interactive agent runtimes.
\end{enumerate}
\noindent\textbf{Answer.} Explicit components support testing, while custom graphs still require lifecycle discipline.
\par\textit{Haystack connects typed components into retrieval, generation, routing, and agent pipelines. Tradeoff: Explicit components support testing, while custom graphs still require lifecycle discipline. Watch for: Allowing incompatible document schemas across components creates runtime surprises.}
\paragraph{MCQ-0676 [medium]} Which explanation would be strongest in a system-design interview about Haystack pipelines?
\begin{enumerate}[label=\Alph*.]
\item Event-driven workflows coordinate retrieval, tools, agents, and structured steps in the LlamaIndex ecosystem.
\item General-purpose data orchestrators schedule DAGs, assets, or flows with retries, lineage, and operational UIs.
\item AutoGen structures multi-agent conversations and tool-using participants for collaborative task execution.
\item Haystack connects typed components into retrieval, generation, routing, and agent pipelines.
\end{enumerate}
\noindent\textbf{Answer.} Haystack connects typed components into retrieval, generation, routing, and agent pipelines.
\par\textit{Haystack connects typed components into retrieval, generation, routing, and agent pipelines. Tradeoff: Explicit components support testing, while custom graphs still require lifecycle discipline. Watch for: Allowing incompatible document schemas across components creates runtime surprises.}
\paragraph{MCQ-0677 [hard]} Which production warning is most directly associated with Haystack pipelines?
\begin{enumerate}[label=\Alph*.]
\item Allowing incompatible document schemas across components creates runtime surprises.
\item Hiding all retrieval behind defaults makes evaluation and migration difficult.
\item Adding agents without independent capabilities creates expensive echo chambers.
\item Running per-token or per-chat-turn work as heavy batch tasks adds unacceptable overhead.
\end{enumerate}
\noindent\textbf{Answer.} Allowing incompatible document schemas across components creates runtime surprises.
\par\textit{Haystack connects typed components into retrieval, generation, routing, and agent pipelines. Tradeoff: Explicit components support testing, while custom graphs still require lifecycle discipline. Watch for: Allowing incompatible document schemas across components creates runtime surprises.}
\paragraph{FC-0678 [medium]} Define Haystack pipelines in interview-ready language.
\noindent\textbf{Answer.} Haystack connects typed components into retrieval, generation, routing, and agent pipelines.
\par\textit{Haystack connects typed components into retrieval, generation, routing, and agent pipelines.}
\paragraph{FC-0679 [hard]} State the main tradeoff and production pitfall for Haystack pipelines.
\noindent\textbf{Answer.} Tradeoff: Explicit components support testing, while custom graphs still require lifecycle discipline. Pitfall: Allowing incompatible document schemas across components creates runtime surprises.
\par\textit{Tradeoff: Explicit components support testing, while custom graphs still require lifecycle discipline. Pitfall: Allowing incompatible document schemas across components creates runtime surprises.}
\paragraph{LA-0680 [hard]} Explain Haystack pipelines from first principles, then show how you would evaluate and operate it in production.
\noindent\textbf{Answer.} Start with the mechanism: Haystack connects typed components into retrieval, generation, routing, and agent pipelines. Make the tradeoff explicit: Explicit components support testing, while custom graphs still require lifecycle discipline. Define workload-specific offline metrics and a latency/cost budget, test important slices, instrument the critical path, and create a rollback criterion. The production review must explicitly guard against this failure: Allowing incompatible document schemas across components creates runtime surprises.
\par\textit{A strong answer connects mechanism, measurable tradeoffs, failure isolation, observability, and rollback rather than listing product features.}
\paragraph{MCQ-0681 [medium]} Which statement best describes PydanticAI?
\begin{enumerate}[label=\Alph*.]
\item Haystack connects typed components into retrieval, generation, routing, and agent pipelines.
\item AutoGen structures multi-agent conversations and tool-using participants for collaborative task execution.
\item PydanticAI uses typed dependencies, tools, results, and validation to build model-agnostic Python agents.
\item Event-driven workflows coordinate retrieval, tools, agents, and structured steps in the LlamaIndex ecosystem.
\end{enumerate}
\noindent\textbf{Answer.} PydanticAI uses typed dependencies, tools, results, and validation to build model-agnostic Python agents.
\par\textit{PydanticAI uses typed dependencies, tools, results, and validation to build model-agnostic Python agents. Tradeoff: Type safety improves integration but cannot validate semantic truth. Watch for: Confusing schema conformance with a correct answer produces false confidence.}
\paragraph{MCQ-0682 [hard]} What is the central engineering tradeoff of PydanticAI?
\begin{enumerate}[label=\Alph*.]
\item Strong data and indexing integrations speed RAG work but can couple application flow to framework abstractions.
\item Type safety improves integration but cannot validate semantic truth.
\item Explicit components support testing, while custom graphs still require lifecycle discipline.
\item Role decomposition enables complex workflows but multiplies calls, coordination failures, and evaluation cost.
\end{enumerate}
\noindent\textbf{Answer.} Type safety improves integration but cannot validate semantic truth.
\par\textit{PydanticAI uses typed dependencies, tools, results, and validation to build model-agnostic Python agents. Tradeoff: Type safety improves integration but cannot validate semantic truth. Watch for: Confusing schema conformance with a correct answer produces false confidence.}
\paragraph{MCQ-0683 [hard]} Which failure mode should an interviewer expect you to identify for PydanticAI?
\begin{enumerate}[label=\Alph*.]
\item Confusing schema conformance with a correct answer produces false confidence.
\item Hiding all retrieval behind defaults makes evaluation and migration difficult.
\item Allowing incompatible document schemas across components creates runtime surprises.
\item Adding agents without independent capabilities creates expensive echo chambers.
\end{enumerate}
\noindent\textbf{Answer.} Confusing schema conformance with a correct answer produces false confidence.
\par\textit{PydanticAI uses typed dependencies, tools, results, and validation to build model-agnostic Python agents. Tradeoff: Type safety improves integration but cannot validate semantic truth. Watch for: Confusing schema conformance with a correct answer produces false confidence.}
\paragraph{MCQ-0684 [medium]} A design review is specifically concerned with this risk: Confusing schema conformance with a correct answer produces false confidence. Which topic should be reviewed first?
\begin{enumerate}[label=\Alph*.]
\item LlamaIndex workflows
\item AutoGen
\item Haystack pipelines
\item PydanticAI
\end{enumerate}
\noindent\textbf{Answer.} PydanticAI
\par\textit{PydanticAI uses typed dependencies, tools, results, and validation to build model-agnostic Python agents. Tradeoff: Type safety improves integration but cannot validate semantic truth. Watch for: Confusing schema conformance with a correct answer produces false confidence.}
\paragraph{MCQ-0685 [hard]} Which design note most accurately balances the benefits and costs of PydanticAI?
\begin{enumerate}[label=\Alph*.]
\item Role decomposition enables complex workflows but multiplies calls, coordination failures, and evaluation cost.
\item Explicit components support testing, while custom graphs still require lifecycle discipline.
\item Type safety improves integration but cannot validate semantic truth.
\item Strong data and indexing integrations speed RAG work but can couple application flow to framework abstractions.
\end{enumerate}
\noindent\textbf{Answer.} Type safety improves integration but cannot validate semantic truth.
\par\textit{PydanticAI uses typed dependencies, tools, results, and validation to build model-agnostic Python agents. Tradeoff: Type safety improves integration but cannot validate semantic truth. Watch for: Confusing schema conformance with a correct answer produces false confidence.}
\paragraph{MCQ-0686 [medium]} Which explanation would be strongest in a system-design interview about PydanticAI?
\begin{enumerate}[label=\Alph*.]
\item AutoGen structures multi-agent conversations and tool-using participants for collaborative task execution.
\item Haystack connects typed components into retrieval, generation, routing, and agent pipelines.
\item PydanticAI uses typed dependencies, tools, results, and validation to build model-agnostic Python agents.
\item Event-driven workflows coordinate retrieval, tools, agents, and structured steps in the LlamaIndex ecosystem.
\end{enumerate}
\noindent\textbf{Answer.} PydanticAI uses typed dependencies, tools, results, and validation to build model-agnostic Python agents.
\par\textit{PydanticAI uses typed dependencies, tools, results, and validation to build model-agnostic Python agents. Tradeoff: Type safety improves integration but cannot validate semantic truth. Watch for: Confusing schema conformance with a correct answer produces false confidence.}
\paragraph{MCQ-0687 [hard]} Which production warning is most directly associated with PydanticAI?
\begin{enumerate}[label=\Alph*.]
\item Confusing schema conformance with a correct answer produces false confidence.
\item Adding agents without independent capabilities creates expensive echo chambers.
\item Allowing incompatible document schemas across components creates runtime surprises.
\item Hiding all retrieval behind defaults makes evaluation and migration difficult.
\end{enumerate}
\noindent\textbf{Answer.} Confusing schema conformance with a correct answer produces false confidence.
\par\textit{PydanticAI uses typed dependencies, tools, results, and validation to build model-agnostic Python agents. Tradeoff: Type safety improves integration but cannot validate semantic truth. Watch for: Confusing schema conformance with a correct answer produces false confidence.}
\paragraph{FC-0688 [medium]} Define PydanticAI in interview-ready language.
\noindent\textbf{Answer.} PydanticAI uses typed dependencies, tools, results, and validation to build model-agnostic Python agents.
\par\textit{PydanticAI uses typed dependencies, tools, results, and validation to build model-agnostic Python agents.}
\paragraph{FC-0689 [hard]} State the main tradeoff and production pitfall for PydanticAI.
\noindent\textbf{Answer.} Tradeoff: Type safety improves integration but cannot validate semantic truth. Pitfall: Confusing schema conformance with a correct answer produces false confidence.
\par\textit{Tradeoff: Type safety improves integration but cannot validate semantic truth. Pitfall: Confusing schema conformance with a correct answer produces false confidence.}
\paragraph{LA-0690 [hard]} Explain PydanticAI from first principles, then show how you would evaluate and operate it in production.
\noindent\textbf{Answer.} Start with the mechanism: PydanticAI uses typed dependencies, tools, results, and validation to build model-agnostic Python agents. Make the tradeoff explicit: Type safety improves integration but cannot validate semantic truth. Define workload-specific offline metrics and a latency/cost budget, test important slices, instrument the critical path, and create a rollback criterion. The production review must explicitly guard against this failure: Confusing schema conformance with a correct answer produces false confidence.
\par\textit{A strong answer connects mechanism, measurable tradeoffs, failure isolation, observability, and rollback rather than listing product features.}
\paragraph{MCQ-0691 [medium]} Which statement best describes Semantic Kernel?
\begin{enumerate}[label=\Alph*.]
\item Haystack connects typed components into retrieval, generation, routing, and agent pipelines.
\item Semantic Kernel offers planners, agents, plugins, memory connectors, and multi-language enterprise integration.
\item Temporal persists workflow history and replays deterministic workflow code around durable activities.
\item CrewAI organizes role-based agents, tasks, crews, and flows for sequential or coordinated work.
\end{enumerate}
\noindent\textbf{Answer.} Semantic Kernel offers planners, agents, plugins, memory connectors, and multi-language enterprise integration.
\par\textit{Semantic Kernel offers planners, agents, plugins, memory connectors, and multi-language enterprise integration. Tradeoff: Its broad integration surface suits platform teams but introduces framework concepts and version coupling. Watch for: Registering overly broad plugins grants the model unnecessary authority.}
\paragraph{MCQ-0692 [hard]} What is the central engineering tradeoff of Semantic Kernel?
\begin{enumerate}[label=\Alph*.]
\item It provides robust retries and long-running execution but requires deterministic workflow discipline.
\item It gives accessible orchestration but still needs explicit state, permissions, and observability.
\item Its broad integration surface suits platform teams but introduces framework concepts and version coupling.
\item Explicit components support testing, while custom graphs still require lifecycle discipline.
\end{enumerate}
\noindent\textbf{Answer.} Its broad integration surface suits platform teams but introduces framework concepts and version coupling.
\par\textit{Semantic Kernel offers planners, agents, plugins, memory connectors, and multi-language enterprise integration. Tradeoff: Its broad integration surface suits platform teams but introduces framework concepts and version coupling. Watch for: Registering overly broad plugins grants the model unnecessary authority.}
\paragraph{MCQ-0693 [hard]} Which failure mode should an interviewer expect you to identify for Semantic Kernel?
\begin{enumerate}[label=\Alph*.]
\item Reading wall-clock time or random state directly in replayed code causes nondeterminism.
\item Persona descriptions cannot substitute for tool-level access control.
\item Allowing incompatible document schemas across components creates runtime surprises.
\item Registering overly broad plugins grants the model unnecessary authority.
\end{enumerate}
\noindent\textbf{Answer.} Registering overly broad plugins grants the model unnecessary authority.
\par\textit{Semantic Kernel offers planners, agents, plugins, memory connectors, and multi-language enterprise integration. Tradeoff: Its broad integration surface suits platform teams but introduces framework concepts and version coupling. Watch for: Registering overly broad plugins grants the model unnecessary authority.}
\paragraph{MCQ-0694 [medium]} A design review is specifically concerned with this risk: Registering overly broad plugins grants the model unnecessary authority. Which topic should be reviewed first?
\begin{enumerate}[label=\Alph*.]
\item Haystack pipelines
\item Temporal
\item CrewAI
\item Semantic Kernel
\end{enumerate}
\noindent\textbf{Answer.} Semantic Kernel
\par\textit{Semantic Kernel offers planners, agents, plugins, memory connectors, and multi-language enterprise integration. Tradeoff: Its broad integration surface suits platform teams but introduces framework concepts and version coupling. Watch for: Registering overly broad plugins grants the model unnecessary authority.}
\paragraph{MCQ-0695 [hard]} Which design note most accurately balances the benefits and costs of Semantic Kernel?
\begin{enumerate}[label=\Alph*.]
\item Its broad integration surface suits platform teams but introduces framework concepts and version coupling.
\item Explicit components support testing, while custom graphs still require lifecycle discipline.
\item It gives accessible orchestration but still needs explicit state, permissions, and observability.
\item It provides robust retries and long-running execution but requires deterministic workflow discipline.
\end{enumerate}
\noindent\textbf{Answer.} Its broad integration surface suits platform teams but introduces framework concepts and version coupling.
\par\textit{Semantic Kernel offers planners, agents, plugins, memory connectors, and multi-language enterprise integration. Tradeoff: Its broad integration surface suits platform teams but introduces framework concepts and version coupling. Watch for: Registering overly broad plugins grants the model unnecessary authority.}
\paragraph{MCQ-0696 [medium]} Which explanation would be strongest in a system-design interview about Semantic Kernel?
\begin{enumerate}[label=\Alph*.]
\item Temporal persists workflow history and replays deterministic workflow code around durable activities.
\item Semantic Kernel offers planners, agents, plugins, memory connectors, and multi-language enterprise integration.
\item CrewAI organizes role-based agents, tasks, crews, and flows for sequential or coordinated work.
\item Haystack connects typed components into retrieval, generation, routing, and agent pipelines.
\end{enumerate}
\noindent\textbf{Answer.} Semantic Kernel offers planners, agents, plugins, memory connectors, and multi-language enterprise integration.
\par\textit{Semantic Kernel offers planners, agents, plugins, memory connectors, and multi-language enterprise integration. Tradeoff: Its broad integration surface suits platform teams but introduces framework concepts and version coupling. Watch for: Registering overly broad plugins grants the model unnecessary authority.}
\paragraph{MCQ-0697 [hard]} Which production warning is most directly associated with Semantic Kernel?
\begin{enumerate}[label=\Alph*.]
\item Registering overly broad plugins grants the model unnecessary authority.
\item Reading wall-clock time or random state directly in replayed code causes nondeterminism.
\item Allowing incompatible document schemas across components creates runtime surprises.
\item Persona descriptions cannot substitute for tool-level access control.
\end{enumerate}
\noindent\textbf{Answer.} Registering overly broad plugins grants the model unnecessary authority.
\par\textit{Semantic Kernel offers planners, agents, plugins, memory connectors, and multi-language enterprise integration. Tradeoff: Its broad integration surface suits platform teams but introduces framework concepts and version coupling. Watch for: Registering overly broad plugins grants the model unnecessary authority.}
\paragraph{FC-0698 [medium]} Define Semantic Kernel in interview-ready language.
\noindent\textbf{Answer.} Semantic Kernel offers planners, agents, plugins, memory connectors, and multi-language enterprise integration.
\par\textit{Semantic Kernel offers planners, agents, plugins, memory connectors, and multi-language enterprise integration.}
\paragraph{FC-0699 [hard]} State the main tradeoff and production pitfall for Semantic Kernel.
\noindent\textbf{Answer.} Tradeoff: Its broad integration surface suits platform teams but introduces framework concepts and version coupling. Pitfall: Registering overly broad plugins grants the model unnecessary authority.
\par\textit{Tradeoff: Its broad integration surface suits platform teams but introduces framework concepts and version coupling. Pitfall: Registering overly broad plugins grants the model unnecessary authority.}
\paragraph{LA-0700 [hard]} Explain Semantic Kernel from first principles, then show how you would evaluate and operate it in production.
\noindent\textbf{Answer.} Start with the mechanism: Semantic Kernel offers planners, agents, plugins, memory connectors, and multi-language enterprise integration. Make the tradeoff explicit: Its broad integration surface suits platform teams but introduces framework concepts and version coupling. Define workload-specific offline metrics and a latency/cost budget, test important slices, instrument the critical path, and create a rollback criterion. The production review must explicitly guard against this failure: Registering overly broad plugins grants the model unnecessary authority.
\par\textit{A strong answer connects mechanism, measurable tradeoffs, failure isolation, observability, and rollback rather than listing product features.}
\paragraph{MCQ-0701 [medium]} Which statement best describes AutoGen?
\begin{enumerate}[label=\Alph*.]
\item Haystack connects typed components into retrieval, generation, routing, and agent pipelines.
\item CrewAI organizes role-based agents, tasks, crews, and flows for sequential or coordinated work.
\item LangGraph models long-running agent workflows as stateful graphs with checkpoints, interrupts, streaming, and durable execution.
\item AutoGen structures multi-agent conversations and tool-using participants for collaborative task execution.
\end{enumerate}
\noindent\textbf{Answer.} AutoGen structures multi-agent conversations and tool-using participants for collaborative task execution.
\par\textit{AutoGen structures multi-agent conversations and tool-using participants for collaborative task execution. Tradeoff: Role decomposition enables complex workflows but multiplies calls, coordination failures, and evaluation cost. Watch for: Adding agents without independent capabilities creates expensive echo chambers.}
\paragraph{MCQ-0702 [hard]} What is the central engineering tradeoff of AutoGen?
\begin{enumerate}[label=\Alph*.]
\item Role decomposition enables complex workflows but multiplies calls, coordination failures, and evaluation cost.
\item Fine-grained control aids production reliability but exposes more state and graph design decisions.
\item It gives accessible orchestration but still needs explicit state, permissions, and observability.
\item Explicit components support testing, while custom graphs still require lifecycle discipline.
\end{enumerate}
\noindent\textbf{Answer.} Role decomposition enables complex workflows but multiplies calls, coordination failures, and evaluation cost.
\par\textit{AutoGen structures multi-agent conversations and tool-using participants for collaborative task execution. Tradeoff: Role decomposition enables complex workflows but multiplies calls, coordination failures, and evaluation cost. Watch for: Adding agents without independent capabilities creates expensive echo chambers.}
\paragraph{MCQ-0703 [hard]} Which failure mode should an interviewer expect you to identify for AutoGen?
\begin{enumerate}[label=\Alph*.]
\item Persona descriptions cannot substitute for tool-level access control.
\item Adding agents without independent capabilities creates expensive echo chambers.
\item Non-idempotent node side effects make checkpoint replay unsafe.
\item Allowing incompatible document schemas across components creates runtime surprises.
\end{enumerate}
\noindent\textbf{Answer.} Adding agents without independent capabilities creates expensive echo chambers.
\par\textit{AutoGen structures multi-agent conversations and tool-using participants for collaborative task execution. Tradeoff: Role decomposition enables complex workflows but multiplies calls, coordination failures, and evaluation cost. Watch for: Adding agents without independent capabilities creates expensive echo chambers.}
\paragraph{MCQ-0704 [medium]} A design review is specifically concerned with this risk: Adding agents without independent capabilities creates expensive echo chambers. Which topic should be reviewed first?
\begin{enumerate}[label=\Alph*.]
\item AutoGen
\item LangGraph
\item CrewAI
\item Haystack pipelines
\end{enumerate}
\noindent\textbf{Answer.} AutoGen
\par\textit{AutoGen structures multi-agent conversations and tool-using participants for collaborative task execution. Tradeoff: Role decomposition enables complex workflows but multiplies calls, coordination failures, and evaluation cost. Watch for: Adding agents without independent capabilities creates expensive echo chambers.}
\paragraph{MCQ-0705 [hard]} Which design note most accurately balances the benefits and costs of AutoGen?
\begin{enumerate}[label=\Alph*.]
\item Role decomposition enables complex workflows but multiplies calls, coordination failures, and evaluation cost.
\item Explicit components support testing, while custom graphs still require lifecycle discipline.
\item Fine-grained control aids production reliability but exposes more state and graph design decisions.
\item It gives accessible orchestration but still needs explicit state, permissions, and observability.
\end{enumerate}
\noindent\textbf{Answer.} Role decomposition enables complex workflows but multiplies calls, coordination failures, and evaluation cost.
\par\textit{AutoGen structures multi-agent conversations and tool-using participants for collaborative task execution. Tradeoff: Role decomposition enables complex workflows but multiplies calls, coordination failures, and evaluation cost. Watch for: Adding agents without independent capabilities creates expensive echo chambers.}
\paragraph{MCQ-0706 [medium]} Which explanation would be strongest in a system-design interview about AutoGen?
\begin{enumerate}[label=\Alph*.]
\item LangGraph models long-running agent workflows as stateful graphs with checkpoints, interrupts, streaming, and durable execution.
\item Haystack connects typed components into retrieval, generation, routing, and agent pipelines.
\item CrewAI organizes role-based agents, tasks, crews, and flows for sequential or coordinated work.
\item AutoGen structures multi-agent conversations and tool-using participants for collaborative task execution.
\end{enumerate}
\noindent\textbf{Answer.} AutoGen structures multi-agent conversations and tool-using participants for collaborative task execution.
\par\textit{AutoGen structures multi-agent conversations and tool-using participants for collaborative task execution. Tradeoff: Role decomposition enables complex workflows but multiplies calls, coordination failures, and evaluation cost. Watch for: Adding agents without independent capabilities creates expensive echo chambers.}
\paragraph{MCQ-0707 [hard]} Which production warning is most directly associated with AutoGen?
\begin{enumerate}[label=\Alph*.]
\item Persona descriptions cannot substitute for tool-level access control.
\item Non-idempotent node side effects make checkpoint replay unsafe.
\item Allowing incompatible document schemas across components creates runtime surprises.
\item Adding agents without independent capabilities creates expensive echo chambers.
\end{enumerate}
\noindent\textbf{Answer.} Adding agents without independent capabilities creates expensive echo chambers.
\par\textit{AutoGen structures multi-agent conversations and tool-using participants for collaborative task execution. Tradeoff: Role decomposition enables complex workflows but multiplies calls, coordination failures, and evaluation cost. Watch for: Adding agents without independent capabilities creates expensive echo chambers.}
\paragraph{FC-0708 [medium]} Define AutoGen in interview-ready language.
\noindent\textbf{Answer.} AutoGen structures multi-agent conversations and tool-using participants for collaborative task execution.
\par\textit{AutoGen structures multi-agent conversations and tool-using participants for collaborative task execution.}
\paragraph{FC-0709 [hard]} State the main tradeoff and production pitfall for AutoGen.
\noindent\textbf{Answer.} Tradeoff: Role decomposition enables complex workflows but multiplies calls, coordination failures, and evaluation cost. Pitfall: Adding agents without independent capabilities creates expensive echo chambers.
\par\textit{Tradeoff: Role decomposition enables complex workflows but multiplies calls, coordination failures, and evaluation cost. Pitfall: Adding agents without independent capabilities creates expensive echo chambers.}
\paragraph{LA-0710 [hard]} Explain AutoGen from first principles, then show how you would evaluate and operate it in production.
\noindent\textbf{Answer.} Start with the mechanism: AutoGen structures multi-agent conversations and tool-using participants for collaborative task execution. Make the tradeoff explicit: Role decomposition enables complex workflows but multiplies calls, coordination failures, and evaluation cost. Define workload-specific offline metrics and a latency/cost budget, test important slices, instrument the critical path, and create a rollback criterion. The production review must explicitly guard against this failure: Adding agents without independent capabilities creates expensive echo chambers.
\par\textit{A strong answer connects mechanism, measurable tradeoffs, failure isolation, observability, and rollback rather than listing product features.}
\paragraph{MCQ-0711 [medium]} Which statement best describes CrewAI?
\begin{enumerate}[label=\Alph*.]
\item Semantic Kernel offers planners, agents, plugins, memory connectors, and multi-language enterprise integration.
\item General-purpose data orchestrators schedule DAGs, assets, or flows with retries, lineage, and operational UIs.
\item Haystack connects typed components into retrieval, generation, routing, and agent pipelines.
\item CrewAI organizes role-based agents, tasks, crews, and flows for sequential or coordinated work.
\end{enumerate}
\noindent\textbf{Answer.} CrewAI organizes role-based agents, tasks, crews, and flows for sequential or coordinated work.
\par\textit{CrewAI organizes role-based agents, tasks, crews, and flows for sequential or coordinated work. Tradeoff: It gives accessible orchestration but still needs explicit state, permissions, and observability. Watch for: Persona descriptions cannot substitute for tool-level access control.}
\paragraph{MCQ-0712 [hard]} What is the central engineering tradeoff of CrewAI?
\begin{enumerate}[label=\Alph*.]
\item It gives accessible orchestration but still needs explicit state, permissions, and observability.
\item Explicit components support testing, while custom graphs still require lifecycle discipline.
\item They excel at batch pipelines but are not automatically low-latency interactive agent runtimes.
\item Its broad integration surface suits platform teams but introduces framework concepts and version coupling.
\end{enumerate}
\noindent\textbf{Answer.} It gives accessible orchestration but still needs explicit state, permissions, and observability.
\par\textit{CrewAI organizes role-based agents, tasks, crews, and flows for sequential or coordinated work. Tradeoff: It gives accessible orchestration but still needs explicit state, permissions, and observability. Watch for: Persona descriptions cannot substitute for tool-level access control.}
\paragraph{MCQ-0713 [hard]} Which failure mode should an interviewer expect you to identify for CrewAI?
\begin{enumerate}[label=\Alph*.]
\item Registering overly broad plugins grants the model unnecessary authority.
\item Running per-token or per-chat-turn work as heavy batch tasks adds unacceptable overhead.
\item Persona descriptions cannot substitute for tool-level access control.
\item Allowing incompatible document schemas across components creates runtime surprises.
\end{enumerate}
\noindent\textbf{Answer.} Persona descriptions cannot substitute for tool-level access control.
\par\textit{CrewAI organizes role-based agents, tasks, crews, and flows for sequential or coordinated work. Tradeoff: It gives accessible orchestration but still needs explicit state, permissions, and observability. Watch for: Persona descriptions cannot substitute for tool-level access control.}
\paragraph{MCQ-0714 [medium]} A design review is specifically concerned with this risk: Persona descriptions cannot substitute for tool-level access control. Which topic should be reviewed first?
\begin{enumerate}[label=\Alph*.]
\item Haystack pipelines
\item CrewAI
\item Airflow, Dagster, and Prefect
\item Semantic Kernel
\end{enumerate}
\noindent\textbf{Answer.} CrewAI
\par\textit{CrewAI organizes role-based agents, tasks, crews, and flows for sequential or coordinated work. Tradeoff: It gives accessible orchestration but still needs explicit state, permissions, and observability. Watch for: Persona descriptions cannot substitute for tool-level access control.}
\paragraph{MCQ-0715 [hard]} Which design note most accurately balances the benefits and costs of CrewAI?
\begin{enumerate}[label=\Alph*.]
\item It gives accessible orchestration but still needs explicit state, permissions, and observability.
\item Explicit components support testing, while custom graphs still require lifecycle discipline.
\item They excel at batch pipelines but are not automatically low-latency interactive agent runtimes.
\item Its broad integration surface suits platform teams but introduces framework concepts and version coupling.
\end{enumerate}
\noindent\textbf{Answer.} It gives accessible orchestration but still needs explicit state, permissions, and observability.
\par\textit{CrewAI organizes role-based agents, tasks, crews, and flows for sequential or coordinated work. Tradeoff: It gives accessible orchestration but still needs explicit state, permissions, and observability. Watch for: Persona descriptions cannot substitute for tool-level access control.}
\paragraph{MCQ-0716 [medium]} Which explanation would be strongest in a system-design interview about CrewAI?
\begin{enumerate}[label=\Alph*.]
\item Haystack connects typed components into retrieval, generation, routing, and agent pipelines.
\item CrewAI organizes role-based agents, tasks, crews, and flows for sequential or coordinated work.
\item General-purpose data orchestrators schedule DAGs, assets, or flows with retries, lineage, and operational UIs.
\item Semantic Kernel offers planners, agents, plugins, memory connectors, and multi-language enterprise integration.
\end{enumerate}
\noindent\textbf{Answer.} CrewAI organizes role-based agents, tasks, crews, and flows for sequential or coordinated work.
\par\textit{CrewAI organizes role-based agents, tasks, crews, and flows for sequential or coordinated work. Tradeoff: It gives accessible orchestration but still needs explicit state, permissions, and observability. Watch for: Persona descriptions cannot substitute for tool-level access control.}
\paragraph{MCQ-0717 [hard]} Which production warning is most directly associated with CrewAI?
\begin{enumerate}[label=\Alph*.]
\item Running per-token or per-chat-turn work as heavy batch tasks adds unacceptable overhead.
\item Allowing incompatible document schemas across components creates runtime surprises.
\item Registering overly broad plugins grants the model unnecessary authority.
\item Persona descriptions cannot substitute for tool-level access control.
\end{enumerate}
\noindent\textbf{Answer.} Persona descriptions cannot substitute for tool-level access control.
\par\textit{CrewAI organizes role-based agents, tasks, crews, and flows for sequential or coordinated work. Tradeoff: It gives accessible orchestration but still needs explicit state, permissions, and observability. Watch for: Persona descriptions cannot substitute for tool-level access control.}
\paragraph{FC-0718 [medium]} Define CrewAI in interview-ready language.
\noindent\textbf{Answer.} CrewAI organizes role-based agents, tasks, crews, and flows for sequential or coordinated work.
\par\textit{CrewAI organizes role-based agents, tasks, crews, and flows for sequential or coordinated work.}
\paragraph{FC-0719 [hard]} State the main tradeoff and production pitfall for CrewAI.
\noindent\textbf{Answer.} Tradeoff: It gives accessible orchestration but still needs explicit state, permissions, and observability. Pitfall: Persona descriptions cannot substitute for tool-level access control.
\par\textit{Tradeoff: It gives accessible orchestration but still needs explicit state, permissions, and observability. Pitfall: Persona descriptions cannot substitute for tool-level access control.}
\paragraph{LA-0720 [hard]} Explain CrewAI from first principles, then show how you would evaluate and operate it in production.
\noindent\textbf{Answer.} Start with the mechanism: CrewAI organizes role-based agents, tasks, crews, and flows for sequential or coordinated work. Make the tradeoff explicit: It gives accessible orchestration but still needs explicit state, permissions, and observability. Define workload-specific offline metrics and a latency/cost budget, test important slices, instrument the critical path, and create a rollback criterion. The production review must explicitly guard against this failure: Persona descriptions cannot substitute for tool-level access control.
\par\textit{A strong answer connects mechanism, measurable tradeoffs, failure isolation, observability, and rollback rather than listing product features.}
\paragraph{MCQ-0721 [medium]} Which statement best describes DSPy?
\begin{enumerate}[label=\Alph*.]
\item DSPy treats prompts and modules as optimizable programs evaluated against examples and metrics.
\item LangGraph models long-running agent workflows as stateful graphs with checkpoints, interrupts, streaming, and durable execution.
\item Haystack connects typed components into retrieval, generation, routing, and agent pipelines.
\item AutoGen structures multi-agent conversations and tool-using participants for collaborative task execution.
\end{enumerate}
\noindent\textbf{Answer.} DSPy treats prompts and modules as optimizable programs evaluated against examples and metrics.
\par\textit{DSPy treats prompts and modules as optimizable programs evaluated against examples and metrics. Tradeoff: Systematic optimization can beat manual prompting, while it depends on representative data and stable metrics. Watch for: Optimizing against a leaky judge overfits wording instead of capability.}
\paragraph{MCQ-0722 [hard]} What is the central engineering tradeoff of DSPy?
\begin{enumerate}[label=\Alph*.]
\item Role decomposition enables complex workflows but multiplies calls, coordination failures, and evaluation cost.
\item Explicit components support testing, while custom graphs still require lifecycle discipline.
\item Fine-grained control aids production reliability but exposes more state and graph design decisions.
\item Systematic optimization can beat manual prompting, while it depends on representative data and stable metrics.
\end{enumerate}
\noindent\textbf{Answer.} Systematic optimization can beat manual prompting, while it depends on representative data and stable metrics.
\par\textit{DSPy treats prompts and modules as optimizable programs evaluated against examples and metrics. Tradeoff: Systematic optimization can beat manual prompting, while it depends on representative data and stable metrics. Watch for: Optimizing against a leaky judge overfits wording instead of capability.}
\paragraph{MCQ-0723 [hard]} Which failure mode should an interviewer expect you to identify for DSPy?
\begin{enumerate}[label=\Alph*.]
\item Optimizing against a leaky judge overfits wording instead of capability.
\item Non-idempotent node side effects make checkpoint replay unsafe.
\item Allowing incompatible document schemas across components creates runtime surprises.
\item Adding agents without independent capabilities creates expensive echo chambers.
\end{enumerate}
\noindent\textbf{Answer.} Optimizing against a leaky judge overfits wording instead of capability.
\par\textit{DSPy treats prompts and modules as optimizable programs evaluated against examples and metrics. Tradeoff: Systematic optimization can beat manual prompting, while it depends on representative data and stable metrics. Watch for: Optimizing against a leaky judge overfits wording instead of capability.}
\paragraph{MCQ-0724 [medium]} A design review is specifically concerned with this risk: Optimizing against a leaky judge overfits wording instead of capability. Which topic should be reviewed first?
\begin{enumerate}[label=\Alph*.]
\item LangGraph
\item DSPy
\item AutoGen
\item Haystack pipelines
\end{enumerate}
\noindent\textbf{Answer.} DSPy
\par\textit{DSPy treats prompts and modules as optimizable programs evaluated against examples and metrics. Tradeoff: Systematic optimization can beat manual prompting, while it depends on representative data and stable metrics. Watch for: Optimizing against a leaky judge overfits wording instead of capability.}
\paragraph{MCQ-0725 [hard]} Which design note most accurately balances the benefits and costs of DSPy?
\begin{enumerate}[label=\Alph*.]
\item Systematic optimization can beat manual prompting, while it depends on representative data and stable metrics.
\item Role decomposition enables complex workflows but multiplies calls, coordination failures, and evaluation cost.
\item Explicit components support testing, while custom graphs still require lifecycle discipline.
\item Fine-grained control aids production reliability but exposes more state and graph design decisions.
\end{enumerate}
\noindent\textbf{Answer.} Systematic optimization can beat manual prompting, while it depends on representative data and stable metrics.
\par\textit{DSPy treats prompts and modules as optimizable programs evaluated against examples and metrics. Tradeoff: Systematic optimization can beat manual prompting, while it depends on representative data and stable metrics. Watch for: Optimizing against a leaky judge overfits wording instead of capability.}
\paragraph{MCQ-0726 [medium]} Which explanation would be strongest in a system-design interview about DSPy?
\begin{enumerate}[label=\Alph*.]
\item DSPy treats prompts and modules as optimizable programs evaluated against examples and metrics.
\item AutoGen structures multi-agent conversations and tool-using participants for collaborative task execution.
\item LangGraph models long-running agent workflows as stateful graphs with checkpoints, interrupts, streaming, and durable execution.
\item Haystack connects typed components into retrieval, generation, routing, and agent pipelines.
\end{enumerate}
\noindent\textbf{Answer.} DSPy treats prompts and modules as optimizable programs evaluated against examples and metrics.
\par\textit{DSPy treats prompts and modules as optimizable programs evaluated against examples and metrics. Tradeoff: Systematic optimization can beat manual prompting, while it depends on representative data and stable metrics. Watch for: Optimizing against a leaky judge overfits wording instead of capability.}
\paragraph{MCQ-0727 [hard]} Which production warning is most directly associated with DSPy?
\begin{enumerate}[label=\Alph*.]
\item Non-idempotent node side effects make checkpoint replay unsafe.
\item Adding agents without independent capabilities creates expensive echo chambers.
\item Allowing incompatible document schemas across components creates runtime surprises.
\item Optimizing against a leaky judge overfits wording instead of capability.
\end{enumerate}
\noindent\textbf{Answer.} Optimizing against a leaky judge overfits wording instead of capability.
\par\textit{DSPy treats prompts and modules as optimizable programs evaluated against examples and metrics. Tradeoff: Systematic optimization can beat manual prompting, while it depends on representative data and stable metrics. Watch for: Optimizing against a leaky judge overfits wording instead of capability.}
\paragraph{FC-0728 [medium]} Define DSPy in interview-ready language.
\noindent\textbf{Answer.} DSPy treats prompts and modules as optimizable programs evaluated against examples and metrics.
\par\textit{DSPy treats prompts and modules as optimizable programs evaluated against examples and metrics.}
\paragraph{FC-0729 [hard]} State the main tradeoff and production pitfall for DSPy.
\noindent\textbf{Answer.} Tradeoff: Systematic optimization can beat manual prompting, while it depends on representative data and stable metrics. Pitfall: Optimizing against a leaky judge overfits wording instead of capability.
\par\textit{Tradeoff: Systematic optimization can beat manual prompting, while it depends on representative data and stable metrics. Pitfall: Optimizing against a leaky judge overfits wording instead of capability.}
\paragraph{LA-0730 [hard]} Explain DSPy from first principles, then show how you would evaluate and operate it in production.
\noindent\textbf{Answer.} Start with the mechanism: DSPy treats prompts and modules as optimizable programs evaluated against examples and metrics. Make the tradeoff explicit: Systematic optimization can beat manual prompting, while it depends on representative data and stable metrics. Define workload-specific offline metrics and a latency/cost budget, test important slices, instrument the critical path, and create a rollback criterion. The production review must explicitly guard against this failure: Optimizing against a leaky judge overfits wording instead of capability.
\par\textit{A strong answer connects mechanism, measurable tradeoffs, failure isolation, observability, and rollback rather than listing product features.}
\paragraph{MCQ-0731 [medium]} Which statement best describes Temporal?
\begin{enumerate}[label=\Alph*.]
\item Temporal persists workflow history and replays deterministic workflow code around durable activities.
\item General-purpose data orchestrators schedule DAGs, assets, or flows with retries, lineage, and operational UIs.
\item PydanticAI uses typed dependencies, tools, results, and validation to build model-agnostic Python agents.
\item AutoGen structures multi-agent conversations and tool-using participants for collaborative task execution.
\end{enumerate}
\noindent\textbf{Answer.} Temporal persists workflow history and replays deterministic workflow code around durable activities.
\par\textit{Temporal persists workflow history and replays deterministic workflow code around durable activities. Tradeoff: It provides robust retries and long-running execution but requires deterministic workflow discipline. Watch for: Reading wall-clock time or random state directly in replayed code causes nondeterminism.}
\paragraph{MCQ-0732 [hard]} What is the central engineering tradeoff of Temporal?
\begin{enumerate}[label=\Alph*.]
\item Type safety improves integration but cannot validate semantic truth.
\item Role decomposition enables complex workflows but multiplies calls, coordination failures, and evaluation cost.
\item It provides robust retries and long-running execution but requires deterministic workflow discipline.
\item They excel at batch pipelines but are not automatically low-latency interactive agent runtimes.
\end{enumerate}
\noindent\textbf{Answer.} It provides robust retries and long-running execution but requires deterministic workflow discipline.
\par\textit{Temporal persists workflow history and replays deterministic workflow code around durable activities. Tradeoff: It provides robust retries and long-running execution but requires deterministic workflow discipline. Watch for: Reading wall-clock time or random state directly in replayed code causes nondeterminism.}
\paragraph{MCQ-0733 [hard]} Which failure mode should an interviewer expect you to identify for Temporal?
\begin{enumerate}[label=\Alph*.]
\item Reading wall-clock time or random state directly in replayed code causes nondeterminism.
\item Adding agents without independent capabilities creates expensive echo chambers.
\item Confusing schema conformance with a correct answer produces false confidence.
\item Running per-token or per-chat-turn work as heavy batch tasks adds unacceptable overhead.
\end{enumerate}
\noindent\textbf{Answer.} Reading wall-clock time or random state directly in replayed code causes nondeterminism.
\par\textit{Temporal persists workflow history and replays deterministic workflow code around durable activities. Tradeoff: It provides robust retries and long-running execution but requires deterministic workflow discipline. Watch for: Reading wall-clock time or random state directly in replayed code causes nondeterminism.}
\paragraph{MCQ-0734 [medium]} A design review is specifically concerned with this risk: Reading wall-clock time or random state directly in replayed code causes nondeterminism. Which topic should be reviewed first?
\begin{enumerate}[label=\Alph*.]
\item PydanticAI
\item AutoGen
\item Airflow, Dagster, and Prefect
\item Temporal
\end{enumerate}
\noindent\textbf{Answer.} Temporal
\par\textit{Temporal persists workflow history and replays deterministic workflow code around durable activities. Tradeoff: It provides robust retries and long-running execution but requires deterministic workflow discipline. Watch for: Reading wall-clock time or random state directly in replayed code causes nondeterminism.}
\paragraph{MCQ-0735 [hard]} Which design note most accurately balances the benefits and costs of Temporal?
\begin{enumerate}[label=\Alph*.]
\item Role decomposition enables complex workflows but multiplies calls, coordination failures, and evaluation cost.
\item Type safety improves integration but cannot validate semantic truth.
\item It provides robust retries and long-running execution but requires deterministic workflow discipline.
\item They excel at batch pipelines but are not automatically low-latency interactive agent runtimes.
\end{enumerate}
\noindent\textbf{Answer.} It provides robust retries and long-running execution but requires deterministic workflow discipline.
\par\textit{Temporal persists workflow history and replays deterministic workflow code around durable activities. Tradeoff: It provides robust retries and long-running execution but requires deterministic workflow discipline. Watch for: Reading wall-clock time or random state directly in replayed code causes nondeterminism.}
\paragraph{MCQ-0736 [medium]} Which explanation would be strongest in a system-design interview about Temporal?
\begin{enumerate}[label=\Alph*.]
\item PydanticAI uses typed dependencies, tools, results, and validation to build model-agnostic Python agents.
\item AutoGen structures multi-agent conversations and tool-using participants for collaborative task execution.
\item Temporal persists workflow history and replays deterministic workflow code around durable activities.
\item General-purpose data orchestrators schedule DAGs, assets, or flows with retries, lineage, and operational UIs.
\end{enumerate}
\noindent\textbf{Answer.} Temporal persists workflow history and replays deterministic workflow code around durable activities.
\par\textit{Temporal persists workflow history and replays deterministic workflow code around durable activities. Tradeoff: It provides robust retries and long-running execution but requires deterministic workflow discipline. Watch for: Reading wall-clock time or random state directly in replayed code causes nondeterminism.}
\paragraph{MCQ-0737 [hard]} Which production warning is most directly associated with Temporal?
\begin{enumerate}[label=\Alph*.]
\item Adding agents without independent capabilities creates expensive echo chambers.
\item Running per-token or per-chat-turn work as heavy batch tasks adds unacceptable overhead.
\item Reading wall-clock time or random state directly in replayed code causes nondeterminism.
\item Confusing schema conformance with a correct answer produces false confidence.
\end{enumerate}
\noindent\textbf{Answer.} Reading wall-clock time or random state directly in replayed code causes nondeterminism.
\par\textit{Temporal persists workflow history and replays deterministic workflow code around durable activities. Tradeoff: It provides robust retries and long-running execution but requires deterministic workflow discipline. Watch for: Reading wall-clock time or random state directly in replayed code causes nondeterminism.}
\paragraph{FC-0738 [medium]} Define Temporal in interview-ready language.
\noindent\textbf{Answer.} Temporal persists workflow history and replays deterministic workflow code around durable activities.
\par\textit{Temporal persists workflow history and replays deterministic workflow code around durable activities.}
\paragraph{FC-0739 [hard]} State the main tradeoff and production pitfall for Temporal.
\noindent\textbf{Answer.} Tradeoff: It provides robust retries and long-running execution but requires deterministic workflow discipline. Pitfall: Reading wall-clock time or random state directly in replayed code causes nondeterminism.
\par\textit{Tradeoff: It provides robust retries and long-running execution but requires deterministic workflow discipline. Pitfall: Reading wall-clock time or random state directly in replayed code causes nondeterminism.}
\paragraph{LA-0740 [hard]} Explain Temporal from first principles, then show how you would evaluate and operate it in production.
\noindent\textbf{Answer.} Start with the mechanism: Temporal persists workflow history and replays deterministic workflow code around durable activities. Make the tradeoff explicit: It provides robust retries and long-running execution but requires deterministic workflow discipline. Define workload-specific offline metrics and a latency/cost budget, test important slices, instrument the critical path, and create a rollback criterion. The production review must explicitly guard against this failure: Reading wall-clock time or random state directly in replayed code causes nondeterminism.
\par\textit{A strong answer connects mechanism, measurable tradeoffs, failure isolation, observability, and rollback rather than listing product features.}
\paragraph{MCQ-0741 [medium]} Which statement best describes Airflow, Dagster, and Prefect?
\begin{enumerate}[label=\Alph*.]
\item General-purpose data orchestrators schedule DAGs, assets, or flows with retries, lineage, and operational UIs.
\item Haystack connects typed components into retrieval, generation, routing, and agent pipelines.
\item LangGraph models long-running agent workflows as stateful graphs with checkpoints, interrupts, streaming, and durable execution.
\item DSPy treats prompts and modules as optimizable programs evaluated against examples and metrics.
\end{enumerate}
\noindent\textbf{Answer.} General-purpose data orchestrators schedule DAGs, assets, or flows with retries, lineage, and operational UIs.
\par\textit{General-purpose data orchestrators schedule DAGs, assets, or flows with retries, lineage, and operational UIs. Tradeoff: They excel at batch pipelines but are not automatically low-latency interactive agent runtimes. Watch for: Running per-token or per-chat-turn work as heavy batch tasks adds unacceptable overhead.}
\paragraph{MCQ-0742 [hard]} What is the central engineering tradeoff of Airflow, Dagster, and Prefect?
\begin{enumerate}[label=\Alph*.]
\item Fine-grained control aids production reliability but exposes more state and graph design decisions.
\item Explicit components support testing, while custom graphs still require lifecycle discipline.
\item Systematic optimization can beat manual prompting, while it depends on representative data and stable metrics.
\item They excel at batch pipelines but are not automatically low-latency interactive agent runtimes.
\end{enumerate}
\noindent\textbf{Answer.} They excel at batch pipelines but are not automatically low-latency interactive agent runtimes.
\par\textit{General-purpose data orchestrators schedule DAGs, assets, or flows with retries, lineage, and operational UIs. Tradeoff: They excel at batch pipelines but are not automatically low-latency interactive agent runtimes. Watch for: Running per-token or per-chat-turn work as heavy batch tasks adds unacceptable overhead.}
\paragraph{MCQ-0743 [hard]} Which failure mode should an interviewer expect you to identify for Airflow, Dagster, and Prefect?
\begin{enumerate}[label=\Alph*.]
\item Allowing incompatible document schemas across components creates runtime surprises.
\item Non-idempotent node side effects make checkpoint replay unsafe.
\item Running per-token or per-chat-turn work as heavy batch tasks adds unacceptable overhead.
\item Optimizing against a leaky judge overfits wording instead of capability.
\end{enumerate}
\noindent\textbf{Answer.} Running per-token or per-chat-turn work as heavy batch tasks adds unacceptable overhead.
\par\textit{General-purpose data orchestrators schedule DAGs, assets, or flows with retries, lineage, and operational UIs. Tradeoff: They excel at batch pipelines but are not automatically low-latency interactive agent runtimes. Watch for: Running per-token or per-chat-turn work as heavy batch tasks adds unacceptable overhead.}
\paragraph{MCQ-0744 [medium]} A design review is specifically concerned with this risk: Running per-token or per-chat-turn work as heavy batch tasks adds unacceptable overhead. Which topic should be reviewed first?
\begin{enumerate}[label=\Alph*.]
\item Haystack pipelines
\item Airflow, Dagster, and Prefect
\item LangGraph
\item DSPy
\end{enumerate}
\noindent\textbf{Answer.} Airflow, Dagster, and Prefect
\par\textit{General-purpose data orchestrators schedule DAGs, assets, or flows with retries, lineage, and operational UIs. Tradeoff: They excel at batch pipelines but are not automatically low-latency interactive agent runtimes. Watch for: Running per-token or per-chat-turn work as heavy batch tasks adds unacceptable overhead.}
\paragraph{MCQ-0745 [hard]} Which design note most accurately balances the benefits and costs of Airflow, Dagster, and Prefect?
\begin{enumerate}[label=\Alph*.]
\item Fine-grained control aids production reliability but exposes more state and graph design decisions.
\item Systematic optimization can beat manual prompting, while it depends on representative data and stable metrics.
\item Explicit components support testing, while custom graphs still require lifecycle discipline.
\item They excel at batch pipelines but are not automatically low-latency interactive agent runtimes.
\end{enumerate}
\noindent\textbf{Answer.} They excel at batch pipelines but are not automatically low-latency interactive agent runtimes.
\par\textit{General-purpose data orchestrators schedule DAGs, assets, or flows with retries, lineage, and operational UIs. Tradeoff: They excel at batch pipelines but are not automatically low-latency interactive agent runtimes. Watch for: Running per-token or per-chat-turn work as heavy batch tasks adds unacceptable overhead.}
\paragraph{MCQ-0746 [medium]} Which explanation would be strongest in a system-design interview about Airflow, Dagster, and Prefect?
\begin{enumerate}[label=\Alph*.]
\item DSPy treats prompts and modules as optimizable programs evaluated against examples and metrics.
\item General-purpose data orchestrators schedule DAGs, assets, or flows with retries, lineage, and operational UIs.
\item Haystack connects typed components into retrieval, generation, routing, and agent pipelines.
\item LangGraph models long-running agent workflows as stateful graphs with checkpoints, interrupts, streaming, and durable execution.
\end{enumerate}
\noindent\textbf{Answer.} General-purpose data orchestrators schedule DAGs, assets, or flows with retries, lineage, and operational UIs.
\par\textit{General-purpose data orchestrators schedule DAGs, assets, or flows with retries, lineage, and operational UIs. Tradeoff: They excel at batch pipelines but are not automatically low-latency interactive agent runtimes. Watch for: Running per-token or per-chat-turn work as heavy batch tasks adds unacceptable overhead.}
\paragraph{MCQ-0747 [hard]} Which production warning is most directly associated with Airflow, Dagster, and Prefect?
\begin{enumerate}[label=\Alph*.]
\item Running per-token or per-chat-turn work as heavy batch tasks adds unacceptable overhead.
\item Non-idempotent node side effects make checkpoint replay unsafe.
\item Allowing incompatible document schemas across components creates runtime surprises.
\item Optimizing against a leaky judge overfits wording instead of capability.
\end{enumerate}
\noindent\textbf{Answer.} Running per-token or per-chat-turn work as heavy batch tasks adds unacceptable overhead.
\par\textit{General-purpose data orchestrators schedule DAGs, assets, or flows with retries, lineage, and operational UIs. Tradeoff: They excel at batch pipelines but are not automatically low-latency interactive agent runtimes. Watch for: Running per-token or per-chat-turn work as heavy batch tasks adds unacceptable overhead.}
\paragraph{FC-0748 [medium]} Define Airflow, Dagster, and Prefect in interview-ready language.
\noindent\textbf{Answer.} General-purpose data orchestrators schedule DAGs, assets, or flows with retries, lineage, and operational UIs.
\par\textit{General-purpose data orchestrators schedule DAGs, assets, or flows with retries, lineage, and operational UIs.}
\paragraph{FC-0749 [hard]} State the main tradeoff and production pitfall for Airflow, Dagster, and Prefect.
\noindent\textbf{Answer.} Tradeoff: They excel at batch pipelines but are not automatically low-latency interactive agent runtimes. Pitfall: Running per-token or per-chat-turn work as heavy batch tasks adds unacceptable overhead.
\par\textit{Tradeoff: They excel at batch pipelines but are not automatically low-latency interactive agent runtimes. Pitfall: Running per-token or per-chat-turn work as heavy batch tasks adds unacceptable overhead.}
\paragraph{LA-0750 [hard]} Explain Airflow, Dagster, and Prefect from first principles, then show how you would evaluate and operate it in production.
\noindent\textbf{Answer.} Start with the mechanism: General-purpose data orchestrators schedule DAGs, assets, or flows with retries, lineage, and operational UIs. Make the tradeoff explicit: They excel at batch pipelines but are not automatically low-latency interactive agent runtimes. Define workload-specific offline metrics and a latency/cost budget, test important slices, instrument the critical path, and create a rollback criterion. The production review must explicitly guard against this failure: Running per-token or per-chat-turn work as heavy batch tasks adds unacceptable overhead.
\par\textit{A strong answer connects mechanism, measurable tradeoffs, failure isolation, observability, and rollback rather than listing product features.}
\subsection{Observability and Monitoring}
\paragraph{MCQ-0751 [medium]} Which statement best describes OpenTelemetry?
\begin{enumerate}[label=\Alph*.]
\item Weave tracks calls, objects, datasets, evaluations, and model application experiments.
\item MLflow connects traces, evaluation datasets, scorers, experiments, and production monitoring.
\item Counters, histograms, gauges, service-level indicators, and objectives quantify reliability and performance.
\item OpenTelemetry provides vendor-neutral APIs, SDKs, collectors, and semantic conventions for traces, metrics, logs, and profiles.
\end{enumerate}
\noindent\textbf{Answer.} OpenTelemetry provides vendor-neutral APIs, SDKs, collectors, and semantic conventions for traces, metrics, logs, and profiles.
\par\textit{OpenTelemetry provides vendor-neutral APIs, SDKs, collectors, and semantic conventions for traces, metrics, logs, and profiles. Tradeoff: Portability improves, while instrumentation quality and backend cost remain engineering work. Watch for: High-cardinality prompt or user attributes can overwhelm telemetry systems.}
\paragraph{MCQ-0752 [hard]} What is the central engineering tradeoff of OpenTelemetry?
\begin{enumerate}[label=\Alph*.]
\item Portability improves, while instrumentation quality and backend cost remain engineering work.
\item A common ML platform improves lineage but may be heavier than a focused tracing tool.
\item Rich experiment lineage helps collaboration but introduces another governance surface.
\item Aggregate signals scale well but cannot explain individual semantic failures.
\end{enumerate}
\noindent\textbf{Answer.} Portability improves, while instrumentation quality and backend cost remain engineering work.
\par\textit{OpenTelemetry provides vendor-neutral APIs, SDKs, collectors, and semantic conventions for traces, metrics, logs, and profiles. Tradeoff: Portability improves, while instrumentation quality and backend cost remain engineering work. Watch for: High-cardinality prompt or user attributes can overwhelm telemetry systems.}
\paragraph{MCQ-0753 [hard]} Which failure mode should an interviewer expect you to identify for OpenTelemetry?
\begin{enumerate}[label=\Alph*.]
\item Saving mutable datasets without immutable versions undermines reproducibility.
\item High-cardinality prompt or user attributes can overwhelm telemetry systems.
\item Mixing incomparable scorer versions in one trend chart produces false regressions.
\item Averaging latency hides long-tail p95 and p99 user pain.
\end{enumerate}
\noindent\textbf{Answer.} High-cardinality prompt or user attributes can overwhelm telemetry systems.
\par\textit{OpenTelemetry provides vendor-neutral APIs, SDKs, collectors, and semantic conventions for traces, metrics, logs, and profiles. Tradeoff: Portability improves, while instrumentation quality and backend cost remain engineering work. Watch for: High-cardinality prompt or user attributes can overwhelm telemetry systems.}
\paragraph{MCQ-0754 [medium]} A design review is specifically concerned with this risk: High-cardinality prompt or user attributes can overwhelm telemetry systems. Which topic should be reviewed first?
\begin{enumerate}[label=\Alph*.]
\item W\&B Weave
\item OpenTelemetry
\item Metrics and SLOs
\item MLflow tracing
\end{enumerate}
\noindent\textbf{Answer.} OpenTelemetry
\par\textit{OpenTelemetry provides vendor-neutral APIs, SDKs, collectors, and semantic conventions for traces, metrics, logs, and profiles. Tradeoff: Portability improves, while instrumentation quality and backend cost remain engineering work. Watch for: High-cardinality prompt or user attributes can overwhelm telemetry systems.}
\paragraph{MCQ-0755 [hard]} Which design note most accurately balances the benefits and costs of OpenTelemetry?
\begin{enumerate}[label=\Alph*.]
\item Portability improves, while instrumentation quality and backend cost remain engineering work.
\item Rich experiment lineage helps collaboration but introduces another governance surface.
\item Aggregate signals scale well but cannot explain individual semantic failures.
\item A common ML platform improves lineage but may be heavier than a focused tracing tool.
\end{enumerate}
\noindent\textbf{Answer.} Portability improves, while instrumentation quality and backend cost remain engineering work.
\par\textit{OpenTelemetry provides vendor-neutral APIs, SDKs, collectors, and semantic conventions for traces, metrics, logs, and profiles. Tradeoff: Portability improves, while instrumentation quality and backend cost remain engineering work. Watch for: High-cardinality prompt or user attributes can overwhelm telemetry systems.}
\paragraph{MCQ-0756 [medium]} Which explanation would be strongest in a system-design interview about OpenTelemetry?
\begin{enumerate}[label=\Alph*.]
\item MLflow connects traces, evaluation datasets, scorers, experiments, and production monitoring.
\item OpenTelemetry provides vendor-neutral APIs, SDKs, collectors, and semantic conventions for traces, metrics, logs, and profiles.
\item Weave tracks calls, objects, datasets, evaluations, and model application experiments.
\item Counters, histograms, gauges, service-level indicators, and objectives quantify reliability and performance.
\end{enumerate}
\noindent\textbf{Answer.} OpenTelemetry provides vendor-neutral APIs, SDKs, collectors, and semantic conventions for traces, metrics, logs, and profiles.
\par\textit{OpenTelemetry provides vendor-neutral APIs, SDKs, collectors, and semantic conventions for traces, metrics, logs, and profiles. Tradeoff: Portability improves, while instrumentation quality and backend cost remain engineering work. Watch for: High-cardinality prompt or user attributes can overwhelm telemetry systems.}
\paragraph{MCQ-0757 [hard]} Which production warning is most directly associated with OpenTelemetry?
\begin{enumerate}[label=\Alph*.]
\item Saving mutable datasets without immutable versions undermines reproducibility.
\item High-cardinality prompt or user attributes can overwhelm telemetry systems.
\item Mixing incomparable scorer versions in one trend chart produces false regressions.
\item Averaging latency hides long-tail p95 and p99 user pain.
\end{enumerate}
\noindent\textbf{Answer.} High-cardinality prompt or user attributes can overwhelm telemetry systems.
\par\textit{OpenTelemetry provides vendor-neutral APIs, SDKs, collectors, and semantic conventions for traces, metrics, logs, and profiles. Tradeoff: Portability improves, while instrumentation quality and backend cost remain engineering work. Watch for: High-cardinality prompt or user attributes can overwhelm telemetry systems.}
\paragraph{FC-0758 [medium]} Define OpenTelemetry in interview-ready language.
\noindent\textbf{Answer.} OpenTelemetry provides vendor-neutral APIs, SDKs, collectors, and semantic conventions for traces, metrics, logs, and profiles.
\par\textit{OpenTelemetry provides vendor-neutral APIs, SDKs, collectors, and semantic conventions for traces, metrics, logs, and profiles.}
\paragraph{FC-0759 [hard]} State the main tradeoff and production pitfall for OpenTelemetry.
\noindent\textbf{Answer.} Tradeoff: Portability improves, while instrumentation quality and backend cost remain engineering work. Pitfall: High-cardinality prompt or user attributes can overwhelm telemetry systems.
\par\textit{Tradeoff: Portability improves, while instrumentation quality and backend cost remain engineering work. Pitfall: High-cardinality prompt or user attributes can overwhelm telemetry systems.}
\paragraph{LA-0760 [hard]} Explain OpenTelemetry from first principles, then show how you would evaluate and operate it in production.
\noindent\textbf{Answer.} Start with the mechanism: OpenTelemetry provides vendor-neutral APIs, SDKs, collectors, and semantic conventions for traces, metrics, logs, and profiles. Make the tradeoff explicit: Portability improves, while instrumentation quality and backend cost remain engineering work. Define workload-specific offline metrics and a latency/cost budget, test important slices, instrument the critical path, and create a rollback criterion. The production review must explicitly guard against this failure: High-cardinality prompt or user attributes can overwhelm telemetry systems.
\par\textit{A strong answer connects mechanism, measurable tradeoffs, failure isolation, observability, and rollback rather than listing product features.}
\paragraph{MCQ-0761 [medium]} Which statement best describes Distributed tracing?
\begin{enumerate}[label=\Alph*.]
\item LLM traces capture prompts, outputs, model settings, token usage, latency, tool calls, retrieval, and evaluations.
\item Counters, histograms, gauges, service-level indicators, and objectives quantify reliability and performance.
\item Traces connect request spans across gateways, retrievers, models, tools, queues, and databases.
\item Weave tracks calls, objects, datasets, evaluations, and model application experiments.
\end{enumerate}
\noindent\textbf{Answer.} Traces connect request spans across gateways, retrievers, models, tools, queues, and databases.
\par\textit{Traces connect request spans across gateways, retrievers, models, tools, queues, and databases. Tradeoff: They reveal critical paths but sampling can hide rare quality failures. Watch for: Reusing trace identifiers across independent user requests breaks isolation.}
\paragraph{MCQ-0762 [hard]} What is the central engineering tradeoff of Distributed tracing?
\begin{enumerate}[label=\Alph*.]
\item Debugging improves, while content capture raises privacy and retention risk.
\item They reveal critical paths but sampling can hide rare quality failures.
\item Aggregate signals scale well but cannot explain individual semantic failures.
\item Rich experiment lineage helps collaboration but introduces another governance surface.
\end{enumerate}
\noindent\textbf{Answer.} They reveal critical paths but sampling can hide rare quality failures.
\par\textit{Traces connect request spans across gateways, retrievers, models, tools, queues, and databases. Tradeoff: They reveal critical paths but sampling can hide rare quality failures. Watch for: Reusing trace identifiers across independent user requests breaks isolation.}
\paragraph{MCQ-0763 [hard]} Which failure mode should an interviewer expect you to identify for Distributed tracing?
\begin{enumerate}[label=\Alph*.]
\item Reusing trace identifiers across independent user requests breaks isolation.
\item Logging secrets or personal data by default creates a second sensitive datastore.
\item Averaging latency hides long-tail p95 and p99 user pain.
\item Saving mutable datasets without immutable versions undermines reproducibility.
\end{enumerate}
\noindent\textbf{Answer.} Reusing trace identifiers across independent user requests breaks isolation.
\par\textit{Traces connect request spans across gateways, retrievers, models, tools, queues, and databases. Tradeoff: They reveal critical paths but sampling can hide rare quality failures. Watch for: Reusing trace identifiers across independent user requests breaks isolation.}
\paragraph{MCQ-0764 [medium]} A design review is specifically concerned with this risk: Reusing trace identifiers across independent user requests breaks isolation. Which topic should be reviewed first?
\begin{enumerate}[label=\Alph*.]
\item W\&B Weave
\item Metrics and SLOs
\item LLM tracing
\item Distributed tracing
\end{enumerate}
\noindent\textbf{Answer.} Distributed tracing
\par\textit{Traces connect request spans across gateways, retrievers, models, tools, queues, and databases. Tradeoff: They reveal critical paths but sampling can hide rare quality failures. Watch for: Reusing trace identifiers across independent user requests breaks isolation.}
\paragraph{MCQ-0765 [hard]} Which design note most accurately balances the benefits and costs of Distributed tracing?
\begin{enumerate}[label=\Alph*.]
\item Debugging improves, while content capture raises privacy and retention risk.
\item Aggregate signals scale well but cannot explain individual semantic failures.
\item They reveal critical paths but sampling can hide rare quality failures.
\item Rich experiment lineage helps collaboration but introduces another governance surface.
\end{enumerate}
\noindent\textbf{Answer.} They reveal critical paths but sampling can hide rare quality failures.
\par\textit{Traces connect request spans across gateways, retrievers, models, tools, queues, and databases. Tradeoff: They reveal critical paths but sampling can hide rare quality failures. Watch for: Reusing trace identifiers across independent user requests breaks isolation.}
\paragraph{MCQ-0766 [medium]} Which explanation would be strongest in a system-design interview about Distributed tracing?
\begin{enumerate}[label=\Alph*.]
\item Counters, histograms, gauges, service-level indicators, and objectives quantify reliability and performance.
\item LLM traces capture prompts, outputs, model settings, token usage, latency, tool calls, retrieval, and evaluations.
\item Weave tracks calls, objects, datasets, evaluations, and model application experiments.
\item Traces connect request spans across gateways, retrievers, models, tools, queues, and databases.
\end{enumerate}
\noindent\textbf{Answer.} Traces connect request spans across gateways, retrievers, models, tools, queues, and databases.
\par\textit{Traces connect request spans across gateways, retrievers, models, tools, queues, and databases. Tradeoff: They reveal critical paths but sampling can hide rare quality failures. Watch for: Reusing trace identifiers across independent user requests breaks isolation.}
\paragraph{MCQ-0767 [hard]} Which production warning is most directly associated with Distributed tracing?
\begin{enumerate}[label=\Alph*.]
\item Saving mutable datasets without immutable versions undermines reproducibility.
\item Logging secrets or personal data by default creates a second sensitive datastore.
\item Reusing trace identifiers across independent user requests breaks isolation.
\item Averaging latency hides long-tail p95 and p99 user pain.
\end{enumerate}
\noindent\textbf{Answer.} Reusing trace identifiers across independent user requests breaks isolation.
\par\textit{Traces connect request spans across gateways, retrievers, models, tools, queues, and databases. Tradeoff: They reveal critical paths but sampling can hide rare quality failures. Watch for: Reusing trace identifiers across independent user requests breaks isolation.}
\paragraph{FC-0768 [medium]} Define Distributed tracing in interview-ready language.
\noindent\textbf{Answer.} Traces connect request spans across gateways, retrievers, models, tools, queues, and databases.
\par\textit{Traces connect request spans across gateways, retrievers, models, tools, queues, and databases.}
\paragraph{FC-0769 [hard]} State the main tradeoff and production pitfall for Distributed tracing.
\noindent\textbf{Answer.} Tradeoff: They reveal critical paths but sampling can hide rare quality failures. Pitfall: Reusing trace identifiers across independent user requests breaks isolation.
\par\textit{Tradeoff: They reveal critical paths but sampling can hide rare quality failures. Pitfall: Reusing trace identifiers across independent user requests breaks isolation.}
\paragraph{LA-0770 [hard]} Explain Distributed tracing from first principles, then show how you would evaluate and operate it in production.
\noindent\textbf{Answer.} Start with the mechanism: Traces connect request spans across gateways, retrievers, models, tools, queues, and databases. Make the tradeoff explicit: They reveal critical paths but sampling can hide rare quality failures. Define workload-specific offline metrics and a latency/cost budget, test important slices, instrument the critical path, and create a rollback criterion. The production review must explicitly guard against this failure: Reusing trace identifiers across independent user requests breaks isolation.
\par\textit{A strong answer connects mechanism, measurable tradeoffs, failure isolation, observability, and rollback rather than listing product features.}
\paragraph{MCQ-0771 [medium]} Which statement best describes LLM tracing?
\begin{enumerate}[label=\Alph*.]
\item OpenTelemetry provides vendor-neutral APIs, SDKs, collectors, and semantic conventions for traces, metrics, logs, and profiles.
\item Drift monitoring detects changes in inputs, retrieved evidence, outputs, model versions, and user behavior.
\item Cost telemetry attributes input, output, cache, retrieval, tool, and infrastructure spend to requests and tenants.
\item LLM traces capture prompts, outputs, model settings, token usage, latency, tool calls, retrieval, and evaluations.
\end{enumerate}
\noindent\textbf{Answer.} LLM traces capture prompts, outputs, model settings, token usage, latency, tool calls, retrieval, and evaluations.
\par\textit{LLM traces capture prompts, outputs, model settings, token usage, latency, tool calls, retrieval, and evaluations. Tradeoff: Debugging improves, while content capture raises privacy and retention risk. Watch for: Logging secrets or personal data by default creates a second sensitive datastore.}
\paragraph{MCQ-0772 [hard]} What is the central engineering tradeoff of LLM tracing?
\begin{enumerate}[label=\Alph*.]
\item Debugging improves, while content capture raises privacy and retention risk.
\item Unit economics become actionable, while provider pricing and caching rules evolve.
\item Portability improves, while instrumentation quality and backend cost remain engineering work.
\item Alerts catch regressions but thresholds must distinguish seasonality from harm.
\end{enumerate}
\noindent\textbf{Answer.} Debugging improves, while content capture raises privacy and retention risk.
\par\textit{LLM traces capture prompts, outputs, model settings, token usage, latency, tool calls, retrieval, and evaluations. Tradeoff: Debugging improves, while content capture raises privacy and retention risk. Watch for: Logging secrets or personal data by default creates a second sensitive datastore.}
\paragraph{MCQ-0773 [hard]} Which failure mode should an interviewer expect you to identify for LLM tracing?
\begin{enumerate}[label=\Alph*.]
\item Tracking only average cost hides pathological long-running agents.
\item High-cardinality prompt or user attributes can overwhelm telemetry systems.
\item Measuring embedding distribution shift alone does not prove answer quality changed.
\item Logging secrets or personal data by default creates a second sensitive datastore.
\end{enumerate}
\noindent\textbf{Answer.} Logging secrets or personal data by default creates a second sensitive datastore.
\par\textit{LLM traces capture prompts, outputs, model settings, token usage, latency, tool calls, retrieval, and evaluations. Tradeoff: Debugging improves, while content capture raises privacy and retention risk. Watch for: Logging secrets or personal data by default creates a second sensitive datastore.}
\paragraph{MCQ-0774 [medium]} A design review is specifically concerned with this risk: Logging secrets or personal data by default creates a second sensitive datastore. Which topic should be reviewed first?
\begin{enumerate}[label=\Alph*.]
\item Cost observability
\item LLM tracing
\item Prompt and model drift
\item OpenTelemetry
\end{enumerate}
\noindent\textbf{Answer.} LLM tracing
\par\textit{LLM traces capture prompts, outputs, model settings, token usage, latency, tool calls, retrieval, and evaluations. Tradeoff: Debugging improves, while content capture raises privacy and retention risk. Watch for: Logging secrets or personal data by default creates a second sensitive datastore.}
\paragraph{MCQ-0775 [hard]} Which design note most accurately balances the benefits and costs of LLM tracing?
\begin{enumerate}[label=\Alph*.]
\item Debugging improves, while content capture raises privacy and retention risk.
\item Portability improves, while instrumentation quality and backend cost remain engineering work.
\item Unit economics become actionable, while provider pricing and caching rules evolve.
\item Alerts catch regressions but thresholds must distinguish seasonality from harm.
\end{enumerate}
\noindent\textbf{Answer.} Debugging improves, while content capture raises privacy and retention risk.
\par\textit{LLM traces capture prompts, outputs, model settings, token usage, latency, tool calls, retrieval, and evaluations. Tradeoff: Debugging improves, while content capture raises privacy and retention risk. Watch for: Logging secrets or personal data by default creates a second sensitive datastore.}
\paragraph{MCQ-0776 [medium]} Which explanation would be strongest in a system-design interview about LLM tracing?
\begin{enumerate}[label=\Alph*.]
\item Cost telemetry attributes input, output, cache, retrieval, tool, and infrastructure spend to requests and tenants.
\item Drift monitoring detects changes in inputs, retrieved evidence, outputs, model versions, and user behavior.
\item OpenTelemetry provides vendor-neutral APIs, SDKs, collectors, and semantic conventions for traces, metrics, logs, and profiles.
\item LLM traces capture prompts, outputs, model settings, token usage, latency, tool calls, retrieval, and evaluations.
\end{enumerate}
\noindent\textbf{Answer.} LLM traces capture prompts, outputs, model settings, token usage, latency, tool calls, retrieval, and evaluations.
\par\textit{LLM traces capture prompts, outputs, model settings, token usage, latency, tool calls, retrieval, and evaluations. Tradeoff: Debugging improves, while content capture raises privacy and retention risk. Watch for: Logging secrets or personal data by default creates a second sensitive datastore.}
\paragraph{MCQ-0777 [hard]} Which production warning is most directly associated with LLM tracing?
\begin{enumerate}[label=\Alph*.]
\item Logging secrets or personal data by default creates a second sensitive datastore.
\item High-cardinality prompt or user attributes can overwhelm telemetry systems.
\item Measuring embedding distribution shift alone does not prove answer quality changed.
\item Tracking only average cost hides pathological long-running agents.
\end{enumerate}
\noindent\textbf{Answer.} Logging secrets or personal data by default creates a second sensitive datastore.
\par\textit{LLM traces capture prompts, outputs, model settings, token usage, latency, tool calls, retrieval, and evaluations. Tradeoff: Debugging improves, while content capture raises privacy and retention risk. Watch for: Logging secrets or personal data by default creates a second sensitive datastore.}
\paragraph{FC-0778 [medium]} Define LLM tracing in interview-ready language.
\noindent\textbf{Answer.} LLM traces capture prompts, outputs, model settings, token usage, latency, tool calls, retrieval, and evaluations.
\par\textit{LLM traces capture prompts, outputs, model settings, token usage, latency, tool calls, retrieval, and evaluations.}
\paragraph{FC-0779 [hard]} State the main tradeoff and production pitfall for LLM tracing.
\noindent\textbf{Answer.} Tradeoff: Debugging improves, while content capture raises privacy and retention risk. Pitfall: Logging secrets or personal data by default creates a second sensitive datastore.
\par\textit{Tradeoff: Debugging improves, while content capture raises privacy and retention risk. Pitfall: Logging secrets or personal data by default creates a second sensitive datastore.}
\paragraph{LA-0780 [hard]} Explain LLM tracing from first principles, then show how you would evaluate and operate it in production.
\noindent\textbf{Answer.} Start with the mechanism: LLM traces capture prompts, outputs, model settings, token usage, latency, tool calls, retrieval, and evaluations. Make the tradeoff explicit: Debugging improves, while content capture raises privacy and retention risk. Define workload-specific offline metrics and a latency/cost budget, test important slices, instrument the critical path, and create a rollback criterion. The production review must explicitly guard against this failure: Logging secrets or personal data by default creates a second sensitive datastore.
\par\textit{A strong answer connects mechanism, measurable tradeoffs, failure isolation, observability, and rollback rather than listing product features.}
\paragraph{MCQ-0781 [medium]} Which statement best describes Metrics and SLOs?
\begin{enumerate}[label=\Alph*.]
\item MLflow connects traces, evaluation datasets, scorers, experiments, and production monitoring.
\item Counters, histograms, gauges, service-level indicators, and objectives quantify reliability and performance.
\item Structured events use stable fields, severity, timestamps, correlation IDs, and redaction.
\item Online evaluators combine sampled judges, deterministic checks, human feedback, and business outcomes.
\end{enumerate}
\noindent\textbf{Answer.} Counters, histograms, gauges, service-level indicators, and objectives quantify reliability and performance.
\par\textit{Counters, histograms, gauges, service-level indicators, and objectives quantify reliability and performance. Tradeoff: Aggregate signals scale well but cannot explain individual semantic failures. Watch for: Averaging latency hides long-tail p95 and p99 user pain.}
\paragraph{MCQ-0782 [hard]} What is the central engineering tradeoff of Metrics and SLOs?
\begin{enumerate}[label=\Alph*.]
\item Aggregate signals scale well but cannot explain individual semantic failures.
\item A common ML platform improves lineage but may be heavier than a focused tracing tool.
\item Search and aggregation improve, while schemas must evolve consistently.
\item Live coverage improves realism but evaluations must not expose or delay users.
\end{enumerate}
\noindent\textbf{Answer.} Aggregate signals scale well but cannot explain individual semantic failures.
\par\textit{Counters, histograms, gauges, service-level indicators, and objectives quantify reliability and performance. Tradeoff: Aggregate signals scale well but cannot explain individual semantic failures. Watch for: Averaging latency hides long-tail p95 and p99 user pain.}
\paragraph{MCQ-0783 [hard]} Which failure mode should an interviewer expect you to identify for Metrics and SLOs?
\begin{enumerate}[label=\Alph*.]
\item Placing unbounded model output in one log field causes cost and ingestion failures.
\item Averaging latency hides long-tail p95 and p99 user pain.
\item Mixing incomparable scorer versions in one trend chart produces false regressions.
\item Treating implicit clicks as unambiguous quality labels creates bias.
\end{enumerate}
\noindent\textbf{Answer.} Averaging latency hides long-tail p95 and p99 user pain.
\par\textit{Counters, histograms, gauges, service-level indicators, and objectives quantify reliability and performance. Tradeoff: Aggregate signals scale well but cannot explain individual semantic failures. Watch for: Averaging latency hides long-tail p95 and p99 user pain.}
\paragraph{MCQ-0784 [medium]} A design review is specifically concerned with this risk: Averaging latency hides long-tail p95 and p99 user pain. Which topic should be reviewed first?
\begin{enumerate}[label=\Alph*.]
\item MLflow tracing
\item Metrics and SLOs
\item Structured logging
\item Evaluation in production
\end{enumerate}
\noindent\textbf{Answer.} Metrics and SLOs
\par\textit{Counters, histograms, gauges, service-level indicators, and objectives quantify reliability and performance. Tradeoff: Aggregate signals scale well but cannot explain individual semantic failures. Watch for: Averaging latency hides long-tail p95 and p99 user pain.}
\paragraph{MCQ-0785 [hard]} Which design note most accurately balances the benefits and costs of Metrics and SLOs?
\begin{enumerate}[label=\Alph*.]
\item A common ML platform improves lineage but may be heavier than a focused tracing tool.
\item Search and aggregation improve, while schemas must evolve consistently.
\item Aggregate signals scale well but cannot explain individual semantic failures.
\item Live coverage improves realism but evaluations must not expose or delay users.
\end{enumerate}
\noindent\textbf{Answer.} Aggregate signals scale well but cannot explain individual semantic failures.
\par\textit{Counters, histograms, gauges, service-level indicators, and objectives quantify reliability and performance. Tradeoff: Aggregate signals scale well but cannot explain individual semantic failures. Watch for: Averaging latency hides long-tail p95 and p99 user pain.}
\paragraph{MCQ-0786 [medium]} Which explanation would be strongest in a system-design interview about Metrics and SLOs?
\begin{enumerate}[label=\Alph*.]
\item MLflow connects traces, evaluation datasets, scorers, experiments, and production monitoring.
\item Online evaluators combine sampled judges, deterministic checks, human feedback, and business outcomes.
\item Structured events use stable fields, severity, timestamps, correlation IDs, and redaction.
\item Counters, histograms, gauges, service-level indicators, and objectives quantify reliability and performance.
\end{enumerate}
\noindent\textbf{Answer.} Counters, histograms, gauges, service-level indicators, and objectives quantify reliability and performance.
\par\textit{Counters, histograms, gauges, service-level indicators, and objectives quantify reliability and performance. Tradeoff: Aggregate signals scale well but cannot explain individual semantic failures. Watch for: Averaging latency hides long-tail p95 and p99 user pain.}
\paragraph{MCQ-0787 [hard]} Which production warning is most directly associated with Metrics and SLOs?
\begin{enumerate}[label=\Alph*.]
\item Mixing incomparable scorer versions in one trend chart produces false regressions.
\item Treating implicit clicks as unambiguous quality labels creates bias.
\item Averaging latency hides long-tail p95 and p99 user pain.
\item Placing unbounded model output in one log field causes cost and ingestion failures.
\end{enumerate}
\noindent\textbf{Answer.} Averaging latency hides long-tail p95 and p99 user pain.
\par\textit{Counters, histograms, gauges, service-level indicators, and objectives quantify reliability and performance. Tradeoff: Aggregate signals scale well but cannot explain individual semantic failures. Watch for: Averaging latency hides long-tail p95 and p99 user pain.}
\paragraph{FC-0788 [medium]} Define Metrics and SLOs in interview-ready language.
\noindent\textbf{Answer.} Counters, histograms, gauges, service-level indicators, and objectives quantify reliability and performance.
\par\textit{Counters, histograms, gauges, service-level indicators, and objectives quantify reliability and performance.}
\paragraph{FC-0789 [hard]} State the main tradeoff and production pitfall for Metrics and SLOs.
\noindent\textbf{Answer.} Tradeoff: Aggregate signals scale well but cannot explain individual semantic failures. Pitfall: Averaging latency hides long-tail p95 and p99 user pain.
\par\textit{Tradeoff: Aggregate signals scale well but cannot explain individual semantic failures. Pitfall: Averaging latency hides long-tail p95 and p99 user pain.}
\paragraph{LA-0790 [hard]} Explain Metrics and SLOs from first principles, then show how you would evaluate and operate it in production.
\noindent\textbf{Answer.} Start with the mechanism: Counters, histograms, gauges, service-level indicators, and objectives quantify reliability and performance. Make the tradeoff explicit: Aggregate signals scale well but cannot explain individual semantic failures. Define workload-specific offline metrics and a latency/cost budget, test important slices, instrument the critical path, and create a rollback criterion. The production review must explicitly guard against this failure: Averaging latency hides long-tail p95 and p99 user pain.
\par\textit{A strong answer connects mechanism, measurable tradeoffs, failure isolation, observability, and rollback rather than listing product features.}
\paragraph{MCQ-0791 [medium]} Which statement best describes Structured logging?
\begin{enumerate}[label=\Alph*.]
\item OpenTelemetry provides vendor-neutral APIs, SDKs, collectors, and semantic conventions for traces, metrics, logs, and profiles.
\item LangSmith provides tracing, datasets, evaluation, prompt management, and deployment-oriented tooling for LLM applications.
\item Structured events use stable fields, severity, timestamps, correlation IDs, and redaction.
\item Phoenix provides open-source tracing and evaluation for LLM, retrieval, and agent applications with OpenTelemetry support.
\end{enumerate}
\noindent\textbf{Answer.} Structured events use stable fields, severity, timestamps, correlation IDs, and redaction.
\par\textit{Structured events use stable fields, severity, timestamps, correlation IDs, and redaction. Tradeoff: Search and aggregation improve, while schemas must evolve consistently. Watch for: Placing unbounded model output in one log field causes cost and ingestion failures.}
\paragraph{MCQ-0792 [hard]} What is the central engineering tradeoff of Structured logging?
\begin{enumerate}[label=\Alph*.]
\item Search and aggregation improve, while schemas must evolve consistently.
\item Portability improves, while instrumentation quality and backend cost remain engineering work.
\item It unifies qualitative traces and quantitative evals, while production operations still need design.
\item Tight LangChain integration accelerates work but requires portability planning.
\end{enumerate}
\noindent\textbf{Answer.} Search and aggregation improve, while schemas must evolve consistently.
\par\textit{Structured events use stable fields, severity, timestamps, correlation IDs, and redaction. Tradeoff: Search and aggregation improve, while schemas must evolve consistently. Watch for: Placing unbounded model output in one log field causes cost and ingestion failures.}
\paragraph{MCQ-0793 [hard]} Which failure mode should an interviewer expect you to identify for Structured logging?
\begin{enumerate}[label=\Alph*.]
\item Depending only on vendor UI annotations leaves evaluation logic unreproducible.
\item Judge scores without slice analysis can mask failures for important cohorts.
\item Placing unbounded model output in one log field causes cost and ingestion failures.
\item High-cardinality prompt or user attributes can overwhelm telemetry systems.
\end{enumerate}
\noindent\textbf{Answer.} Placing unbounded model output in one log field causes cost and ingestion failures.
\par\textit{Structured events use stable fields, severity, timestamps, correlation IDs, and redaction. Tradeoff: Search and aggregation improve, while schemas must evolve consistently. Watch for: Placing unbounded model output in one log field causes cost and ingestion failures.}
\paragraph{MCQ-0794 [medium]} A design review is specifically concerned with this risk: Placing unbounded model output in one log field causes cost and ingestion failures. Which topic should be reviewed first?
\begin{enumerate}[label=\Alph*.]
\item Structured logging
\item LangSmith
\item OpenTelemetry
\item Arize Phoenix
\end{enumerate}
\noindent\textbf{Answer.} Structured logging
\par\textit{Structured events use stable fields, severity, timestamps, correlation IDs, and redaction. Tradeoff: Search and aggregation improve, while schemas must evolve consistently. Watch for: Placing unbounded model output in one log field causes cost and ingestion failures.}
\paragraph{MCQ-0795 [hard]} Which design note most accurately balances the benefits and costs of Structured logging?
\begin{enumerate}[label=\Alph*.]
\item It unifies qualitative traces and quantitative evals, while production operations still need design.
\item Tight LangChain integration accelerates work but requires portability planning.
\item Search and aggregation improve, while schemas must evolve consistently.
\item Portability improves, while instrumentation quality and backend cost remain engineering work.
\end{enumerate}
\noindent\textbf{Answer.} Search and aggregation improve, while schemas must evolve consistently.
\par\textit{Structured events use stable fields, severity, timestamps, correlation IDs, and redaction. Tradeoff: Search and aggregation improve, while schemas must evolve consistently. Watch for: Placing unbounded model output in one log field causes cost and ingestion failures.}
\paragraph{MCQ-0796 [medium]} Which explanation would be strongest in a system-design interview about Structured logging?
\begin{enumerate}[label=\Alph*.]
\item Structured events use stable fields, severity, timestamps, correlation IDs, and redaction.
\item Phoenix provides open-source tracing and evaluation for LLM, retrieval, and agent applications with OpenTelemetry support.
\item LangSmith provides tracing, datasets, evaluation, prompt management, and deployment-oriented tooling for LLM applications.
\item OpenTelemetry provides vendor-neutral APIs, SDKs, collectors, and semantic conventions for traces, metrics, logs, and profiles.
\end{enumerate}
\noindent\textbf{Answer.} Structured events use stable fields, severity, timestamps, correlation IDs, and redaction.
\par\textit{Structured events use stable fields, severity, timestamps, correlation IDs, and redaction. Tradeoff: Search and aggregation improve, while schemas must evolve consistently. Watch for: Placing unbounded model output in one log field causes cost and ingestion failures.}
\paragraph{MCQ-0797 [hard]} Which production warning is most directly associated with Structured logging?
\begin{enumerate}[label=\Alph*.]
\item Depending only on vendor UI annotations leaves evaluation logic unreproducible.
\item High-cardinality prompt or user attributes can overwhelm telemetry systems.
\item Placing unbounded model output in one log field causes cost and ingestion failures.
\item Judge scores without slice analysis can mask failures for important cohorts.
\end{enumerate}
\noindent\textbf{Answer.} Placing unbounded model output in one log field causes cost and ingestion failures.
\par\textit{Structured events use stable fields, severity, timestamps, correlation IDs, and redaction. Tradeoff: Search and aggregation improve, while schemas must evolve consistently. Watch for: Placing unbounded model output in one log field causes cost and ingestion failures.}
\paragraph{FC-0798 [medium]} Define Structured logging in interview-ready language.
\noindent\textbf{Answer.} Structured events use stable fields, severity, timestamps, correlation IDs, and redaction.
\par\textit{Structured events use stable fields, severity, timestamps, correlation IDs, and redaction.}
\paragraph{FC-0799 [hard]} State the main tradeoff and production pitfall for Structured logging.
\noindent\textbf{Answer.} Tradeoff: Search and aggregation improve, while schemas must evolve consistently. Pitfall: Placing unbounded model output in one log field causes cost and ingestion failures.
\par\textit{Tradeoff: Search and aggregation improve, while schemas must evolve consistently. Pitfall: Placing unbounded model output in one log field causes cost and ingestion failures.}
\paragraph{LA-0800 [hard]} Explain Structured logging from first principles, then show how you would evaluate and operate it in production.
\noindent\textbf{Answer.} Start with the mechanism: Structured events use stable fields, severity, timestamps, correlation IDs, and redaction. Make the tradeoff explicit: Search and aggregation improve, while schemas must evolve consistently. Define workload-specific offline metrics and a latency/cost budget, test important slices, instrument the critical path, and create a rollback criterion. The production review must explicitly guard against this failure: Placing unbounded model output in one log field causes cost and ingestion failures.
\par\textit{A strong answer connects mechanism, measurable tradeoffs, failure isolation, observability, and rollback rather than listing product features.}
\paragraph{MCQ-0801 [medium]} Which statement best describes LangSmith?
\begin{enumerate}[label=\Alph*.]
\item LangSmith provides tracing, datasets, evaluation, prompt management, and deployment-oriented tooling for LLM applications.
\item OpenTelemetry provides vendor-neutral APIs, SDKs, collectors, and semantic conventions for traces, metrics, logs, and profiles.
\item Online evaluators combine sampled judges, deterministic checks, human feedback, and business outcomes.
\item Structured events use stable fields, severity, timestamps, correlation IDs, and redaction.
\end{enumerate}
\noindent\textbf{Answer.} LangSmith provides tracing, datasets, evaluation, prompt management, and deployment-oriented tooling for LLM applications.
\par\textit{LangSmith provides tracing, datasets, evaluation, prompt management, and deployment-oriented tooling for LLM applications. Tradeoff: Tight LangChain integration accelerates work but requires portability planning. Watch for: Depending only on vendor UI annotations leaves evaluation logic unreproducible.}
\paragraph{MCQ-0802 [hard]} What is the central engineering tradeoff of LangSmith?
\begin{enumerate}[label=\Alph*.]
\item Portability improves, while instrumentation quality and backend cost remain engineering work.
\item Tight LangChain integration accelerates work but requires portability planning.
\item Search and aggregation improve, while schemas must evolve consistently.
\item Live coverage improves realism but evaluations must not expose or delay users.
\end{enumerate}
\noindent\textbf{Answer.} Tight LangChain integration accelerates work but requires portability planning.
\par\textit{LangSmith provides tracing, datasets, evaluation, prompt management, and deployment-oriented tooling for LLM applications. Tradeoff: Tight LangChain integration accelerates work but requires portability planning. Watch for: Depending only on vendor UI annotations leaves evaluation logic unreproducible.}
\paragraph{MCQ-0803 [hard]} Which failure mode should an interviewer expect you to identify for LangSmith?
\begin{enumerate}[label=\Alph*.]
\item Treating implicit clicks as unambiguous quality labels creates bias.
\item Placing unbounded model output in one log field causes cost and ingestion failures.
\item High-cardinality prompt or user attributes can overwhelm telemetry systems.
\item Depending only on vendor UI annotations leaves evaluation logic unreproducible.
\end{enumerate}
\noindent\textbf{Answer.} Depending only on vendor UI annotations leaves evaluation logic unreproducible.
\par\textit{LangSmith provides tracing, datasets, evaluation, prompt management, and deployment-oriented tooling for LLM applications. Tradeoff: Tight LangChain integration accelerates work but requires portability planning. Watch for: Depending only on vendor UI annotations leaves evaluation logic unreproducible.}
\paragraph{MCQ-0804 [medium]} A design review is specifically concerned with this risk: Depending only on vendor UI annotations leaves evaluation logic unreproducible. Which topic should be reviewed first?
\begin{enumerate}[label=\Alph*.]
\item LangSmith
\item OpenTelemetry
\item Structured logging
\item Evaluation in production
\end{enumerate}
\noindent\textbf{Answer.} LangSmith
\par\textit{LangSmith provides tracing, datasets, evaluation, prompt management, and deployment-oriented tooling for LLM applications. Tradeoff: Tight LangChain integration accelerates work but requires portability planning. Watch for: Depending only on vendor UI annotations leaves evaluation logic unreproducible.}
\paragraph{MCQ-0805 [hard]} Which design note most accurately balances the benefits and costs of LangSmith?
\begin{enumerate}[label=\Alph*.]
\item Live coverage improves realism but evaluations must not expose or delay users.
\item Tight LangChain integration accelerates work but requires portability planning.
\item Search and aggregation improve, while schemas must evolve consistently.
\item Portability improves, while instrumentation quality and backend cost remain engineering work.
\end{enumerate}
\noindent\textbf{Answer.} Tight LangChain integration accelerates work but requires portability planning.
\par\textit{LangSmith provides tracing, datasets, evaluation, prompt management, and deployment-oriented tooling for LLM applications. Tradeoff: Tight LangChain integration accelerates work but requires portability planning. Watch for: Depending only on vendor UI annotations leaves evaluation logic unreproducible.}
\paragraph{MCQ-0806 [medium]} Which explanation would be strongest in a system-design interview about LangSmith?
\begin{enumerate}[label=\Alph*.]
\item Structured events use stable fields, severity, timestamps, correlation IDs, and redaction.
\item Online evaluators combine sampled judges, deterministic checks, human feedback, and business outcomes.
\item LangSmith provides tracing, datasets, evaluation, prompt management, and deployment-oriented tooling for LLM applications.
\item OpenTelemetry provides vendor-neutral APIs, SDKs, collectors, and semantic conventions for traces, metrics, logs, and profiles.
\end{enumerate}
\noindent\textbf{Answer.} LangSmith provides tracing, datasets, evaluation, prompt management, and deployment-oriented tooling for LLM applications.
\par\textit{LangSmith provides tracing, datasets, evaluation, prompt management, and deployment-oriented tooling for LLM applications. Tradeoff: Tight LangChain integration accelerates work but requires portability planning. Watch for: Depending only on vendor UI annotations leaves evaluation logic unreproducible.}
\paragraph{MCQ-0807 [hard]} Which production warning is most directly associated with LangSmith?
\begin{enumerate}[label=\Alph*.]
\item Depending only on vendor UI annotations leaves evaluation logic unreproducible.
\item Placing unbounded model output in one log field causes cost and ingestion failures.
\item High-cardinality prompt or user attributes can overwhelm telemetry systems.
\item Treating implicit clicks as unambiguous quality labels creates bias.
\end{enumerate}
\noindent\textbf{Answer.} Depending only on vendor UI annotations leaves evaluation logic unreproducible.
\par\textit{LangSmith provides tracing, datasets, evaluation, prompt management, and deployment-oriented tooling for LLM applications. Tradeoff: Tight LangChain integration accelerates work but requires portability planning. Watch for: Depending only on vendor UI annotations leaves evaluation logic unreproducible.}
\paragraph{FC-0808 [medium]} Define LangSmith in interview-ready language.
\noindent\textbf{Answer.} LangSmith provides tracing, datasets, evaluation, prompt management, and deployment-oriented tooling for LLM applications.
\par\textit{LangSmith provides tracing, datasets, evaluation, prompt management, and deployment-oriented tooling for LLM applications.}
\paragraph{FC-0809 [hard]} State the main tradeoff and production pitfall for LangSmith.
\noindent\textbf{Answer.} Tradeoff: Tight LangChain integration accelerates work but requires portability planning. Pitfall: Depending only on vendor UI annotations leaves evaluation logic unreproducible.
\par\textit{Tradeoff: Tight LangChain integration accelerates work but requires portability planning. Pitfall: Depending only on vendor UI annotations leaves evaluation logic unreproducible.}
\paragraph{LA-0810 [hard]} Explain LangSmith from first principles, then show how you would evaluate and operate it in production.
\noindent\textbf{Answer.} Start with the mechanism: LangSmith provides tracing, datasets, evaluation, prompt management, and deployment-oriented tooling for LLM applications. Make the tradeoff explicit: Tight LangChain integration accelerates work but requires portability planning. Define workload-specific offline metrics and a latency/cost budget, test important slices, instrument the critical path, and create a rollback criterion. The production review must explicitly guard against this failure: Depending only on vendor UI annotations leaves evaluation logic unreproducible.
\par\textit{A strong answer connects mechanism, measurable tradeoffs, failure isolation, observability, and rollback rather than listing product features.}
\paragraph{MCQ-0811 [medium]} Which statement best describes Langfuse?
\begin{enumerate}[label=\Alph*.]
\item Langfuse is an open-source platform for LLM traces, sessions, prompts, scores, datasets, and cost analytics.
\item Traces connect request spans across gateways, retrievers, models, tools, queues, and databases.
\item Phoenix provides open-source tracing and evaluation for LLM, retrieval, and agent applications with OpenTelemetry support.
\item Weave tracks calls, objects, datasets, evaluations, and model application experiments.
\end{enumerate}
\noindent\textbf{Answer.} Langfuse is an open-source platform for LLM traces, sessions, prompts, scores, datasets, and cost analytics.
\par\textit{Langfuse is an open-source platform for LLM traces, sessions, prompts, scores, datasets, and cost analytics. Tradeoff: Self-hosting offers control but shifts scaling, upgrades, and security to the team. Watch for: Capturing every token forever creates expensive sensitive retention.}
\paragraph{MCQ-0812 [hard]} What is the central engineering tradeoff of Langfuse?
\begin{enumerate}[label=\Alph*.]
\item They reveal critical paths but sampling can hide rare quality failures.
\item It unifies qualitative traces and quantitative evals, while production operations still need design.
\item Rich experiment lineage helps collaboration but introduces another governance surface.
\item Self-hosting offers control but shifts scaling, upgrades, and security to the team.
\end{enumerate}
\noindent\textbf{Answer.} Self-hosting offers control but shifts scaling, upgrades, and security to the team.
\par\textit{Langfuse is an open-source platform for LLM traces, sessions, prompts, scores, datasets, and cost analytics. Tradeoff: Self-hosting offers control but shifts scaling, upgrades, and security to the team. Watch for: Capturing every token forever creates expensive sensitive retention.}
\paragraph{MCQ-0813 [hard]} Which failure mode should an interviewer expect you to identify for Langfuse?
\begin{enumerate}[label=\Alph*.]
\item Capturing every token forever creates expensive sensitive retention.
\item Judge scores without slice analysis can mask failures for important cohorts.
\item Reusing trace identifiers across independent user requests breaks isolation.
\item Saving mutable datasets without immutable versions undermines reproducibility.
\end{enumerate}
\noindent\textbf{Answer.} Capturing every token forever creates expensive sensitive retention.
\par\textit{Langfuse is an open-source platform for LLM traces, sessions, prompts, scores, datasets, and cost analytics. Tradeoff: Self-hosting offers control but shifts scaling, upgrades, and security to the team. Watch for: Capturing every token forever creates expensive sensitive retention.}
\paragraph{MCQ-0814 [medium]} A design review is specifically concerned with this risk: Capturing every token forever creates expensive sensitive retention. Which topic should be reviewed first?
\begin{enumerate}[label=\Alph*.]
\item Distributed tracing
\item W\&B Weave
\item Langfuse
\item Arize Phoenix
\end{enumerate}
\noindent\textbf{Answer.} Langfuse
\par\textit{Langfuse is an open-source platform for LLM traces, sessions, prompts, scores, datasets, and cost analytics. Tradeoff: Self-hosting offers control but shifts scaling, upgrades, and security to the team. Watch for: Capturing every token forever creates expensive sensitive retention.}
\paragraph{MCQ-0815 [hard]} Which design note most accurately balances the benefits and costs of Langfuse?
\begin{enumerate}[label=\Alph*.]
\item Rich experiment lineage helps collaboration but introduces another governance surface.
\item Self-hosting offers control but shifts scaling, upgrades, and security to the team.
\item It unifies qualitative traces and quantitative evals, while production operations still need design.
\item They reveal critical paths but sampling can hide rare quality failures.
\end{enumerate}
\noindent\textbf{Answer.} Self-hosting offers control but shifts scaling, upgrades, and security to the team.
\par\textit{Langfuse is an open-source platform for LLM traces, sessions, prompts, scores, datasets, and cost analytics. Tradeoff: Self-hosting offers control but shifts scaling, upgrades, and security to the team. Watch for: Capturing every token forever creates expensive sensitive retention.}
\paragraph{MCQ-0816 [medium]} Which explanation would be strongest in a system-design interview about Langfuse?
\begin{enumerate}[label=\Alph*.]
\item Weave tracks calls, objects, datasets, evaluations, and model application experiments.
\item Phoenix provides open-source tracing and evaluation for LLM, retrieval, and agent applications with OpenTelemetry support.
\item Traces connect request spans across gateways, retrievers, models, tools, queues, and databases.
\item Langfuse is an open-source platform for LLM traces, sessions, prompts, scores, datasets, and cost analytics.
\end{enumerate}
\noindent\textbf{Answer.} Langfuse is an open-source platform for LLM traces, sessions, prompts, scores, datasets, and cost analytics.
\par\textit{Langfuse is an open-source platform for LLM traces, sessions, prompts, scores, datasets, and cost analytics. Tradeoff: Self-hosting offers control but shifts scaling, upgrades, and security to the team. Watch for: Capturing every token forever creates expensive sensitive retention.}
\paragraph{MCQ-0817 [hard]} Which production warning is most directly associated with Langfuse?
\begin{enumerate}[label=\Alph*.]
\item Capturing every token forever creates expensive sensitive retention.
\item Reusing trace identifiers across independent user requests breaks isolation.
\item Saving mutable datasets without immutable versions undermines reproducibility.
\item Judge scores without slice analysis can mask failures for important cohorts.
\end{enumerate}
\noindent\textbf{Answer.} Capturing every token forever creates expensive sensitive retention.
\par\textit{Langfuse is an open-source platform for LLM traces, sessions, prompts, scores, datasets, and cost analytics. Tradeoff: Self-hosting offers control but shifts scaling, upgrades, and security to the team. Watch for: Capturing every token forever creates expensive sensitive retention.}
\paragraph{FC-0818 [medium]} Define Langfuse in interview-ready language.
\noindent\textbf{Answer.} Langfuse is an open-source platform for LLM traces, sessions, prompts, scores, datasets, and cost analytics.
\par\textit{Langfuse is an open-source platform for LLM traces, sessions, prompts, scores, datasets, and cost analytics.}
\paragraph{FC-0819 [hard]} State the main tradeoff and production pitfall for Langfuse.
\noindent\textbf{Answer.} Tradeoff: Self-hosting offers control but shifts scaling, upgrades, and security to the team. Pitfall: Capturing every token forever creates expensive sensitive retention.
\par\textit{Tradeoff: Self-hosting offers control but shifts scaling, upgrades, and security to the team. Pitfall: Capturing every token forever creates expensive sensitive retention.}
\paragraph{LA-0820 [hard]} Explain Langfuse from first principles, then show how you would evaluate and operate it in production.
\noindent\textbf{Answer.} Start with the mechanism: Langfuse is an open-source platform for LLM traces, sessions, prompts, scores, datasets, and cost analytics. Make the tradeoff explicit: Self-hosting offers control but shifts scaling, upgrades, and security to the team. Define workload-specific offline metrics and a latency/cost budget, test important slices, instrument the critical path, and create a rollback criterion. The production review must explicitly guard against this failure: Capturing every token forever creates expensive sensitive retention.
\par\textit{A strong answer connects mechanism, measurable tradeoffs, failure isolation, observability, and rollback rather than listing product features.}
\paragraph{MCQ-0821 [medium]} Which statement best describes Arize Phoenix?
\begin{enumerate}[label=\Alph*.]
\item Phoenix provides open-source tracing and evaluation for LLM, retrieval, and agent applications with OpenTelemetry support.
\item Drift monitoring detects changes in inputs, retrieved evidence, outputs, model versions, and user behavior.
\item Weave tracks calls, objects, datasets, evaluations, and model application experiments.
\item LLM traces capture prompts, outputs, model settings, token usage, latency, tool calls, retrieval, and evaluations.
\end{enumerate}
\noindent\textbf{Answer.} Phoenix provides open-source tracing and evaluation for LLM, retrieval, and agent applications with OpenTelemetry support.
\par\textit{Phoenix provides open-source tracing and evaluation for LLM, retrieval, and agent applications with OpenTelemetry support. Tradeoff: It unifies qualitative traces and quantitative evals, while production operations still need design. Watch for: Judge scores without slice analysis can mask failures for important cohorts.}
\paragraph{MCQ-0822 [hard]} What is the central engineering tradeoff of Arize Phoenix?
\begin{enumerate}[label=\Alph*.]
\item It unifies qualitative traces and quantitative evals, while production operations still need design.
\item Alerts catch regressions but thresholds must distinguish seasonality from harm.
\item Rich experiment lineage helps collaboration but introduces another governance surface.
\item Debugging improves, while content capture raises privacy and retention risk.
\end{enumerate}
\noindent\textbf{Answer.} It unifies qualitative traces and quantitative evals, while production operations still need design.
\par\textit{Phoenix provides open-source tracing and evaluation for LLM, retrieval, and agent applications with OpenTelemetry support. Tradeoff: It unifies qualitative traces and quantitative evals, while production operations still need design. Watch for: Judge scores without slice analysis can mask failures for important cohorts.}
\paragraph{MCQ-0823 [hard]} Which failure mode should an interviewer expect you to identify for Arize Phoenix?
\begin{enumerate}[label=\Alph*.]
\item Judge scores without slice analysis can mask failures for important cohorts.
\item Logging secrets or personal data by default creates a second sensitive datastore.
\item Saving mutable datasets without immutable versions undermines reproducibility.
\item Measuring embedding distribution shift alone does not prove answer quality changed.
\end{enumerate}
\noindent\textbf{Answer.} Judge scores without slice analysis can mask failures for important cohorts.
\par\textit{Phoenix provides open-source tracing and evaluation for LLM, retrieval, and agent applications with OpenTelemetry support. Tradeoff: It unifies qualitative traces and quantitative evals, while production operations still need design. Watch for: Judge scores without slice analysis can mask failures for important cohorts.}
\paragraph{MCQ-0824 [medium]} A design review is specifically concerned with this risk: Judge scores without slice analysis can mask failures for important cohorts. Which topic should be reviewed first?
\begin{enumerate}[label=\Alph*.]
\item W\&B Weave
\item Arize Phoenix
\item LLM tracing
\item Prompt and model drift
\end{enumerate}
\noindent\textbf{Answer.} Arize Phoenix
\par\textit{Phoenix provides open-source tracing and evaluation for LLM, retrieval, and agent applications with OpenTelemetry support. Tradeoff: It unifies qualitative traces and quantitative evals, while production operations still need design. Watch for: Judge scores without slice analysis can mask failures for important cohorts.}
\paragraph{MCQ-0825 [hard]} Which design note most accurately balances the benefits and costs of Arize Phoenix?
\begin{enumerate}[label=\Alph*.]
\item Debugging improves, while content capture raises privacy and retention risk.
\item It unifies qualitative traces and quantitative evals, while production operations still need design.
\item Rich experiment lineage helps collaboration but introduces another governance surface.
\item Alerts catch regressions but thresholds must distinguish seasonality from harm.
\end{enumerate}
\noindent\textbf{Answer.} It unifies qualitative traces and quantitative evals, while production operations still need design.
\par\textit{Phoenix provides open-source tracing and evaluation for LLM, retrieval, and agent applications with OpenTelemetry support. Tradeoff: It unifies qualitative traces and quantitative evals, while production operations still need design. Watch for: Judge scores without slice analysis can mask failures for important cohorts.}
\paragraph{MCQ-0826 [medium]} Which explanation would be strongest in a system-design interview about Arize Phoenix?
\begin{enumerate}[label=\Alph*.]
\item Phoenix provides open-source tracing and evaluation for LLM, retrieval, and agent applications with OpenTelemetry support.
\item Weave tracks calls, objects, datasets, evaluations, and model application experiments.
\item Drift monitoring detects changes in inputs, retrieved evidence, outputs, model versions, and user behavior.
\item LLM traces capture prompts, outputs, model settings, token usage, latency, tool calls, retrieval, and evaluations.
\end{enumerate}
\noindent\textbf{Answer.} Phoenix provides open-source tracing and evaluation for LLM, retrieval, and agent applications with OpenTelemetry support.
\par\textit{Phoenix provides open-source tracing and evaluation for LLM, retrieval, and agent applications with OpenTelemetry support. Tradeoff: It unifies qualitative traces and quantitative evals, while production operations still need design. Watch for: Judge scores without slice analysis can mask failures for important cohorts.}
\paragraph{MCQ-0827 [hard]} Which production warning is most directly associated with Arize Phoenix?
\begin{enumerate}[label=\Alph*.]
\item Logging secrets or personal data by default creates a second sensitive datastore.
\item Saving mutable datasets without immutable versions undermines reproducibility.
\item Measuring embedding distribution shift alone does not prove answer quality changed.
\item Judge scores without slice analysis can mask failures for important cohorts.
\end{enumerate}
\noindent\textbf{Answer.} Judge scores without slice analysis can mask failures for important cohorts.
\par\textit{Phoenix provides open-source tracing and evaluation for LLM, retrieval, and agent applications with OpenTelemetry support. Tradeoff: It unifies qualitative traces and quantitative evals, while production operations still need design. Watch for: Judge scores without slice analysis can mask failures for important cohorts.}
\paragraph{FC-0828 [medium]} Define Arize Phoenix in interview-ready language.
\noindent\textbf{Answer.} Phoenix provides open-source tracing and evaluation for LLM, retrieval, and agent applications with OpenTelemetry support.
\par\textit{Phoenix provides open-source tracing and evaluation for LLM, retrieval, and agent applications with OpenTelemetry support.}
\paragraph{FC-0829 [hard]} State the main tradeoff and production pitfall for Arize Phoenix.
\noindent\textbf{Answer.} Tradeoff: It unifies qualitative traces and quantitative evals, while production operations still need design. Pitfall: Judge scores without slice analysis can mask failures for important cohorts.
\par\textit{Tradeoff: It unifies qualitative traces and quantitative evals, while production operations still need design. Pitfall: Judge scores without slice analysis can mask failures for important cohorts.}
\paragraph{LA-0830 [hard]} Explain Arize Phoenix from first principles, then show how you would evaluate and operate it in production.
\noindent\textbf{Answer.} Start with the mechanism: Phoenix provides open-source tracing and evaluation for LLM, retrieval, and agent applications with OpenTelemetry support. Make the tradeoff explicit: It unifies qualitative traces and quantitative evals, while production operations still need design. Define workload-specific offline metrics and a latency/cost budget, test important slices, instrument the critical path, and create a rollback criterion. The production review must explicitly guard against this failure: Judge scores without slice analysis can mask failures for important cohorts.
\par\textit{A strong answer connects mechanism, measurable tradeoffs, failure isolation, observability, and rollback rather than listing product features.}
\paragraph{MCQ-0831 [medium]} Which statement best describes MLflow tracing?
\begin{enumerate}[label=\Alph*.]
\item Drift monitoring detects changes in inputs, retrieved evidence, outputs, model versions, and user behavior.
\item MLflow connects traces, evaluation datasets, scorers, experiments, and production monitoring.
\item Cost telemetry attributes input, output, cache, retrieval, tool, and infrastructure spend to requests and tenants.
\item Structured events use stable fields, severity, timestamps, correlation IDs, and redaction.
\end{enumerate}
\noindent\textbf{Answer.} MLflow connects traces, evaluation datasets, scorers, experiments, and production monitoring.
\par\textit{MLflow connects traces, evaluation datasets, scorers, experiments, and production monitoring. Tradeoff: A common ML platform improves lineage but may be heavier than a focused tracing tool. Watch for: Mixing incomparable scorer versions in one trend chart produces false regressions.}
\paragraph{MCQ-0832 [hard]} What is the central engineering tradeoff of MLflow tracing?
\begin{enumerate}[label=\Alph*.]
\item Unit economics become actionable, while provider pricing and caching rules evolve.
\item Search and aggregation improve, while schemas must evolve consistently.
\item A common ML platform improves lineage but may be heavier than a focused tracing tool.
\item Alerts catch regressions but thresholds must distinguish seasonality from harm.
\end{enumerate}
\noindent\textbf{Answer.} A common ML platform improves lineage but may be heavier than a focused tracing tool.
\par\textit{MLflow connects traces, evaluation datasets, scorers, experiments, and production monitoring. Tradeoff: A common ML platform improves lineage but may be heavier than a focused tracing tool. Watch for: Mixing incomparable scorer versions in one trend chart produces false regressions.}
\paragraph{MCQ-0833 [hard]} Which failure mode should an interviewer expect you to identify for MLflow tracing?
\begin{enumerate}[label=\Alph*.]
\item Measuring embedding distribution shift alone does not prove answer quality changed.
\item Mixing incomparable scorer versions in one trend chart produces false regressions.
\item Tracking only average cost hides pathological long-running agents.
\item Placing unbounded model output in one log field causes cost and ingestion failures.
\end{enumerate}
\noindent\textbf{Answer.} Mixing incomparable scorer versions in one trend chart produces false regressions.
\par\textit{MLflow connects traces, evaluation datasets, scorers, experiments, and production monitoring. Tradeoff: A common ML platform improves lineage but may be heavier than a focused tracing tool. Watch for: Mixing incomparable scorer versions in one trend chart produces false regressions.}
\paragraph{MCQ-0834 [medium]} A design review is specifically concerned with this risk: Mixing incomparable scorer versions in one trend chart produces false regressions. Which topic should be reviewed first?
\begin{enumerate}[label=\Alph*.]
\item Structured logging
\item Prompt and model drift
\item Cost observability
\item MLflow tracing
\end{enumerate}
\noindent\textbf{Answer.} MLflow tracing
\par\textit{MLflow connects traces, evaluation datasets, scorers, experiments, and production monitoring. Tradeoff: A common ML platform improves lineage but may be heavier than a focused tracing tool. Watch for: Mixing incomparable scorer versions in one trend chart produces false regressions.}
\paragraph{MCQ-0835 [hard]} Which design note most accurately balances the benefits and costs of MLflow tracing?
\begin{enumerate}[label=\Alph*.]
\item Unit economics become actionable, while provider pricing and caching rules evolve.
\item Search and aggregation improve, while schemas must evolve consistently.
\item A common ML platform improves lineage but may be heavier than a focused tracing tool.
\item Alerts catch regressions but thresholds must distinguish seasonality from harm.
\end{enumerate}
\noindent\textbf{Answer.} A common ML platform improves lineage but may be heavier than a focused tracing tool.
\par\textit{MLflow connects traces, evaluation datasets, scorers, experiments, and production monitoring. Tradeoff: A common ML platform improves lineage but may be heavier than a focused tracing tool. Watch for: Mixing incomparable scorer versions in one trend chart produces false regressions.}
\paragraph{MCQ-0836 [medium]} Which explanation would be strongest in a system-design interview about MLflow tracing?
\begin{enumerate}[label=\Alph*.]
\item Structured events use stable fields, severity, timestamps, correlation IDs, and redaction.
\item MLflow connects traces, evaluation datasets, scorers, experiments, and production monitoring.
\item Drift monitoring detects changes in inputs, retrieved evidence, outputs, model versions, and user behavior.
\item Cost telemetry attributes input, output, cache, retrieval, tool, and infrastructure spend to requests and tenants.
\end{enumerate}
\noindent\textbf{Answer.} MLflow connects traces, evaluation datasets, scorers, experiments, and production monitoring.
\par\textit{MLflow connects traces, evaluation datasets, scorers, experiments, and production monitoring. Tradeoff: A common ML platform improves lineage but may be heavier than a focused tracing tool. Watch for: Mixing incomparable scorer versions in one trend chart produces false regressions.}
\paragraph{MCQ-0837 [hard]} Which production warning is most directly associated with MLflow tracing?
\begin{enumerate}[label=\Alph*.]
\item Measuring embedding distribution shift alone does not prove answer quality changed.
\item Mixing incomparable scorer versions in one trend chart produces false regressions.
\item Tracking only average cost hides pathological long-running agents.
\item Placing unbounded model output in one log field causes cost and ingestion failures.
\end{enumerate}
\noindent\textbf{Answer.} Mixing incomparable scorer versions in one trend chart produces false regressions.
\par\textit{MLflow connects traces, evaluation datasets, scorers, experiments, and production monitoring. Tradeoff: A common ML platform improves lineage but may be heavier than a focused tracing tool. Watch for: Mixing incomparable scorer versions in one trend chart produces false regressions.}
\paragraph{FC-0838 [medium]} Define MLflow tracing in interview-ready language.
\noindent\textbf{Answer.} MLflow connects traces, evaluation datasets, scorers, experiments, and production monitoring.
\par\textit{MLflow connects traces, evaluation datasets, scorers, experiments, and production monitoring.}
\paragraph{FC-0839 [hard]} State the main tradeoff and production pitfall for MLflow tracing.
\noindent\textbf{Answer.} Tradeoff: A common ML platform improves lineage but may be heavier than a focused tracing tool. Pitfall: Mixing incomparable scorer versions in one trend chart produces false regressions.
\par\textit{Tradeoff: A common ML platform improves lineage but may be heavier than a focused tracing tool. Pitfall: Mixing incomparable scorer versions in one trend chart produces false regressions.}
\paragraph{LA-0840 [hard]} Explain MLflow tracing from first principles, then show how you would evaluate and operate it in production.
\noindent\textbf{Answer.} Start with the mechanism: MLflow connects traces, evaluation datasets, scorers, experiments, and production monitoring. Make the tradeoff explicit: A common ML platform improves lineage but may be heavier than a focused tracing tool. Define workload-specific offline metrics and a latency/cost budget, test important slices, instrument the critical path, and create a rollback criterion. The production review must explicitly guard against this failure: Mixing incomparable scorer versions in one trend chart produces false regressions.
\par\textit{A strong answer connects mechanism, measurable tradeoffs, failure isolation, observability, and rollback rather than listing product features.}
\paragraph{MCQ-0841 [medium]} Which statement best describes W\&B Weave?
\begin{enumerate}[label=\Alph*.]
\item Langfuse is an open-source platform for LLM traces, sessions, prompts, scores, datasets, and cost analytics.
\item Weave tracks calls, objects, datasets, evaluations, and model application experiments.
\item MLflow connects traces, evaluation datasets, scorers, experiments, and production monitoring.
\item Drift monitoring detects changes in inputs, retrieved evidence, outputs, model versions, and user behavior.
\end{enumerate}
\noindent\textbf{Answer.} Weave tracks calls, objects, datasets, evaluations, and model application experiments.
\par\textit{Weave tracks calls, objects, datasets, evaluations, and model application experiments. Tradeoff: Rich experiment lineage helps collaboration but introduces another governance surface. Watch for: Saving mutable datasets without immutable versions undermines reproducibility.}
\paragraph{MCQ-0842 [hard]} What is the central engineering tradeoff of W\&B Weave?
\begin{enumerate}[label=\Alph*.]
\item Alerts catch regressions but thresholds must distinguish seasonality from harm.
\item Rich experiment lineage helps collaboration but introduces another governance surface.
\item A common ML platform improves lineage but may be heavier than a focused tracing tool.
\item Self-hosting offers control but shifts scaling, upgrades, and security to the team.
\end{enumerate}
\noindent\textbf{Answer.} Rich experiment lineage helps collaboration but introduces another governance surface.
\par\textit{Weave tracks calls, objects, datasets, evaluations, and model application experiments. Tradeoff: Rich experiment lineage helps collaboration but introduces another governance surface. Watch for: Saving mutable datasets without immutable versions undermines reproducibility.}
\paragraph{MCQ-0843 [hard]} Which failure mode should an interviewer expect you to identify for W\&B Weave?
\begin{enumerate}[label=\Alph*.]
\item Saving mutable datasets without immutable versions undermines reproducibility.
\item Capturing every token forever creates expensive sensitive retention.
\item Measuring embedding distribution shift alone does not prove answer quality changed.
\item Mixing incomparable scorer versions in one trend chart produces false regressions.
\end{enumerate}
\noindent\textbf{Answer.} Saving mutable datasets without immutable versions undermines reproducibility.
\par\textit{Weave tracks calls, objects, datasets, evaluations, and model application experiments. Tradeoff: Rich experiment lineage helps collaboration but introduces another governance surface. Watch for: Saving mutable datasets without immutable versions undermines reproducibility.}
\paragraph{MCQ-0844 [medium]} A design review is specifically concerned with this risk: Saving mutable datasets without immutable versions undermines reproducibility. Which topic should be reviewed first?
\begin{enumerate}[label=\Alph*.]
\item MLflow tracing
\item W\&B Weave
\item Langfuse
\item Prompt and model drift
\end{enumerate}
\noindent\textbf{Answer.} W\&B Weave
\par\textit{Weave tracks calls, objects, datasets, evaluations, and model application experiments. Tradeoff: Rich experiment lineage helps collaboration but introduces another governance surface. Watch for: Saving mutable datasets without immutable versions undermines reproducibility.}
\paragraph{MCQ-0845 [hard]} Which design note most accurately balances the benefits and costs of W\&B Weave?
\begin{enumerate}[label=\Alph*.]
\item Rich experiment lineage helps collaboration but introduces another governance surface.
\item Alerts catch regressions but thresholds must distinguish seasonality from harm.
\item A common ML platform improves lineage but may be heavier than a focused tracing tool.
\item Self-hosting offers control but shifts scaling, upgrades, and security to the team.
\end{enumerate}
\noindent\textbf{Answer.} Rich experiment lineage helps collaboration but introduces another governance surface.
\par\textit{Weave tracks calls, objects, datasets, evaluations, and model application experiments. Tradeoff: Rich experiment lineage helps collaboration but introduces another governance surface. Watch for: Saving mutable datasets without immutable versions undermines reproducibility.}
\paragraph{MCQ-0846 [medium]} Which explanation would be strongest in a system-design interview about W\&B Weave?
\begin{enumerate}[label=\Alph*.]
\item Weave tracks calls, objects, datasets, evaluations, and model application experiments.
\item Drift monitoring detects changes in inputs, retrieved evidence, outputs, model versions, and user behavior.
\item Langfuse is an open-source platform for LLM traces, sessions, prompts, scores, datasets, and cost analytics.
\item MLflow connects traces, evaluation datasets, scorers, experiments, and production monitoring.
\end{enumerate}
\noindent\textbf{Answer.} Weave tracks calls, objects, datasets, evaluations, and model application experiments.
\par\textit{Weave tracks calls, objects, datasets, evaluations, and model application experiments. Tradeoff: Rich experiment lineage helps collaboration but introduces another governance surface. Watch for: Saving mutable datasets without immutable versions undermines reproducibility.}
\paragraph{MCQ-0847 [hard]} Which production warning is most directly associated with W\&B Weave?
\begin{enumerate}[label=\Alph*.]
\item Measuring embedding distribution shift alone does not prove answer quality changed.
\item Mixing incomparable scorer versions in one trend chart produces false regressions.
\item Saving mutable datasets without immutable versions undermines reproducibility.
\item Capturing every token forever creates expensive sensitive retention.
\end{enumerate}
\noindent\textbf{Answer.} Saving mutable datasets without immutable versions undermines reproducibility.
\par\textit{Weave tracks calls, objects, datasets, evaluations, and model application experiments. Tradeoff: Rich experiment lineage helps collaboration but introduces another governance surface. Watch for: Saving mutable datasets without immutable versions undermines reproducibility.}
\paragraph{FC-0848 [medium]} Define W\&B Weave in interview-ready language.
\noindent\textbf{Answer.} Weave tracks calls, objects, datasets, evaluations, and model application experiments.
\par\textit{Weave tracks calls, objects, datasets, evaluations, and model application experiments.}
\paragraph{FC-0849 [hard]} State the main tradeoff and production pitfall for W\&B Weave.
\noindent\textbf{Answer.} Tradeoff: Rich experiment lineage helps collaboration but introduces another governance surface. Pitfall: Saving mutable datasets without immutable versions undermines reproducibility.
\par\textit{Tradeoff: Rich experiment lineage helps collaboration but introduces another governance surface. Pitfall: Saving mutable datasets without immutable versions undermines reproducibility.}
\paragraph{LA-0850 [hard]} Explain W\&B Weave from first principles, then show how you would evaluate and operate it in production.
\noindent\textbf{Answer.} Start with the mechanism: Weave tracks calls, objects, datasets, evaluations, and model application experiments. Make the tradeoff explicit: Rich experiment lineage helps collaboration but introduces another governance surface. Define workload-specific offline metrics and a latency/cost budget, test important slices, instrument the critical path, and create a rollback criterion. The production review must explicitly guard against this failure: Saving mutable datasets without immutable versions undermines reproducibility.
\par\textit{A strong answer connects mechanism, measurable tradeoffs, failure isolation, observability, and rollback rather than listing product features.}
\paragraph{MCQ-0851 [medium]} Which statement best describes Prompt and model drift?
\begin{enumerate}[label=\Alph*.]
\item Drift monitoring detects changes in inputs, retrieved evidence, outputs, model versions, and user behavior.
\item MLflow connects traces, evaluation datasets, scorers, experiments, and production monitoring.
\item Online evaluators combine sampled judges, deterministic checks, human feedback, and business outcomes.
\item Structured events use stable fields, severity, timestamps, correlation IDs, and redaction.
\end{enumerate}
\noindent\textbf{Answer.} Drift monitoring detects changes in inputs, retrieved evidence, outputs, model versions, and user behavior.
\par\textit{Drift monitoring detects changes in inputs, retrieved evidence, outputs, model versions, and user behavior. Tradeoff: Alerts catch regressions but thresholds must distinguish seasonality from harm. Watch for: Measuring embedding distribution shift alone does not prove answer quality changed.}
\paragraph{MCQ-0852 [hard]} What is the central engineering tradeoff of Prompt and model drift?
\begin{enumerate}[label=\Alph*.]
\item Alerts catch regressions but thresholds must distinguish seasonality from harm.
\item Live coverage improves realism but evaluations must not expose or delay users.
\item A common ML platform improves lineage but may be heavier than a focused tracing tool.
\item Search and aggregation improve, while schemas must evolve consistently.
\end{enumerate}
\noindent\textbf{Answer.} Alerts catch regressions but thresholds must distinguish seasonality from harm.
\par\textit{Drift monitoring detects changes in inputs, retrieved evidence, outputs, model versions, and user behavior. Tradeoff: Alerts catch regressions but thresholds must distinguish seasonality from harm. Watch for: Measuring embedding distribution shift alone does not prove answer quality changed.}
\paragraph{MCQ-0853 [hard]} Which failure mode should an interviewer expect you to identify for Prompt and model drift?
\begin{enumerate}[label=\Alph*.]
\item Measuring embedding distribution shift alone does not prove answer quality changed.
\item Treating implicit clicks as unambiguous quality labels creates bias.
\item Placing unbounded model output in one log field causes cost and ingestion failures.
\item Mixing incomparable scorer versions in one trend chart produces false regressions.
\end{enumerate}
\noindent\textbf{Answer.} Measuring embedding distribution shift alone does not prove answer quality changed.
\par\textit{Drift monitoring detects changes in inputs, retrieved evidence, outputs, model versions, and user behavior. Tradeoff: Alerts catch regressions but thresholds must distinguish seasonality from harm. Watch for: Measuring embedding distribution shift alone does not prove answer quality changed.}
\paragraph{MCQ-0854 [medium]} A design review is specifically concerned with this risk: Measuring embedding distribution shift alone does not prove answer quality changed. Which topic should be reviewed first?
\begin{enumerate}[label=\Alph*.]
\item Evaluation in production
\item Prompt and model drift
\item Structured logging
\item MLflow tracing
\end{enumerate}
\noindent\textbf{Answer.} Prompt and model drift
\par\textit{Drift monitoring detects changes in inputs, retrieved evidence, outputs, model versions, and user behavior. Tradeoff: Alerts catch regressions but thresholds must distinguish seasonality from harm. Watch for: Measuring embedding distribution shift alone does not prove answer quality changed.}
\paragraph{MCQ-0855 [hard]} Which design note most accurately balances the benefits and costs of Prompt and model drift?
\begin{enumerate}[label=\Alph*.]
\item Alerts catch regressions but thresholds must distinguish seasonality from harm.
\item A common ML platform improves lineage but may be heavier than a focused tracing tool.
\item Live coverage improves realism but evaluations must not expose or delay users.
\item Search and aggregation improve, while schemas must evolve consistently.
\end{enumerate}
\noindent\textbf{Answer.} Alerts catch regressions but thresholds must distinguish seasonality from harm.
\par\textit{Drift monitoring detects changes in inputs, retrieved evidence, outputs, model versions, and user behavior. Tradeoff: Alerts catch regressions but thresholds must distinguish seasonality from harm. Watch for: Measuring embedding distribution shift alone does not prove answer quality changed.}
\paragraph{MCQ-0856 [medium]} Which explanation would be strongest in a system-design interview about Prompt and model drift?
\begin{enumerate}[label=\Alph*.]
\item MLflow connects traces, evaluation datasets, scorers, experiments, and production monitoring.
\item Drift monitoring detects changes in inputs, retrieved evidence, outputs, model versions, and user behavior.
\item Structured events use stable fields, severity, timestamps, correlation IDs, and redaction.
\item Online evaluators combine sampled judges, deterministic checks, human feedback, and business outcomes.
\end{enumerate}
\noindent\textbf{Answer.} Drift monitoring detects changes in inputs, retrieved evidence, outputs, model versions, and user behavior.
\par\textit{Drift monitoring detects changes in inputs, retrieved evidence, outputs, model versions, and user behavior. Tradeoff: Alerts catch regressions but thresholds must distinguish seasonality from harm. Watch for: Measuring embedding distribution shift alone does not prove answer quality changed.}
\paragraph{MCQ-0857 [hard]} Which production warning is most directly associated with Prompt and model drift?
\begin{enumerate}[label=\Alph*.]
\item Treating implicit clicks as unambiguous quality labels creates bias.
\item Measuring embedding distribution shift alone does not prove answer quality changed.
\item Placing unbounded model output in one log field causes cost and ingestion failures.
\item Mixing incomparable scorer versions in one trend chart produces false regressions.
\end{enumerate}
\noindent\textbf{Answer.} Measuring embedding distribution shift alone does not prove answer quality changed.
\par\textit{Drift monitoring detects changes in inputs, retrieved evidence, outputs, model versions, and user behavior. Tradeoff: Alerts catch regressions but thresholds must distinguish seasonality from harm. Watch for: Measuring embedding distribution shift alone does not prove answer quality changed.}
\paragraph{FC-0858 [medium]} Define Prompt and model drift in interview-ready language.
\noindent\textbf{Answer.} Drift monitoring detects changes in inputs, retrieved evidence, outputs, model versions, and user behavior.
\par\textit{Drift monitoring detects changes in inputs, retrieved evidence, outputs, model versions, and user behavior.}
\paragraph{FC-0859 [hard]} State the main tradeoff and production pitfall for Prompt and model drift.
\noindent\textbf{Answer.} Tradeoff: Alerts catch regressions but thresholds must distinguish seasonality from harm. Pitfall: Measuring embedding distribution shift alone does not prove answer quality changed.
\par\textit{Tradeoff: Alerts catch regressions but thresholds must distinguish seasonality from harm. Pitfall: Measuring embedding distribution shift alone does not prove answer quality changed.}
\paragraph{LA-0860 [hard]} Explain Prompt and model drift from first principles, then show how you would evaluate and operate it in production.
\noindent\textbf{Answer.} Start with the mechanism: Drift monitoring detects changes in inputs, retrieved evidence, outputs, model versions, and user behavior. Make the tradeoff explicit: Alerts catch regressions but thresholds must distinguish seasonality from harm. Define workload-specific offline metrics and a latency/cost budget, test important slices, instrument the critical path, and create a rollback criterion. The production review must explicitly guard against this failure: Measuring embedding distribution shift alone does not prove answer quality changed.
\par\textit{A strong answer connects mechanism, measurable tradeoffs, failure isolation, observability, and rollback rather than listing product features.}
\paragraph{MCQ-0861 [medium]} Which statement best describes Cost observability?
\begin{enumerate}[label=\Alph*.]
\item LangSmith provides tracing, datasets, evaluation, prompt management, and deployment-oriented tooling for LLM applications.
\item Phoenix provides open-source tracing and evaluation for LLM, retrieval, and agent applications with OpenTelemetry support.
\item Cost telemetry attributes input, output, cache, retrieval, tool, and infrastructure spend to requests and tenants.
\item Structured events use stable fields, severity, timestamps, correlation IDs, and redaction.
\end{enumerate}
\noindent\textbf{Answer.} Cost telemetry attributes input, output, cache, retrieval, tool, and infrastructure spend to requests and tenants.
\par\textit{Cost telemetry attributes input, output, cache, retrieval, tool, and infrastructure spend to requests and tenants. Tradeoff: Unit economics become actionable, while provider pricing and caching rules evolve. Watch for: Tracking only average cost hides pathological long-running agents.}
\paragraph{MCQ-0862 [hard]} What is the central engineering tradeoff of Cost observability?
\begin{enumerate}[label=\Alph*.]
\item Unit economics become actionable, while provider pricing and caching rules evolve.
\item Tight LangChain integration accelerates work but requires portability planning.
\item Search and aggregation improve, while schemas must evolve consistently.
\item It unifies qualitative traces and quantitative evals, while production operations still need design.
\end{enumerate}
\noindent\textbf{Answer.} Unit economics become actionable, while provider pricing and caching rules evolve.
\par\textit{Cost telemetry attributes input, output, cache, retrieval, tool, and infrastructure spend to requests and tenants. Tradeoff: Unit economics become actionable, while provider pricing and caching rules evolve. Watch for: Tracking only average cost hides pathological long-running agents.}
\paragraph{MCQ-0863 [hard]} Which failure mode should an interviewer expect you to identify for Cost observability?
\begin{enumerate}[label=\Alph*.]
\item Placing unbounded model output in one log field causes cost and ingestion failures.
\item Judge scores without slice analysis can mask failures for important cohorts.
\item Tracking only average cost hides pathological long-running agents.
\item Depending only on vendor UI annotations leaves evaluation logic unreproducible.
\end{enumerate}
\noindent\textbf{Answer.} Tracking only average cost hides pathological long-running agents.
\par\textit{Cost telemetry attributes input, output, cache, retrieval, tool, and infrastructure spend to requests and tenants. Tradeoff: Unit economics become actionable, while provider pricing and caching rules evolve. Watch for: Tracking only average cost hides pathological long-running agents.}
\paragraph{MCQ-0864 [medium]} A design review is specifically concerned with this risk: Tracking only average cost hides pathological long-running agents. Which topic should be reviewed first?
\begin{enumerate}[label=\Alph*.]
\item Structured logging
\item Arize Phoenix
\item Cost observability
\item LangSmith
\end{enumerate}
\noindent\textbf{Answer.} Cost observability
\par\textit{Cost telemetry attributes input, output, cache, retrieval, tool, and infrastructure spend to requests and tenants. Tradeoff: Unit economics become actionable, while provider pricing and caching rules evolve. Watch for: Tracking only average cost hides pathological long-running agents.}
\paragraph{MCQ-0865 [hard]} Which design note most accurately balances the benefits and costs of Cost observability?
\begin{enumerate}[label=\Alph*.]
\item Tight LangChain integration accelerates work but requires portability planning.
\item It unifies qualitative traces and quantitative evals, while production operations still need design.
\item Unit economics become actionable, while provider pricing and caching rules evolve.
\item Search and aggregation improve, while schemas must evolve consistently.
\end{enumerate}
\noindent\textbf{Answer.} Unit economics become actionable, while provider pricing and caching rules evolve.
\par\textit{Cost telemetry attributes input, output, cache, retrieval, tool, and infrastructure spend to requests and tenants. Tradeoff: Unit economics become actionable, while provider pricing and caching rules evolve. Watch for: Tracking only average cost hides pathological long-running agents.}
\paragraph{MCQ-0866 [medium]} Which explanation would be strongest in a system-design interview about Cost observability?
\begin{enumerate}[label=\Alph*.]
\item Cost telemetry attributes input, output, cache, retrieval, tool, and infrastructure spend to requests and tenants.
\item Phoenix provides open-source tracing and evaluation for LLM, retrieval, and agent applications with OpenTelemetry support.
\item Structured events use stable fields, severity, timestamps, correlation IDs, and redaction.
\item LangSmith provides tracing, datasets, evaluation, prompt management, and deployment-oriented tooling for LLM applications.
\end{enumerate}
\noindent\textbf{Answer.} Cost telemetry attributes input, output, cache, retrieval, tool, and infrastructure spend to requests and tenants.
\par\textit{Cost telemetry attributes input, output, cache, retrieval, tool, and infrastructure spend to requests and tenants. Tradeoff: Unit economics become actionable, while provider pricing and caching rules evolve. Watch for: Tracking only average cost hides pathological long-running agents.}
\paragraph{MCQ-0867 [hard]} Which production warning is most directly associated with Cost observability?
\begin{enumerate}[label=\Alph*.]
\item Placing unbounded model output in one log field causes cost and ingestion failures.
\item Depending only on vendor UI annotations leaves evaluation logic unreproducible.
\item Judge scores without slice analysis can mask failures for important cohorts.
\item Tracking only average cost hides pathological long-running agents.
\end{enumerate}
\noindent\textbf{Answer.} Tracking only average cost hides pathological long-running agents.
